% Contraction de texte : rédiger un résumé de texte — Français Tle Tech — Cours signé OnBuch+
% Programme officiel MINESEC. Compiler avec : tectonic N10-contraction-texte-rediger-resume-texte.tex
\newcommand{\DOCMATIERE}{Français}
\newcommand{\DOCNIVEAU}{Tle Tech}
\newcommand{\DOCMODULE}{Français – classe de Terminale technique}
\newcommand{\DOCLECON}{N10}
\newcommand{\DOCTITRE}{Contraction de texte : rédiger un résumé de texte}
% OnBuch+ — préambule « cours signés » ENRICHI (compilé avec tectonic / XeLaTeX)
% Système visuel « Pop » d'OnBuch+ : fond crème (jamais blanc), bordures encre
% épaisses, ombres « dures » décalées, titres Archivo Black en capitales.
% Version MATHÉMATIQUES : graphes (pgfplots/tikz), boîtes pédagogiques
% (définition, propriété, méthode, attention, le savais-tu, exercice résolu,
% exercice, TP, bilan), page de garde et signature OnBuch+ ancrée en bas.
%
% Macros attendues dans main.tex AVANT \input{preamble} :
%   \DOCMATIERE  \DOCNIVEAU  \DOCMODULE  \DOCLECON (numéro)  \DOCTITRE
\documentclass[11pt]{article}

\usepackage{fontspec}
\usepackage[french]{babel}
\usepackage[a4paper,margin=2cm,top=3.4cm,bottom=2.8cm,headheight=1.4cm,footskip=1.3cm]{geometry}
\usepackage{amsmath,amssymb,amsfonts}
\usepackage{mathtools} % \xmapsto, \coloneqq... souvent utilisés par les modèles
\usepackage{cancel} % simplifications barrées (bilans de forces)
% \abs{z} (module, valeur absolue) et \norm{u} (norme), utilisés par les modèles
\providecommand{\abs}{}\let\abs\relax\DeclarePairedDelimiter{\abs}{\lvert}{\rvert}
\providecommand{\norm}{}\let\norm\relax\DeclarePairedDelimiter{\norm}{\lVert}{\rVert}
\usepackage[table]{xcolor} % [table] : fournit aussi \rowcolors (alternance de lignes)
\usepackage{pagecolor}
\usepackage{tikz}
\usetikzlibrary{arrows.meta,positioning,calc,shapes.geometric,shapes.misc,shapes.arrows,decorations.pathmorphing,decorations.pathreplacing,decorations.markings,patterns,shadows,mindmap,trees,fit,backgrounds,matrix,intersections,angles,quotes}
\usepackage{pgfplots}
\usepackage[version=4]{mhchem} % équations nucléaires \ce{^{235}_{92}U}
\usepackage[european,siunitx]{circuitikz} % schémas électriques
\pgfplotsset{compat=1.18}
\usepgfplotslibrary{fillbetween,groupplots} % aires entre courbes (calcul intégral)
\usepackage{enumitem}
\usepackage{fancyhdr}
% [most] charge toutes les bibliothèques tcolorbox (listings, minted, raster,
% documentation...) et gonfle inutilement la mémoire TeX sur un document long
% (~300 boîtes) : on ne charge que ce qui est réellement utilisé.
\usepackage[skins,breakable]{tcolorbox}
\usepackage{graphicx}
\usepackage{colortbl}
\usepackage{booktabs}
\usepackage{tabularx}
\usepackage{multirow}
\usepackage{diagbox} % case d'en-tête diagonale des tableaux à double entrée
% Colonnes extensibles centrée / gauche / droite (C, L, R), souvent utilisées par les modèles
\newcolumntype{C}{>{\centering\arraybackslash}X}
\newcolumntype{L}{>{\raggedright\arraybackslash}X}
\newcolumntype{R}{>{\raggedleft\arraybackslash}X}
\usepackage{array}
\usepackage{multicol}
\usepackage{siunitx}
% parse-numbers=false : accepte aussi \SI{2,5 \times 10^{-3}}{...}
\sisetup{locale=FR,output-decimal-marker={,},group-minimum-digits=4,parse-numbers=false}
\DeclareSIUnit\ohm{\ensuremath{\symup{\Omega}}} % U+2126 absent de la police
\DeclareSIUnit\division{div} % réglages d'oscilloscope (ms/div, V/div)
\DeclareSIUnit\tr{tr} % tours (vitesse de rotation, ex. \SI{1500}{\tr\per\minute})
\DeclareSIUnit\molar{M} % concentration molaire (ex. \SI{3}{\molar}, \SI{1}{\micro\molar})
\DeclareSIUnit\an{an} % année (le modèle confond parfois avec \year, un compteur TeX) — ex. \SI{5}{\kilo\meter\per\an}
\DeclareSIUnit\FCFA{FCFA} % franc CFA (ex. \SI{500}{\FCFA\per\kilo\gram})
\DeclareSIUnit\degreeBrix{\textdegree Brix} % teneur en sucre d'un sirop/jus (réfractométrie)
\DeclareSIUnit\month{mois} % le modèle utilise parfois \month, un compteur TeX, comme unité de durée
\DeclareSIUnit\microliter{\micro\liter} % le modèle écrit parfois \microliter en un seul token
\DeclareSIUnit\million{millions} % dénombrement (ex. \SI{15}{\million\per\mL} spermatozoïdes)
\usepackage{qrcode}
% \degree : commande fréquemment utilisée par les modèles pour un angle ou une
% température en dehors de \SI{}{} (non fournie par les paquets chargés).
\providecommand{\degree}{^{\circ}}
% \qed (fin de démonstration), utilisé par les modèles sans amsthm
\providecommand{\qed}{\hfill$\square$}
% \barre : double barre des valeurs interdites dans un tableau de variations (tkz-tab)
\providecommand{\barre}{\ensuremath{\|}}
% environnement proof (amsthm non chargé) utilisé par les modèles
\ifdefined\proof\else\newenvironment{proof}[1][Démonstration]{\par\noindent\textit{#1.}\ }{\qed\par}\fi
\usepackage{lastpage}
\usepackage{hyperref}
\hypersetup{hidelinks}

% ============================================================
% Palette « Pop » OnBuch+ — mêmes hex que la classe P (plus_ui.dart)
% ============================================================
\definecolor{poporange}{HTML}{F59321}
\definecolor{popdark}{HTML}{E07A0C}
\definecolor{popcream}{HTML}{FFF6E8}
\definecolor{popink}{HTML}{1C1714}
\definecolor{poppurple}{HTML}{7A5AE0}
\definecolor{popgreen}{HTML}{2E9E6A}
\definecolor{poppink}{HTML}{E0457B}
\definecolor{popblue}{HTML}{2F7DE1}
\definecolor{popgold}{HTML}{C98A12}
\definecolor{popmuted}{HTML}{8A7F76}
\definecolor{popline}{HTML}{EADFD2}
\definecolor{popgray}{HTML}{8A7F76}
\colorlet{popgrey}{popgray}
% Alias tolérants (le modèle nomme parfois les couleurs différemment)
\colorlet{popwhite}{white}
\colorlet{popblack}{popink}
\colorlet{popred}{poppink}
\colorlet{popyellow}{popgold}
% Teintes claires (fonds de tableaux/encadrés) — le modèle nomme parfois
% les couleurs avec un suffixe L (ex. popblueL) sans les avoir définies.
\colorlet{poporangeL}{poporange!20}
\colorlet{popdarkL}{popdark!20}
\colorlet{poppurpleL}{poppurple!20}
\colorlet{popgreenL}{popgreen!20}
\colorlet{poppinkL}{poppink!20}
\colorlet{popblueL}{popblue!20}
\colorlet{popgoldL}{popgold!20}
\colorlet{popredL}{popred!20}
\colorlet{popgrayL}{popgray!20}
% Teintes claires pour fonds de tableaux / zones
\colorlet{popblueL}{popblue!10!white}
\colorlet{poporangeL}{poporange!14!white}
\colorlet{popgreenL}{popgreen!12!white}
\colorlet{poppurpleL}{poppurple!12!white}
\colorlet{poppinkL}{poppink!12!white}
\colorlet{popgoldL}{popgold!14!white}

\pagecolor{popcream}

% ============================================================
% Typographie
% ============================================================
\defaultfontfeatures{Ligatures=TeX}
\setmainfont{PlusJakartaSans-Medium.ttf}[Path=../../../fonts/,BoldFont=PlusJakartaSans-Bold.ttf]
% Maths en sans-serif (Fira Math) avec chiffres et lettres droites de Plus
% Jakarta, pour que les formules se fondent dans le texte.
\usepackage[math-style=ISO,bold-style=ISO]{unicode-math}
\setmathfont{FiraMath-Regular.otf}[Path=../../../fonts/]
\setmathfont{PlusJakartaSans-Medium.ttf}[Path=../../../fonts/,range={up/num,up/{Latin,latin},\mathup}]
\usepackage{icomma} % virgule décimale sans espace en mode math
% Caractères mathématiques Unicode absents des polices de texte (Plus Jakarta,
% Archivo Black) : rendus par leur équivalent mathématique, en texte comme en
% formule — sinon ils s'affichent en carré vide (ex. « Arithmétique dans ℤ »).
\usepackage{newunicodechar}
\newunicodechar{ℝ}{\ensuremath{\mathbb{R}}}
\newunicodechar{ℤ}{\ensuremath{\mathbb{Z}}}
\newunicodechar{ℂ}{\ensuremath{\mathbb{C}}}
\newunicodechar{ℕ}{\ensuremath{\mathbb{N}}}
\newunicodechar{ℚ}{\ensuremath{\mathbb{Q}}}
\newunicodechar{𝔻}{\ensuremath{\mathbb{D}}}
\newunicodechar{α}{\ensuremath{\alpha}}
\newunicodechar{β}{\ensuremath{\beta}}
\newunicodechar{γ}{\ensuremath{\gamma}}
\newunicodechar{δ}{\ensuremath{\delta}}
\newunicodechar{ε}{\ensuremath{\varepsilon}}
\newunicodechar{θ}{\ensuremath{\theta}}
\newunicodechar{λ}{\ensuremath{\lambda}}
\newunicodechar{μ}{\ensuremath{\mu}}
\newunicodechar{σ}{\ensuremath{\sigma}}
\newunicodechar{τ}{\ensuremath{\tau}}
\newunicodechar{φ}{\ensuremath{\varphi}}
\newunicodechar{ψ}{\ensuremath{\psi}}
\newunicodechar{ω}{\ensuremath{\omega}}
\newunicodechar{Σ}{\ensuremath{\Sigma}}
% majuscules produites par \MakeUppercase dans les titres (α-aminés -> Α)
\newunicodechar{Α}{\ensuremath{\alpha}}
\newunicodechar{Β}{\ensuremath{\beta}}
\newunicodechar{Γ}{\ensuremath{\Gamma}}
\newunicodechar{Δ}{\ensuremath{\Delta}}
\newunicodechar{Ε}{\ensuremath{\varepsilon}}
\newunicodechar{Θ}{\ensuremath{\Theta}}
\newunicodechar{Λ}{\ensuremath{\Lambda}}
\newunicodechar{Μ}{\ensuremath{\mu}}
\newunicodechar{Τ}{\ensuremath{\tau}}
\newunicodechar{Φ}{\ensuremath{\Phi}}
\newunicodechar{Ψ}{\ensuremath{\Psi}}
\newunicodechar{Ω}{\ensuremath{\Omega}}
\newunicodechar{↦}{\ensuremath{\mapsto}}
\newunicodechar{∈}{\ensuremath{\in}}
\newunicodechar{∉}{\ensuremath{\notin}}
\newunicodechar{∧}{\ensuremath{\wedge}}
\newunicodechar{⇔}{\ensuremath{\Leftrightarrow}}
\newunicodechar{⟺}{\ensuremath{\Longleftrightarrow}}
\newunicodechar{⇒}{\ensuremath{\Rightarrow}}
\newunicodechar{⟹}{\ensuremath{\Longrightarrow}}
\newunicodechar{∘}{\ensuremath{\circ}}
\newunicodechar{∪}{\ensuremath{\cup}}
\newunicodechar{∩}{\ensuremath{\cap}}
\newunicodechar{≡}{\ensuremath{\equiv}}
\newunicodechar{⊂}{\ensuremath{\subset}}
\newunicodechar{⊥}{\ensuremath{\perp}}
\newunicodechar{∃}{\ensuremath{\exists}}
\newunicodechar{∀}{\ensuremath{\forall}}
\newunicodechar{∛}{\ensuremath{\sqrt[3]{\,}}}
\newunicodechar{∜}{\ensuremath{\sqrt[4]{\,}}}
\newunicodechar{⃗}{\ensuremath{{}^{\to}}}
\newunicodechar{ⁿ}{\ifmmode^{n}\else\textsuperscript{\ensuremath{n}}\fi}
\newunicodechar{ˣ}{\ifmmode^{x}\else\textsuperscript{\ensuremath{x}}\fi}
\newunicodechar{ᵇ}{\ifmmode^{b}\else\textsuperscript{\ensuremath{b}}\fi}
\newunicodechar{ʳ}{\ifmmode^{r}\else\textsuperscript{\ensuremath{r}}\fi}
\newunicodechar{ᵘ}{\ifmmode^{u}\else\textsuperscript{\ensuremath{u}}\fi}
\newunicodechar{ʸ}{\ifmmode^{y}\else\textsuperscript{\ensuremath{y}}\fi}
\newunicodechar{ᵃ}{\ifmmode^{a}\else\textsuperscript{\ensuremath{a}}\fi}
\newunicodechar{ᵉ}{\ifmmode^{e}\else\textsuperscript{\ensuremath{e}}\fi}
\newunicodechar{ᵏ}{\ifmmode^{k}\else\textsuperscript{\ensuremath{k}}\fi}
\newunicodechar{ᵖ}{\ifmmode^{p}\else\textsuperscript{\ensuremath{p}}\fi}
\newunicodechar{ᴺ}{\ifmmode^{N}\else\textsuperscript{\ensuremath{N}}\fi}
\newunicodechar{ᶜ}{\ifmmode^{c}\else\textsuperscript{\ensuremath{c}}\fi}
\newunicodechar{ʰ}{\ifmmode^{h}\else\textsuperscript{\ensuremath{h}}\fi}
\newunicodechar{ᵠ}{\ifmmode^{q}\else\textsuperscript{\ensuremath{q}}\fi}
\newunicodechar{ₙ}{\ifmmode_{n}\else\textsubscript{\ensuremath{n}}\fi}
\newunicodechar{ₐ}{\ifmmode_{a}\else\textsubscript{\ensuremath{a}}\fi}
\newunicodechar{ᵢ}{\ifmmode_{i}\else\textsubscript{\ensuremath{i}}\fi}
\newunicodechar{ᵣ}{\ifmmode_{r}\else\textsubscript{\ensuremath{r}}\fi}
\newunicodechar{ₖ}{\ifmmode_{k}\else\textsubscript{\ensuremath{k}}\fi}
\newunicodechar{ᵦ}{\ifmmode_{\beta}\else\textsubscript{\ensuremath{\beta}}\fi}
\newunicodechar{ₓ}{\ifmmode_{x}\else\textsubscript{\ensuremath{x}}\fi}
\newfontfamily\popdisplay{ArchivoBlack-Regular.ttf}[Path=../../../fonts/]
\newfontfamily\poplogo{SpaceGrotesk-Bold.ttf}[Path=../../../fonts/]
\newfontfamily\popsemi{PlusJakartaSans-SemiBold.ttf}[Path=../../../fonts/]
\newfontfamily\popxbold{PlusJakartaSans-ExtraBold.ttf}[Path=../../../fonts/]
\newfontfamily\popxboldls{PlusJakartaSans-ExtraBold.ttf}[Path=../../../fonts/,LetterSpace=3.0]

\setlist[enumerate]{leftmargin=1.7em,itemsep=5pt,topsep=4pt}
\setlist[itemize]{leftmargin=1.7em,itemsep=4pt,topsep=4pt,label={\color{poporange}\textbullet}}
\setlength{\parindent}{0pt}
\setlength{\parskip}{5pt}
\renewcommand{\arraystretch}{1.25}

% pgfplots : style Pop par défaut
\pgfplotsset{
  every axis/.append style={axis line style={popink,line width=1pt},tick style={popink},
    grid style={popline,line width=0.5pt},label style={font=\small},tick label style={font=\footnotesize}},
  popcurve/.style={poporange,line width=1.8pt,no markers,smooth},
}
% Même style disponible dans un \draw TikZ simple (hors pgfplots)
\tikzset{popcurve/.style={poporange,line width=1.8pt}}
\tikzset{popgrid/.style={popline,very thin}}
\colorlet{popcurve}{poporange} % le nom du style est parfois utilisé comme couleur

% ============================================================
% Pill / squiggle
% ============================================================
\newcommand{\poppill}[3]{%
  \tikz[baseline=-0.55ex]\node[draw=popink,line width=1.2pt,rounded corners=2.6pt,%
    fill=#2,inner xsep=6.5pt,inner ysep=3pt]{%
    \popxboldls\fontsize{7}{7}\selectfont\color{#3}\MakeUppercase{#1}};%
}
\newcommand{\popsquiggle}[1]{%
  \tikz[baseline=-0.4ex]{
    \draw[#1,line width=2.6pt,line cap=round]
      plot [smooth] coordinates {(0,0.09) (0.34,0.01) (0.68,0.21) (1.02,0.01) (1.36,0.10)};
  }%
}
% Mot-clé surligné (marqueur orange sous le texte)
\newcommand{\cle}[1]{{\popxbold\color{popdark}#1}}

% ============================================================
% En-tête / pied
% ============================================================
\pagestyle{fancy}
\fancyhf{}
\renewcommand{\headrulewidth}{0pt}
\renewcommand{\footrulewidth}{0pt}
\lhead{\raisebox{-0.15ex}{\poplogo\fontsize{13}{13}\selectfont\color{popink}OnBuch\color{poporange}+}}
\rhead{\poppill{\DOCMATIERE\ \textbullet\ \DOCNIVEAU\ \textbullet\ Leçon \DOCLECON}{white}{popink}}
\lfoot{\popxbold\fontsize{7.6}{7.6}\selectfont\color{popmuted}\MakeUppercase{Cours signé OnBuch+}\ \textbullet\ {\popsemi\DOCTITRE}}
\rfoot{\tikz[baseline=-0.6ex]\node[rounded corners=8.5pt,fill=popink,minimum height=17pt,minimum width=17pt,inner xsep=5pt,inner sep=0]{\popxbold\fontsize{8}{8}\selectfont\color{popcream}\thepage\,/\,\pageref*{LastPage}};}

% ============================================================
% Titres
% ============================================================
\newcommand{\chaptitre}[2]{%
  \begin{tcolorbox}[enhanced,colback=poporange,colframe=popink,boxrule=2.6pt,arc=11pt,
      drop shadow=popink,left=15pt,right=15pt,top=12pt,bottom=11pt,before skip=3pt,after skip=18pt]
    \poppill{#2}{popink}{popcream}\par\vspace{9pt}
    {\raggedright\hyphenpenalty=10000\exhyphenpenalty=10000\popdisplay\fontsize{22}{24}\selectfont\color{popcream}\MakeUppercase{#1}\par}
  \end{tcolorbox}
}
% Section numérotée : pastille orange avec le numéro + titre Archivo
\newcounter{popsec}
\newcommand{\coursec}[1]{%
  \stepcounter{popsec}\setcounter{popsub}{0}%
  \par\vspace{16pt}\noindent
  \begin{minipage}{\linewidth}
  \tikz[baseline=-0.7ex]\node[draw=popink,line width=1.6pt,fill=poporange,rounded corners=4pt,minimum size=22pt,inner sep=0]{\popdisplay\fontsize{12}{12}\selectfont\color{popcream}\thepopsec};%
  \hspace{8pt}{\popdisplay\fontsize{15}{17}\selectfont\color{popink}\MakeUppercase{#1}}\par
  \vspace{4pt}\noindent{\color{poporange}\rule{36pt}{3.6pt}}
  \end{minipage}\par\nopagebreak\vspace{8pt}
}
\newcounter{popsub}
\newcommand{\courssub}[1]{%
  \stepcounter{popsub}%
  \par\vspace{9pt}\noindent{\popxbold\fontsize{11.5}{13}\selectfont\color{popink}{\color{poporange}\thepopsec.\thepopsub}\hspace{6pt}#1}\par\nopagebreak\vspace{4pt}
}

% ============================================================
% Boîtes pédagogiques — même recette : fond blanc, bordure épaisse colorée,
% ombre dure encre, badge plein. Titre optionnel [#1].
% ============================================================
\tcbset{popbase/.style={breakable,enhanced,colback=white,boxrule=2.3pt,arc=9pt,
  shadow={3pt}{-3pt}{0pt}{popink},left=14pt,right=14pt,top=11pt,bottom=11pt,before skip=11pt,after skip=15pt,
  fontupper=\popsemi\fontsize{10.6}{14.5}\selectfont\color{popink},
  fontlower=\popsemi\fontsize{10.6}{14.5}\selectfont\color{popink}}}
% Coche dessinée (le glyphe ✓ n'existe pas dans les polices du document)
\newcommand{\popcheck}[1]{\tikz[baseline=-0.5ex]\draw[#1,line width=1.6pt,line cap=round,line join=round](-0.13,0.0)--(-0.04,-0.09)--(0.13,0.1);}
\providecommand{\coche}{\popcheck{popgreen}}
\providecommand{\resultat}[1]{\par\smallskip\noindent\fcolorbox{popgreen}{popgreenL}{\parbox{\dimexpr\linewidth-2\fboxsep-2\fboxrule\relax}{\textbf{Résultat.} #1}}\par\smallskip}
\newcommand{\popboxhead}[3]{\poppill{#1}{#2}{white}\ifx\relax#3\relax\else\hspace{6pt}{\popxbold\fontsize{10.6}{12}\selectfont\color{popink}#3}\fi\par\vspace{7pt}}

\newtcolorbox{aretenir}[1][]{popbase,colframe=popgreen,before upper={\popboxhead{À retenir}{popgreen}{#1}}}
\newtcolorbox{exemplebox}[1][]{popbase,colframe=poporange,before upper={\popboxhead{Exemple}{poporange}{#1}}}
\newtcolorbox{definition}[1][]{popbase,colframe=popblue,before upper={\popboxhead{Définition}{popblue}{#1}}}
\newtcolorbox{propriete}[1][]{popbase,colframe=poppurple,before upper={\popboxhead{Propriété}{poppurple}{#1}}}
\newtcolorbox{methode}[1][]{popbase,colframe=popgold,before upper={\popboxhead{Méthode}{popgold}{#1}}}
\newtcolorbox{attention}[1][]{popbase,colframe=poppink,before upper={\popboxhead{Attention}{poppink}{#1}}}
\newtcolorbox{experience}[1][]{popbase,colframe=popblue,colback=popblueL,before upper={\popboxhead{Expérience}{popblue}{#1}}}
% « Le savais-tu ? » : carte violette pleine (contexte, culture, Cameroun)
\newtcolorbox{savaistu}[1][]{popbase,colback=poppurple,colframe=popink,
  fontupper=\popsemi\fontsize{10.4}{14.2}\selectfont\color{white},
  before upper={\poppill{Le savais-tu\,?}{white}{poppurple}\ifx\relax#1\relax\else\hspace{6pt}{\popxbold\color{white}#1}\fi\par\vspace{7pt}}}
% Exercice résolu : énoncé en haut, \tcblower puis la solution sur fond crème
\newcounter{popexo}\newcounter{popexr}
\newtcolorbox{exoresolu}[1][]{popbase,colframe=poporange,colbacklower=popcream,
  segmentation style={popink,line width=1.2pt,dashed},
  before upper={\stepcounter{popexr}\popboxhead{Exercice résolu \thepopexr}{poporange}{#1}},
  before lower={\poppill{Solution}{popink}{popcream}\par\vspace{6pt}}}
% Exercice (non résolu en place) — la correction est dans la partie Corrigés
\newtcolorbox{exercice}[1][]{popbase,colframe=popink,
  before upper={\stepcounter{popexo}\popboxhead{Exercice \thepopexo}{popink}{#1}}}
% Corrigé
\newtcolorbox{corrige}[1][]{popbase,colframe=popgreen,colback=popgreenL,
  before upper={\popboxhead{Corrigé}{popgreen}{#1}}}
% Objectifs / prérequis / bilan
\newtcolorbox{objectifs}{popbase,colframe=popink,colback=poporangeL,
  before upper={\popboxhead{Objectifs de la leçon}{popink}{}}}
\newtcolorbox{prerequis}{popbase,colframe=popink,colback=popblueL,
  before upper={\popboxhead{Prérequis}{popblue}{}}}
\newtcolorbox{bilan}{popbase,colframe=popink,colback=white,boxrule=2.8pt,
  before upper={\popboxhead{Fiche bilan}{poporange}{L'essentiel à connaître pour l'examen}}}
% Figure encadrée avec légende
\newtcolorbox{popfigure}[1][]{enhanced,breakable=false,colback=white,colframe=popline,boxrule=1.4pt,arc=8pt,
  left=10pt,right=10pt,top=10pt,bottom=8pt,before skip=10pt,after skip=12pt,halign=center,
  lower separated=false,halign lower=center,
  fontlower=\popsemi\fontsize{9}{11.5}\selectfont\color{popmuted},
  #1}
\newcommand{\legende}[1]{\par\vspace{6pt}{\popsemi\fontsize{9}{11.5}\selectfont\color{popmuted}#1}}

% ============================================================
% Signature OnBuch+
% ============================================================
\newcommand{\poplogoinline}[1]{{\poplogo\fontsize{#1}{#1}\selectfont\color{popink}OnBuch\color{poporange}+}}
\newcommand{\signatureonbuch}{%
  \begin{tcolorbox}[enhanced,colback=popink,colframe=popink,boxrule=2.2pt,arc=10pt,
      drop shadow=poporange,left=14pt,right=14pt,top=10pt,bottom=10pt,before skip=0pt,after skip=0pt]
    \tikz[baseline=-0.7ex]\node[circle,fill=popgreen,draw=popcream,line width=1.3pt,minimum size=22pt,inner sep=0]{\popcheck{white}};%
    \hspace{9pt}\begin{minipage}[c]{0.62\linewidth}
      {\popxbold\fontsize{12}{13}\selectfont\color{popcream}Signé\ \textbullet\ {\poplogo OnBuch\color{poporange}+}}\par\vspace{2pt}
      {\raggedright\popsemi\fontsize{8}{10.5}\selectfont\color{popline}\MakeUppercase{Équipe pédagogique OnBuch+}\\\MakeUppercase{Conforme au programme officiel MINESEC}\par}
    \end{minipage}\hfill
    {\poplogo\fontsize{22}{22}\selectfont\color{popcream}OnBuch\color{poporange}+}
  \end{tcolorbox}%
}

% ============================================================
% Page de garde
% #1 titre · #2 matière · #3 niveau · #4 module · #5 n° leçon · #6 durée ·
% #7 description / objectifs (texte ou itemize)
% ============================================================
\newcommand{\pagegarde}[7]{%
  \thispagestyle{empty}
  \newgeometry{a4paper,margin=1.7cm,top=1.6cm,bottom=1.4cm}%
  \begin{tikzpicture}[remember picture,overlay]
    \fill[poporange!18] ([xshift=-3.2cm,yshift=-3.6cm]current page.north east) circle (4.6cm);
    \fill[poppurple!14] ([xshift=2.4cm,yshift=7.6cm]current page.south west) circle (3.8cm);
  \end{tikzpicture}%
  {\poplogo\fontsize{30}{30}\selectfont\color{popink}OnBuch\color{poporange}+}\hfill
  \raisebox{0.6ex}{\poppill{Cours signé \textbullet\ Premium}{popink}{popcream}}\par
  \vspace{1pt}\popsquiggle{poppurple}\par
  \vspace{8pt}
  \begin{tcolorbox}[enhanced,colback=poporange,colframe=popink,boxrule=2.8pt,arc=13pt,
      drop shadow=popink,left=18pt,right=18pt,top=12pt,bottom=13pt,after skip=10pt]
    \poppill{#2 \textbullet\ #3}{popink}{popcream}\hspace{5pt}\poppill{#4}{white}{popink}\par\vspace{8pt}
    {\popdisplay\fontsize{14}{14}\selectfont\color{popink}LEÇON #5}\par\vspace{4pt}
    {\raggedright\hyphenpenalty=10000\exhyphenpenalty=10000\popdisplay\fontsize{25}{27}\selectfont\color{popcream}\MakeUppercase{#1}\par}
    \vspace{8pt}
    {\popxbold\fontsize{9.5}{9.5}\selectfont\color{popink}\MakeUppercase{Durée conseillée : #6}}
  \end{tcolorbox}
  \begin{tcolorbox}[enhanced,colback=white,colframe=popink,boxrule=2.2pt,arc=10pt,
      drop shadow=popink,left=16pt,right=16pt,top=10pt,bottom=10pt,after skip=8pt]
    \poppill{Dans cette leçon}{poppurple}{white}\par\vspace{5pt}
    \popsemi\fontsize{9.3}{12.6}\selectfont\color{popink}#7
  \end{tcolorbox}
  \noindent\begin{minipage}[t]{0.487\linewidth}
    \begin{tcolorbox}[enhanced,colback=popgreen,colframe=popink,boxrule=2.2pt,arc=10pt,
        drop shadow=popink,left=11pt,right=11pt,top=10pt,bottom=10pt,height=3.1cm,valign=center]
      \tcbox[colback=white,colframe=popink,boxrule=1.3pt,arc=5pt,boxsep=0pt,nobeforeafter,
        left=3pt,right=3pt,top=3pt,bottom=3pt,valign=center]{\qrcode[height=1.9cm]{https://play.google.com/store/apps/details?id=com.luvvix.onbuch&hl=fr}}%
      \hspace{8pt}\begin{minipage}[c]{0.5\linewidth}
        {\popdisplay\fontsize{11}{12.5}\selectfont\color{white}\MakeUppercase{Télécharge OnBuch}}\par\vspace{4pt}
        {\raggedright\popsemi\fontsize{8.4}{11}\selectfont\color{white}Gratuit \textbullet\ Google Play\\Tuteur IA, annales, concours, résultats\par}
      \end{minipage}
    \end{tcolorbox}
  \end{minipage}\hfill
  \begin{minipage}[t]{0.487\linewidth}
    \begin{tcolorbox}[enhanced,colback=poppurple,colframe=popink,boxrule=2.2pt,arc=10pt,
        drop shadow=popink,left=11pt,right=11pt,top=10pt,bottom=10pt,height=3.1cm,valign=center]
      \tikz[baseline=-0.7ex]\node[circle,fill=white,draw=popink,line width=1.3pt,minimum size=1.6cm,inner sep=0]{\popdisplay\fontsize{24}{24}\selectfont\color{poppurple}+};%
      \hspace{8pt}\begin{minipage}[c]{0.6\linewidth}
        {\popdisplay\fontsize{11}{12.5}\selectfont\color{white}\MakeUppercase{Passe à OnBuch+}}\par\vspace{4pt}
        {\raggedright\popsemi\fontsize{8.4}{11}\selectfont\color{white}Tous les cours signés, Tuteur IA illimité, fiches PDF — onglet OnBuch+ de l'app\par}
      \end{minipage}
    \end{tcolorbox}
  \end{minipage}
  \vspace{8pt}
  \popcontenu
  \vfill
  \signatureonbuch
  \restoregeometry
  \newpage
}

% Rangée « Ce cours contient » (page de garde)
\newcommand{\poptile}[3]{% #1 couleur #2 gros texte #3 légende
  \begin{tcolorbox}[enhanced,colback=#1,colframe=popink,boxrule=2pt,arc=9pt,shadow={2.5pt}{-2.5pt}{0pt}{popink},
      left=6pt,right=6pt,top=9pt,bottom=9pt,halign=center,valign=center,height=1.75cm,width=\linewidth,nobeforeafter]
    {\popdisplay\fontsize{12.5}{13}\selectfont\color{white}\MakeUppercase{#2}}\par\vspace{3pt}
    {\centering\popsemi\fontsize{7.8}{9.5}\selectfont\color{white}#3\par}
  \end{tcolorbox}}
\newcommand{\popcontenu}{%
  \par\noindent{\popxboldls\fontsize{7.5}{7.5}\selectfont\color{popmuted}CE COURS SIGNÉ CONTIENT}\par\vspace{7pt}
  \noindent\begin{minipage}[t]{0.235\linewidth}\poptile{popblue}{Cours}{complet, illustré, pas à pas}\end{minipage}\hfill
  \begin{minipage}[t]{0.235\linewidth}\poptile{poporange}{Méthodes}{et exercices résolus}\end{minipage}\hfill
  \begin{minipage}[t]{0.235\linewidth}\poptile{poppink}{Exercices}{type examen + corrigés}\end{minipage}\hfill
  \begin{minipage}[t]{0.235\linewidth}\poptile{popgreen}{Fiche bilan}{l'essentiel à retenir}\end{minipage}\par}

% Sommaire
\newtcolorbox{sommaire}{enhanced,colback=white,colframe=popink,boxrule=2.2pt,arc=10pt,
  drop shadow=popink,left=18pt,right=18pt,top=13pt,bottom=13pt,before skip=6pt,after skip=20pt,
  before upper={\poppill{Sommaire}{poppurple}{white}\par\vspace{9pt}\popsemi\fontsize{10.6}{14.5}\selectfont\color{popink}}}

% Signature de fin de cours — poussée tout en bas de la dernière page
\newcommand{\findecours}{%
  \par\vfill\nopagebreak
  \begin{center}\popsquiggle{poporange}\end{center}\vspace{4pt}
  \signatureonbuch
}
\AtBeginDocument{\def\llbracket{\mathopen{[\mkern-3.5mu[}}\def\rrbracket{\mathclose{]\mkern-3.5mu]}}} % intervalles d'entiers (glyphe ⟦ absent de la police)
% Glyphes absents de Fira Math : redéfinis à partir de glyphes disponibles
\AtBeginDocument{%
  \def\setminus{\mathbin{\backslash}}\let\smallsetminus\setminus
  \def\vdots{\mathord{\vbox{\baselineskip3.2pt\lineskiplimit0pt\kern4pt\hbox{.}\hbox{.}\hbox{.}}}}%
  \def\ddots{\mathinner{\mkern1mu\raise6.5pt\hbox{.}\mkern2mu\raise3.6pt\hbox{.}\mkern2mu\raise0.7pt\hbox{.}\mkern1mu}}%
  \def\triangle{\mathord{\scalebox{0.9}{$\symup{\Delta}$}}}%
  \def\nabla{\mathord{\raisebox{\depth}{\scalebox{1}[-1]{$\symup{\Delta}$}}}}%
  \def\checkmark{\popcheck{popdark}}%
  \def\star{\mathbin{\ast}}\def\bigstar{\mathord{\ast}}%
  \def\diamond{\mathbin{\scalebox{0.6}{$\Diamond$}}}%
  \def\bigcup{\mathop{\mathchoice{\raisebox{-0.3em}{\scalebox{1.7}{$\cup$}}}{\scalebox{1.25}{$\cup$}}{\cup}{\cup}}\displaylimits}%
  \def\bigcap{\mathop{\mathchoice{\raisebox{-0.3em}{\scalebox{1.7}{$\cap$}}}{\scalebox{1.25}{$\cap$}}{\cap}{\cap}}\displaylimits}%
  \def\lesseqqgtr{\mathrel{\lessgtr}}\def\bigoplus{\mathop{\oplus}\displaylimits}%
}
\providecommand{\ssi}{si et seulement si} % abréviation fréquente
\AtBeginDocument{\providecommand{\wideparen}{\overparen}} % arc de cercle

%% FIGFIX-BEGIN v4 — correctifs globaux des figures (docs/onbuch-generation/tools/figfix)
\usepackage{adjustbox}\usepackage{etoolbox}
% toute figure TikZ/pgfplots plus large que la ligne est réduite à la largeur disponible
\BeforeBeginEnvironment{tikzpicture}{\begin{adjustbox}{max width=\linewidth}}
\AfterEndEnvironment{tikzpicture}{\end{adjustbox}}
% légendes pgfplots lisibles par-dessus les courbes
\pgfplotsset{every axis/.append style={legend style={fill=white,fill opacity=0.92,text opacity=1,draw=black!15,inner sep=3pt}}}
% étiquettes d'arêtes (syntaxe edge["..."]) avec fond blanc
\tikzset{every edge quotes/.append style={fill=white,inner sep=1.2pt}}
%% FIGFIX-END
\begin{document}

\pagegarde{Contraction de texte : rédiger un résumé de texte}{Français}{Tle Tech}{Français – classe de Terminale technique}{N10}{15 séances (fiche de progression harmonisée)}{Cette leçon outille l'élève de Terminale technique pour maîtriser la contraction de texte : repérage de la structure argumentative, techniques de réduction et de reformulation, élaboration du plan du résumé, rédaction, confrontation et amélioration des productions. Elle s'appuie sur l'étude et le bilan du Vieux nègre et la médaille de Ferdinand Oyono et consolide les notions de cohésion, de cohérence, de progression thématique et de tonalités satirique et polémique.
\vspace{4pt}
\begin{multicols}{2}\small
\begin{itemize}
\item À la fin de cette leçon, je sais identifier la structure argumentative d'un texte (thèse, arguments, exemples, conclusion).
\item À la fin de cette leçon, je sais distinguer l'essentiel de l'accessoire dans un texte explicatif ou argumentatif.
\item À la fin de cette leçon, je sais appliquer les techniques de contraction : suppression, substitution, fusion, généralisation.
\item À la fin de cette leçon, je sais élaborer le plan d'un texte et rédiger un résumé respectant le quart du volume initial.
\item À la fin de cette leçon, je sais garantir la cohésion et la cohérence de mon résumé grâce aux connecteurs et à la progression thématique.
\item À la fin de cette leçon, je sais reconnaître et restituer les tonalités satirique et polémique d'un texte sans les dénaturer.
\end{itemize}\end{multicols}}

\begin{sommaire}
\begin{multicols}{2}
\begin{enumerate}
\item Le texte explicatif et la notion de progression
\item Cohésion, cohérence et sigles
\item Le Vieux nègre et la médaille : contexte, tonalités et bilan de l'œuvre
\item La structure argumentative du texte
\item Techniques de rédaction du résumé
\item Plan du texte, rédaction, confrontation et amélioration
\item Activité d'intégration
\item Méthodes et exercices résolus
\item Exercices
\item Corrigés des exercices
\item Fiche bilan
\end{enumerate}
\end{multicols}
\end{sommaire}

\begin{objectifs}
\begin{itemize}
\item À la fin de cette leçon, je sais identifier la structure argumentative d'un texte (thèse, arguments, exemples, conclusion).
\item À la fin de cette leçon, je sais distinguer l'essentiel de l'accessoire dans un texte explicatif ou argumentatif.
\item À la fin de cette leçon, je sais appliquer les techniques de contraction : suppression, substitution, fusion, généralisation.
\item À la fin de cette leçon, je sais élaborer le plan d'un texte et rédiger un résumé respectant le quart du volume initial.
\item À la fin de cette leçon, je sais garantir la cohésion et la cohérence de mon résumé grâce aux connecteurs et à la progression thématique.
\item À la fin de cette leçon, je sais reconnaître et restituer les tonalités satirique et polémique d'un texte sans les dénaturer.
\item À la fin de cette leçon, je sais situer Le Vieux nègre et la médaille dans son contexte colonial et évaluer mon étude de l'œuvre.
\item À la fin de cette leçon, je sais comparer, corriger et améliorer ma production et celle de mes camarades à l'aide d'une grille critériée.
\item À la fin de cette leçon, je sais rédiger et justifier une démarche, une réponse ou une conclusion en situation d'examen.
\end{itemize}
\end{objectifs}

\begin{prerequis}
\begin{itemize}
\item Les types de textes (narratif, descriptif, explicatif, argumentatif) vus en classe de Première.
\item Les connecteurs logiques et les reprises anaphoriques (cohésion textuelle) vus en début d'année.
\item La lecture intégrale du Vieux nègre et la médaille de Ferdinand Oyono (œuvre au programme).
\item Le schéma narratif et le schéma actantiel (classes de Seconde et Première).
\item Les figures de style de base : ironie, antithèse, hyperbole, comparaison.
\end{itemize}
\end{prerequis}

\newpage

\coursec{Le texte explicatif et la notion de progression}

\courssub{Situation d'accroche et problématique}

À la fin de l'année, l'épreuve de Français du Baccalauréat technique vous demandera de résumer un texte de 500 mots en 125 mots environ, sans trahir l'auteur ni recopier ses phrases. Imaginez un journaliste de la CRTV qui doit rendre compte en trois minutes d'un discours ministériel d'une heure, ou un élève de Yaoundé qui doit transmettre par SMS l'essentiel d'une circulaire du MINESEC à ses camarades de village~: comment dire l'essentiel, fidèlement, en peu de mots~? C'est exactement le défi de la \cle{contraction de texte}.

Mais avant de compresser un texte, encore faut-il comprendre comment il est construit. Un texte n'est pas un sac de phrases jetées en désordre~: c'est une machine organisée, où chaque phrase s'appuie sur la précédente pour faire avancer le sens. Celui qui sait repérer cette mécanique sait aussi distinguer l'essentiel de l'accessoire — et donc résumer sans trahir.

\textbf{Problématique~:} comment un texte transmet-il son information de phrase en phrase, et comment la connaissance de ce fonctionnement nous aide-t-elle à en dégager les idées essentielles~?

Pour répondre à cette question, nous étudierons d'abord le \cle{texte explicatif}, genre de texte très fréquent dans les sujets de résumé, puis la \cle{progression thématique}, c'est-à-dire la manière dont l'information circule dans un texte.

\courssub{Le texte explicatif~: informer sans convaincre}

\textbf{Intuition.} Lisez ces deux phrases~: « Le paludisme se transmet par la piqûre d'un moustique femelle infecté » et « Il faut absolument que chaque famille camerounaise dorme sous moustiquaire imprégnée ». La première vous donne une information, la seconde cherche à vous faire agir. Vous sentez spontanément la différence~: c'est celle qui sépare le texte explicatif du texte argumentatif.

\begin{definition}[Le texte explicatif]
Un \cle{texte explicatif} est un texte dont la fonction principale est d'\textbf{informer}, d'\textbf{expliquer} et de \textbf{faire comprendre} un phénomène, un objet, une notion ou un fonctionnement. L'énonciateur reste en retrait~: il ne cherche ni à \textbf{convaincre} (faire adhérer à une opinion par le raisonnement) ni à \textbf{persuader} (toucher les émotions pour faire agir). Il transmet un savoir présenté comme objectif et vérifiable.
\end{definition}

On trouve des textes explicatifs partout dans la vie quotidienne et scolaire~: notices d'utilisation d'un téléphone ou d'une machine agricole, articles encyclopédiques, manuels scolaires, comptes rendus d'expériences, articles de presse d'information (à distinguer des éditoriaux, qui argumentent), circulaires administratives du MINESEC, modes d'emploi des médicaments.

\begin{propriete}[Caractéristiques du texte explicatif]
Un texte explicatif se reconnaît à plusieurs indices linguistiques~:
\begin{itemize}
\item un \textbf{lexique précis, technique et dénotatif}~: les mots sont pris dans leur sens propre, premier, sans intention de faire rêver ni d'émouvoir (on dit « le moteur à explosion comporte quatre temps », non « le cœur palpitant de la machine »)~;
\item l'emploi du \textbf{présent de vérité générale} (ou présent gnomique)~: « L'eau bout à 100 degrés », « La photosynthèse transforme l'énergie lumineuse en énergie chimique »~; ce présent affirme des faits valables en tout temps~;
\item la présence de \textbf{définitions} (souvent introduites par \emph{c'est-à-dire, on appelle, cela désigne}) et d'\textbf{exemples} (introduits par \emph{par exemple, ainsi, c'est le cas de})~;
\item des \textbf{connecteurs d'articulation logique} qui organisent l'explication~: \emph{d'abord, ensuite, puis, enfin} (ordre chronologique ou ordinal)~; \emph{en effet, car, parce que} (cause)~; \emph{c'est pourquoi, ainsi, donc} (conséquence)~; \emph{mais, cependant, en revanche} (opposition)~;
\item une \textbf{énonciation effacée}~: le « je » de l'auteur disparaît au profit de tournures impersonnelles (\emph{il est, on constate, il convient de}) ou du « on » général.
\end{itemize}
\end{propriete}

\begin{exemplebox}[Un court texte explicatif]
« Le cacao est la première culture d'exportation du Cameroun. On appelle \emph{cabosse} le fruit du cacaoyer. Celle-ci contient des fèves qui, après fermentation et séchage, sont vendues aux unités de transformation. D'abord fermentées pendant environ six jours, les fèves sont ensuite séchées au soleil, puis enfin ensachées pour l'exportation. »

On observe~: lexique technique dénotatif (\emph{cabosse, fermentation, fèves}), présent de vérité générale (\emph{est, contient, sont vendues}), une définition (\emph{on appelle cabosse...}), des connecteurs ordinaux (\emph{d'abord, ensuite, puis enfin}) qui rythment l'explication du procédé. Aucune opinion n'est défendue~: le texte informe, c'est tout.
\end{exemplebox}

\courssub{Texte explicatif et texte argumentatif~: ne pas confondre}

La distinction est capitale pour le résumé~: si vous résumez un texte explicatif comme s'il défendait une thèse, vous trahissez l'auteur~; si vous résumez un texte argumentatif en oubliant sa thèse, vous passez à côté de l'essentiel.

\begin{center}
\begin{tabularx}{\linewidth}{|l|X|X|}
\hline
\rowcolor{popblueL} \textbf{Critère} & \textbf{Texte explicatif} & \textbf{Texte argumentatif} \\
\hline
\textbf{Intention} & Informer, expliquer, faire comprendre & Convaincre ou persuader, défendre une thèse \\
\hline
\textbf{Position de l'auteur} & Neutre, effacée, objective & Engagée~: il prend parti \\
\hline
\textbf{Lexique} & Précis, technique, dénotatif & Souvent connotatif, valorisant ou péjoratif, champs lexicaux de l'opinion \\
\hline
\textbf{Temps verbaux} & Présent de vérité générale & Présent de l'énonciation, futur, conditionnel, impératifs possibles \\
\hline
\textbf{Connecteurs} & \emph{d'abord, ensuite, enfin, c'est-à-dire, par exemple} & \emph{donc, or, en effet, certes... mais, par conséquent} (au service d'une démonstration) \\
\hline
\textbf{Marques typiques} & Définitions, exemples, chiffres, procédés décrits & Thèse, arguments, exemples-preuves, concession, réfutation \\
\hline
\textbf{Question à se poser} & « Qu'est-ce que ce texte m'apprend~? » & « Que veut me faire penser ou faire ce texte~? » \\
\hline
\end{tabularx}
\end{center}

\begin{attention}[Piège fréquent~: confondre explicatif et argumentatif]
L'erreur la plus grave au Bac consiste à lire un texte neutre comme un texte engagé, ou l'inverse. Deux réflexes de sécurité~:
\begin{enumerate}
\item Cherchez une \textbf{thèse} (une opinion que l'auteur veut vous faire partager). S'il n'y en a pas, le texte est explicatif ou informatif.
\item Méfiez-vous des mots-outils~: \emph{en effet} et \emph{donc} apparaissent dans les deux types de textes. Ce n'est pas le connecteur qui décide, c'est l'\textbf{intention globale}. « L'eau bout à 100 degrés, donc la vapeur peut cuire les aliments » explique~; « La peine de mort est inefficace, donc il faut l'abolir » argumente.
\end{enumerate}
Dans un résumé, restituez toujours le \textbf{type de texte} tel qu'il est~: ne transformez jamais une explication en plaidoyer.
\end{attention}

\courssub{La progression thématique~: thème et rhème}

\textbf{Intuition.} Quand vous racontez un match de football à un ami, vous ne partez pas de zéro à chaque phrase~: « Les Lions ont dominé la première période. \emph{Ils} ont marqué deux buts. \emph{Le premier} a été inscrit par Aboubakar... » Chaque phrase reprend un élément déjà connu (\emph{ils, le premier}) et ajoute une information nouvelle. Cette mécanique — du connu vers le nouveau —, c'est la progression thématique.

\begin{definition}[Thème, rhème, progression thématique]
Dans une phrase, on distingue~:
\begin{itemize}
\item le \cle{thème}~: ce dont on parle, l'information \textbf{déjà connue} ou donnée par le contexte, le point de départ de la phrase (souvent placé au début)~;
\item le \cle{rhème}~: ce qu'on en dit, l'information \textbf{nouvelle}, l'apport propre de la phrase.
\end{itemize}
La \cle{progression thématique} est l'\textbf{enchaînement des thèmes et des rhèmes de phrase en phrase}, qui fait avancer le sens du texte. C'est le fil conducteur qui relie les informations entre elles et assure la cohérence du propos.
\end{definition}

\begin{savaistu}[D'où vient cette notion~?]
L'analyse de la phrase en thème et rhème est née au sein de l'\textbf{École de Prague} (Cercle linguistique de Prague), dans les années 1920-1930. Le linguiste tchèque \textbf{Vilém Mathesius} a le premier décrit ce qu'il appelait la « perspective actuelle de la phrase »~: toute phrase organise son information en partant d'un donné (le thème) vers un nouveau (le rhème). Cette découverte, vieille d'un siècle, est aujourd'hui un outil précieux pour lire vite et résumer fidèlement~: les rhèmes successifs portent l'essentiel du contenu d'un texte.
\end{savaistu}

\courssub{Les trois types de progression}

Selon la manière dont le thème de chaque phrase est choisi, on distingue trois types de progression.

\textbf{1. La progression linéaire (ou en chaîne).} Le rhème de la première phrase devient le thème de la suivante, et ainsi de suite~: le texte avance comme une chaîne, maillon après maillon. C'est la progression typique du récit et des explications de procédés.

\emph{Exemple~:} « Meka reçut une médaille. Cette \textbf{médaille}, il la montra à tout le village. Le \textbf{village} organisait une grande fête en son honneur. » (médaille $\rightarrow$ médaille~; village $\rightarrow$ village).

\textbf{2. La progression à thème constant.} Toutes les phrases parlent de la même chose~: le thème reste identique (souvent repris par un pronom), seuls les rhèmes changent. C'est la progression de la description, du portrait, de la définition.

\emph{Exemple~:} « Le vieux Meka était un homme simple. \textbf{Il} avait travaillé toute sa vie pour les Blancs. \textbf{Il} ne savait ni lire ni écrire. \textbf{Il} croyait que la médaille allait changer sa vie. »

\textbf{3. La progression à thèmes dérivés (ou en éventail).} Les thèmes des phrases successives dérivent tous d'un thème général annoncé au début (l'hyperthème), comme les rayons d'un éventail. C'est la progression des textes qui énumèrent les aspects d'une notion.

\emph{Exemple~:} « La cérémonie fut grandiose. \textbf{Le discours du commandant} dura une heure. \textbf{Les danses traditionnelles} suivirent. \textbf{Le banquet}, enfin, rassembla tout le village. » (discours, danses, banquet = aspects dérivés de « la cérémonie »).

\begin{popfigure}
\begin{center}
\begin{tikzpicture}[scale=0.92, transform shape,
  theme/.style={circle, draw=popblue, fill=popblueL, minimum size=0.72cm, font=\small\bfseries},
  rheme/.style={circle, draw=poporange, fill=poporangeL, minimum size=0.72cm, font=\small},
  fleche/.style={-{Stealth[length=2.2mm]}, thick, popdark},
  titre/.style={font=\small\bfseries, popdark, anchor=west}]
% --- Progression linéaire ---
\node[titre] at (-0.4,3.4) {1. Progression linéaire (en chaîne)};
\node[theme] (T1) at (0.8,2.6) {T1};
\node[rheme] (R1) at (2.3,2.6) {R1};
\node[theme] (T2) at (3.8,2.6) {T2};
\node[rheme] (R2) at (5.3,2.6) {R2};
\node[theme] (T3) at (6.8,2.6) {T3};
\node[rheme] (R3) at (8.3,2.6) {R3};
\draw[fleche] (T1) -- (R1);
\draw[fleche] (T2) -- (R2);
\draw[fleche] (T3) -- (R3);
\draw[fleche, popgreen, dashed] (R1) to[bend left=18] (T2);
\draw[fleche, popgreen, dashed] (R2) to[bend left=18] (T3);
% --- Progression à thème constant ---
\node[titre] at (-0.4,1.3) {2. Progression à thème constant};
\node[theme] (Tc1) at (0.8,0.4) {T};
\node[rheme] (Rc1) at (2.3,0.4) {R1};
\node[theme] (Tc2) at (3.8,0.4) {T};
\node[rheme] (Rc2) at (5.3,0.4) {R2};
\node[theme] (Tc3) at (6.8,0.4) {T};
\node[rheme] (Rc3) at (8.3,0.4) {R3};
\draw[fleche] (Tc1) -- (Rc1);
\draw[fleche] (Tc2) -- (Rc2);
\draw[fleche] (Tc3) -- (Rc3);
\draw[fleche, popgreen, dashed] (Tc1) to[bend left=18] (Tc2);
\draw[fleche, popgreen, dashed] (Tc2) to[bend left=18] (Tc3);
% --- Progression à thèmes dérivés ---
\node[titre] at (-0.4,-0.9) {3. Progression à thèmes dérivés (en éventail)};
\node[theme, fill=poppurpleL, draw=poppurple] (HT) at (0.8,-2.1) {HT};
\node[theme] (Td1) at (3.0,-1.5) {T1};
\node[theme] (Td2) at (3.0,-2.1) {T2};
\node[theme] (Td3) at (3.0,-2.7) {T3};
\node[rheme] (Rd1) at (5.2,-1.5) {R1};
\node[rheme] (Rd2) at (5.2,-2.1) {R2};
\node[rheme] (Rd3) at (5.2,-2.7) {R3};
\draw[fleche, poppurple, dashed] (HT) -- (Td1);
\draw[fleche, poppurple, dashed] (HT) -- (Td2);
\draw[fleche, poppurple, dashed] (HT) -- (Td3);
\draw[fleche] (Td1) -- (Rd1);
\draw[fleche] (Td2) -- (Rd2);
\draw[fleche] (Td3) -- (Rd3);
% Légende
\node[theme, minimum size=0.5cm] at (7.4,-1.5) {};
\node[anchor=west, font=\small] at (7.8,-1.5) {thème (connu)};
\node[rheme, minimum size=0.5cm] at (7.4,-2.1) {};
\node[anchor=west, font=\small] at (7.8,-2.1) {rhème (nouveau)};
\node[anchor=west, font=\small, text=poppurple] at (7.4,-2.7) {HT : hyperthème};
\end{tikzpicture}
\end{center}
\legende{Les trois types de progression thématique. En 1, chaque rhème devient le thème suivant (chaîne). En 2, le thème reste identique et les rhèmes se succèdent (description). En 3, les thèmes dérivent d'un hyperthème annoncé au départ (énumération). Les flèches pleines relient thème et rhème à l'intérieur d'une phrase~; les flèches en pointillés montrent la reprise d'information d'une phrase à l'autre.}
\end{popfigure}

\begin{aretenir}[L'essentiel sur la progression thématique]
\begin{itemize}
\item Le \textbf{thème} = ce dont on parle (connu) ; le \textbf{rhème} = ce qu'on en dit (nouveau).
\item Trois types~: \textbf{linéaire} (R devient T), \textbf{à thème constant} (même T partout), \textbf{à thèmes dérivés} (les T dérivent d'un hyperthème).
\item Un texte réel combine souvent plusieurs types de progression.
\item Les rhèmes portent l'information nouvelle~: ce sont eux qu'il faut conserver dans un résumé~; les reprises thématiques (pronoms, répétitions) peuvent être supprimées sans perte.
\end{itemize}
\end{aretenir}

\courssub{Repérer la progression dans un texte}

\begin{methode}[Repérer la progression thématique d'un texte]
\begin{enumerate}
\item Lisez le texte phrase par phrase et \textbf{soulignez le thème} de chaque phrase~: demandez-vous « de quoi cette phrase part-elle~? » (souvent le sujet grammatical, ou un pronom, ou un groupe nominal repris).
\item Encadrez le \textbf{rhème}~: « quelle information nouvelle cette phrase apporte-t-elle~? »
\item Pour chaque thème, vérifiez \textbf{d'où il vient}~: est-il repris du rhème précédent (progression linéaire)~? est-il identique au thème précédent (thème constant)~? renvoie-t-il à une notion générale annoncée au début (thème dérivé)~?
\item Concluez sur le type dominant de progression et sur ce qu'il révèle~: un récit (linéaire), un portrait (thème constant), un exposé (thèmes dérivés).
\item Listez les \textbf{rhèmes successifs}~: vous obtenez le squelette des idées essentielles, prêt pour le résumé.
\end{enumerate}
\end{methode}

\begin{exemplebox}[La remise de la médaille chez Ferdinand Oyono]
Dans \emph{Le Vieux nègre et la médaille}, la scène de la remise de la médaille à Meka peut se lire ainsi (analyse simplifiée)~:

« \textbf{Meka} [T1] s'avança vers l'estrade, tremblant d'émotion [R1]. \textbf{Le commandant} [T2] lui accrocha la médaille sur la poitrine [R2]. \textbf{La médaille} [T3] brillait au soleil comme une promesse [R3]. \textbf{Elle} [T4] allait faire de lui un homme important, croyait-il [R4]. »

Analyse~: T2 (« le commandant ») est un élément nouveau introduit par la scène~; mais T3 reprend le rhème R2 (« la médaille » $\rightarrow$ « la médaille ») et T4 reprend T3 par le pronom « elle ». Nous avons donc une \textbf{progression linéaire} centrée sur la médaille~: le texte enchaîne médaille $\rightarrow$ médaille $\rightarrow$ elle. Ce fil montre que la médaille est le véritable centre du passage — et le résumé devra la conserver comme idée essentielle, avec l'illusion qu'elle représente pour Meka.
\end{exemplebox}

\begin{exoresolu}[Repérage de la progression dans un article de presse]
Énoncé. Voici un extrait adapté d'un article d'information du type de ceux que publie \emph{Cameroon Tribune}~:

« Le gouvernement a lancé lundi un vaste programme de réhabilitation des routes rurales. Ce programme concerne dix régions du pays. Il devrait faciliter l'évacuation des produits agricoles vers les villes. Les produits agricoles, en effet, perdent souvent de leur valeur faute de pistes praticables. »

1) Soulignez le thème de chaque phrase. 2) Identifiez le type de progression. 3) Dégagez les idées essentielles pour un résumé.
\tcblower
Solution détaillée.
\begin{enumerate}
\item Thèmes~: phrase 1~: « \textbf{Le gouvernement} » ; phrase 2~: « \textbf{Ce programme} » ; phrase 3~: « \textbf{Il} » (= le programme) ; phrase 4~: « \textbf{Les produits agricoles} ».
\item Analyse des reprises~: « Ce programme » (T2) reprend le rhème de la phrase 1 (« un vaste programme ») ; « Il » (T3) reprend T2 ; « Les produits agricoles » (T4) reprend le rhème de la phrase 3. Chaque rhème devient le thème suivant~: c'est une \textbf{progression linéaire} (en chaîne), typique d'un article d'information qui déroule les conséquences d'un fait.
\item Idées essentielles (les rhèmes)~: lancement d'un programme de réhabilitation des routes rurales ; extension à dix régions ; objectif~: évacuer les produits agricoles ; cause~: la déperdition due au mauvais état des pistes. Résumé possible en une phrase~: « Le gouvernement a lancé un programme de réhabilitation des routes rurales dans dix régions afin de faciliter l'évacuation des produits agricoles, souvent dévalorisés faute de pistes praticables. »
\end{enumerate}
On remarque que les reprises thématiques (« ce programme », « il ») ont disparu du résumé~: elles assuraient la liaison, pas le contenu. C'est exactement le travail de contraction.
\end{exoresolu}

\courssub{Pourquoi la progression est-elle utile pour résumer~?}

La progression thématique n'est pas un exercice scolaire gratuit~: c'est la clé de voûte de la contraction de texte, pour trois raisons.

\begin{itemize}
\item \textbf{Elle révèle la hiérarchie des idées.} Ce qui revient comme thème est l'ossature du texte~; ce qui n'apparaît qu'une fois, dans un rhème secondaire (exemple, détail, anecdote), est souvent supprimable.
\item \textbf{Elle garantit la fidélité.} En suivant la chaîne thème-rhème, vous respectez l'ordre logique de l'auteur~: votre résumé ne mélange pas les idées, il les condense.
\item \textbf{Elle guide la reformulation.} Les reprises (pronoms, synonymes, répétitions) sont autant de mots économisables~: au lieu de répéter « le programme... ce programme... il... », le résumé dit une seule fois « le programme ».
\end{itemize}

\begin{attention}[Erreurs fréquentes dans le repérage]
\begin{itemize}
\item Ne confondez pas le \textbf{thème} (ce dont on parle, fonction textuelle) avec le \textbf{sujet grammatical}~: ils coïncident souvent, mais pas toujours (« Quant à la médaille, Meka la garda précieusement »~: le thème est « la médaille », le sujet est « Meka »).
\item Ne concluez pas sur un seul couple de phrases~: un type de progression se vérifie sur l'ensemble du paragraphe.
\item Un texte peut \textbf{alterner} les types~: annoncez le type \emph{dominant} et signalez les variations si elles sont significatives.
\end{itemize}
\end{attention}

\begin{exercice}[À vous de jouer]
Relevez les thèmes et les rhèmes du passage suivant, puis indiquez le type de progression dominant~:

« Le lycée technique de la ville possède un nouvel atelier. Cet atelier abrite des machines-outils modernes. Ces machines permettent aux élèves de se former aux métiers de l'industrie. La formation industrielle est aujourd'hui une priorité nationale. »

Proposez ensuite un résumé du passage en une seule phrase.
\end{exercice}

\begin{corrige}[Exercice 1]
Thèmes~: « Le lycée technique » (T1) ; « Cet atelier » (T2, repris du rhème R1 « un nouvel atelier ») ; « Ces machines » (T3, repris de R2 « des machines-outils ») ; « La formation industrielle » (T4, dérivé de R3 « se former aux métiers de l'industrie »). Chaque thème reprend le rhème précédent~: \textbf{progression linéaire}. Résumé possible~: « Le nouvel atelier du lycée technique, équipé de machines-outils modernes, forme les élèves aux métiers de l'industrie, secteur désormais prioritaire. »
\end{corrige}

En résumé, cette première étape de la leçon vous a donné deux outils complémentaires~: savoir \textbf{reconnaître un texte explicatif} (informer sans convaincre, lexique dénotatif, présent de vérité générale) et savoir \textbf{suivre le fil thème-rhème} qui tisse le texte. La section suivante approfondira la cohésion et la cohérence, c'est-à-dire les autres coutures invisibles qui tiennent un texte ensemble.

\coursec{Cohésion, cohérence et sigles}

Après avoir analysé le texte explicatif et ses progressions, nous devons désormais ausculter la « mécanique interne » de tout texte bien formé. Un résumé réussi ne se contente pas de condenser l'information ; il doit restituer la \textbf{texture} du texte source. Cette texture repose sur deux piliers indissociables : la \cle{cohésion} (les coutures visibles entre les phrases) et la \cle{cohérence} (l'architecture logique invisible qui tient l'édifice). Parallèlement, la maîtrise des \cle{sigles} — omniprésents dans les textes techniques, administratifs et journalistiques camerounais — conditionne la justesse terminologique de votre production. Cette section outille votre regard critique pour produire des résumés à la fois fluides, logiques et terminologiquement exacts.

\courssub{La cohésion : les liens de surface du texte}

\begin{definition}[Cohésion textuelle]
La \textbf{cohésion} désigne l'ensemble des procédés linguistiques explicites (marqueurs de surface) qui assurent la liaison \emph{formelle} entre les énoncés successifs d'un texte. Elle opère au niveau de la \emph{phrase} et de l'\emph{énoncé}, créant une chaîne ininterrompue qui guide le lecteur d'une idée à la suivante sans rupture syntaxique ni référentielle.
\end{definition}

La cohésion se manifeste par cinq grands leviers que tout résumeur doit savoir repérer, préserver ou reconstruire :

\begin{enumerate}
    \item \textbf{Les reprises anaphoriques (pronominales et nominales) :} Elles évitent la répétition lourde tout en maintenant la référence.
    \item \textbf{Les connecteurs (logiques, temporels, énumératifs) :} Ils signalent la relation sémantique entre propositions (cause, conséquence, opposition, concession, ajout, chronologie).
    \item \textbf{Les champs lexicaux et les chaînes de référence :} La répétition maîtrisée de mots-clés, l'usage de synonymes, d'hyperonymes ou de métonymies tissent une trame thématique continue.
    \item \textbf{La concordance des temps et des modes :} Elle ancre le discours dans une temporalité cohérente (récit au passé, explication au présent de vérité, argumentation au futur/conditionnel).
    \item \textbf{L'ellipse et la substitution :} L'omission calculée d'un élément récupérable dans le contexte immédiat (ex. \emph{« Le premier ministre a parlé ; le second [a parlé] aussi »}) allège la syntaxe sans perte de sens.
\end{enumerate}

\begin{popfigure}
\begin{tikzpicture}[
    node distance=1.2cm and 1.5cm,
    box/.style={draw=popblue, thick, rounded corners, fill=popblue!5, minimum width=2.8cm, minimum height=1.2cm, align=center, font=\small},
    arrow/.style={-Stealth, thick, popdark}
]
% Phrases
\node[box, align=center] (P1){Phrase 1 :\\ \textit{« Le Cameroun organise la CAN. »}};
\node[box, right=of P1, align=center] (P2){Phrase 2 :\\ \textit{« \textbf{Il} a construit des stades. »}};
\node[box, right=of P2, align=center] (P3){Phrase 3 :\\ \textit{« \textbf{Cet investissement} booste l'économie. »}};
\node[box, right=of P3, align=center] (P4){Phrase 4 :\\ \textit{« \textbf{Par conséquent}, l'emploi augmente. »}};
% Flèches anaphoriques (cohésion de surface)
\draw[arrow, bend left=20] (P1.east) to node[above, font=\tiny, popblue] {Pronom \textbf{Il} = Cameroun} (P2.west);
\draw[arrow, bend left=20] (P2.east) to node[above, font=\tiny, popblue] {Reprise nominale \textbf{Cet investissement}} (P3.west);
\draw[arrow, bend left=20] (P3.east) to node[above, font=\tiny, popblue] {Connecteur \textbf{Par conséquent}} (P4.west);
% Étiquette "Cohésion" au-dessus (remplace l'accolade)
\coordinate (midtop) at ($(P1.north)!0.5!(P4.north)$);
\node[above=1cm of midtop, poporange, font=\bfseries\large] {COHÉSION (Surface)};
\end{tikzpicture}
\legende{Schéma des liens cohésifs de surface : pronoms, reprises nominales et connecteurs assurent la chaîne formelle entre phrases successives.}
\end{popfigure}

\begin{methode}[Repérer la cohésion dans un texte source]
\begin{enumerate}
    \item \textbf{Surlignez} tous les pronoms personnels, démonstratifs, possessifs, relatifs et identifiez leur antécédent (flèche retour).
    \item \textbf{Encadrez} les connecteurs logiques (\emph{donc, cependant, en outre, d'abord, ensuite}) et temporels (\emph{puis, antérieurement, simultanément}).
    \item \textbf{Cerclez} les chaînes lexicales : mots répétés, synonymes, hyperonymes (ex. \emph{stade $\to$ enceinte $\to$ infrastructure}).
    \item \textbf{Vérifiez} la cohérence temporelle : les temps se succèdent-ils logiquement ? (Pas de saut brutal du passé simple au présent sans marqueur).
    \item \textbf{Notez} les ellipses : quel verbe ou nom est sous-entendu ?
\end{enumerate}
\end{methode}

\begin{exemplebox}[Analyse cohésive d'un extrait technique]
\textbf{Texte :} \emph{« La SONEL a lancé un vaste programme d'électrification rurale. \textbf{Ce projet} vise 500 villages. \textbf{Il} nécessite 200 milliards FCFA. \textbf{Toutefois}, le financement n'est pas encore bouclé. »}

\begin{itemize}
    \item \textbf{Reprise nominale :} \emph{Ce projet} = \emph{programme d'électrification rurale}.
    \item \textbf{Pronom personnel :} \emph{Il} = \emph{Ce projet} (sujet de \emph{nécessite}).
    \item \textbf{Connecteur adversatif :} \emph{Toutefois} signale une opposition (projet ambitieux vs financement incertain).
    \item \textbf{Chaîne lexicale :} \emph{programme / projet / financement / milliards FCFA} (champ sémantique de la planification budgétaire).
\end{itemize}
\end{exemplebox}

\courssub{La cohérence : l'unité de sens en profondeur}

\begin{definition}[Cohérence textuelle]
La \textbf{cohérence} est la propriété sémantique et pragmatique qui fait qu'un texte « tient ensemble » par le \emph{sens}. Elle repose sur l'organisation logique des idées (non-contradiction, hiérarchie thèse/arguments/exemples, progression thématique) et sur la \emph{pertinence} de l'enchaînement par rapport à l'intention communicative et au contexte encyclopédique du lecteur. Un texte cohérent construit une \cle{représentation mentale} unifiée et non contradictoire.
\end{definition}

Si la cohésion est le \emph{ciment} (visible), la cohérence est l'\emph{architecture} (invisible). Un texte peut être parfaitement cohésif (beaux connecteurs, belles reprises) mais incohérent (idées absurdes, contradictions, ordre illogique).

\begin{propriete}[Critères de cohérence]
Un texte (et donc un résumé) est cohérent s'il satisfait :
\begin{enumerate}
    \item \textbf{Non-contradiction :} Aucune proposition ne contredit une précédente sans marqueur de concession/rectification explicite.
    \item \textbf{Hiérarchisation :} Idée directrice (thèse) $\to$ Arguments $\to$ Exemples/Illustrations. Le résumé doit respecter cette pyramide.
    \item \textbf{Progression thématique :} Le thème (sujet du discours) et le rhème (information nouvelle) s'enchaînent selon les schémas vus en Section 1 (linéaire, à thème constant, hyperthématique).
    \item \textbf{Adéquation au genre :} Un résumé d'article scientifique suit la logique IMRaD (Introduction, Méthodes, Résultats, Discussion) ; un éditorial suit une logique argumentative.
\end{enumerate}
\end{propriete}

\begin{popfigure}
\begin{tikzpicture}[
    node distance=1.5cm and 2cm,
    level/.style={align=center, font=\small},
    l1/.style={level, fill=poppurple!10, draw=poppurple, thick, rounded corners, minimum width=4cm},
    l2/.style={level, fill=popblue!10, draw=popblue, thick, rounded corners, minimum width=3.5cm},
    l3/.style={level, fill=popgreen!10, draw=popgreen, thick, rounded corners, minimum width=3cm},
    arrow/.style={-Stealth, thick, popdark}
]
% Niveau 1 : Intention / Thèse
\node[l1, align=center] (N1){INTENTION COMMUNICATIVE\\ \textbf{Thèse centrale}\\(ex. « L'électrification rurale\\ est un levier de développement »)};
% Niveau 2 : Arguments principaux
\node[l2, below left=of N1, align=center] (N2a){ARGUMENT 1\\ \textit{Accès à l'énergie}\\\textit{ $\to$ activités génératrices de revenus}};
\node[l2, below right=of N1, align=center] (N2b){ARGUMENT 2\\ \textit{Amélioration services sociaux}\\\textit{ (santé, éducation)}};
% Niveau 3 : Exemples / Données
\node[l3, below left=of N2a, align=center] (N3a){EXEMPLE\\ \textit{Village de Batchenga :}\\\textit{ +40\% revenus ménage}};
\node[l3, below right=of N2a, align=center] (N3b){EXEMPLE\\ \textit{Centre de santé solaire :}\\\textit{ vaccins conservés}};
\node[l3, below left=of N2b, align=center] (N3c){EXEMPLE\\ \textit{École connectée :}\\\textit{ accès ressources numériques}};
\node[l3, below right=of N2b, align=center] (N3d){DONNÉE\\ \textit{Étude MINEDUB 2023 :}\\\textit{ taux réussite +15\%}};
% Flèches
\draw[arrow] (N1.south) -- ++(0,-0.3) -| (N2a.north);
\draw[arrow] (N1.south) -- ++(0,-0.3) -| (N2b.north);
\draw[arrow] (N2a.south) -- ++(0,-0.3) -| (N3a.north);
\draw[arrow] (N2a.south) -- ++(0,-0.3) -| (N3b.north);
\draw[arrow] (N2b.south) -- ++(0,-0.3) -| (N3c.north);
\draw[arrow] (N2b.south) -- ++(0,-0.3) -| (N3d.north);
% Étiquette "Cohérence" en dessous (remplace l'accolade)
\coordinate (midbot) at ($(N3a.south)!0.5!(N3d.south)$);
\node[below=1cm of midbot, poporange, font=\bfseries\large, align=center] {COHÉRENCE (Profondeur :\\ Logique thèse $\to$ arguments $\to$ preuves)};
\end{tikzpicture}
\legende{Architecture de la cohérence : organisation hiérarchique du sens (thèse, arguments, preuves). Le résumé doit préserver cette arborescence.}
\end{popfigure}

\begin{attention}[Piège majeur : Cohésion sans cohérence]
Un résumé peut être \textbf{cohésif en surface} (beaux connecteurs, pronoms bien résolus) mais \textbf{incohérent en profondeur} si l'on supprime une étape logique essentielle.
\textbf{Exemple :} Texte source : \emph{« Pour réduire la mortalité infantile, le MINSANTE vaccinate les enfants. \textbf{Or}, la chaîne du froid est défaillante. \textbf{Donc}, les vaccins perdent leur efficacité. \textbf{Aussi}, la mortalité ne baisse pas. »}
\emph{Mauvais résumé (cohésif mais incohérent) :} \emph{« Le MINSANTE vaccine les enfants. \textbf{Donc}, la mortalité ne baisse pas. »} $\to$ Le connecteur « donc » crée une cohésion de surface, mais la suppression de la cause (chaîne du froid) brise la cohérence causale. Le lecteur ne comprend \emph{pas pourquoi} la vaccination échoue.
\end{attention}

\courssub{Outils de la cohésion : boîte à outils du résumeur}

Pour garantir la cohésion de votre propre résumé, vous devez maîtriser activement trois mécanismes de substitution et de liaison.

\begin{definition}[Substitution pronominale]
Remplacement d'un syntagme nominal (ou proposition) par un pronom (personnel, démonstratif, relatif, indéfini) pour éviter la répétition tout en maintenant la référence.
\begin{itemize}
    \item \textbf{Personnel :} \emph{Le MINEDUB a publié les résultats. \textbf{Il} les a mis en ligne.}
    \item \textbf{Démonstratif :} \emph{Cet arrêté concerne les lycées techniques. \textbf{Celui-ci} entre en vigueur lundi.}
    \item \textbf{Relatif :} \emph{Le rapport \textbf{que} vous avez lu est confidentiel.}
\end{itemize}
\end{definition}

\begin{definition}[Reprise nominale (ou nomination variée)]
Remplacement par un synonyme, un hyperonyme, une périphrase ou un métonyme. Plus formelle que le pronom, elle relance l'attention et précise le référent.
\begin{itemize}
    \item \emph{Le chef de l'État $\to$ \textbf{le Président de la République} $\to$ \textbf{le chef de l'exécutif} $\to$ \textbf{Paul Biya}.}
    \item \emph{La CAN 2021 $\to$ \textbf{la compétition} $\to$ \textbf{le tournoi} $\to$ \textbf{cet événement continental}.}
\end{itemize}
\end{definition}

\begin{definition}[Ellipse et connecteurs]
\begin{itemize}
    \item \textbf{Ellipse :} Suppression d'un élément prévisible. \emph{« Yaoundé a 4 millions d'habitants ; Douala, \textbf{3,5 millions} [a]. »}
    \item \textbf{Connecteurs logiques :} \emph{donc, car, mais, cependant, par conséquent, en outre, en effet, c'est-à-dire}.
    \item \textbf{Connecteurs temporels/énumératifs :} \emph{d'abord, ensuite, puis, enfin ; premierement, deuxièmement}.
\end{itemize}
\end{definition}

\begin{methode}[Vérifier la cohésion de son résumé (auto-correction)]
\begin{enumerate}
    \item \textbf{Lisez votre résumé à voix haute.} L'oreille détecte les ruptures de rythme que l'œil manque.
    \item \textbf{Traquez chaque pronom :} A-t-il un antécédent clair \emph{dans le résumé} (pas seulement dans le texte source) ? Si l'antécédent a été coupé, remplacez le pronom par la reprise nominale.
    \item \textbf{Testez les connecteurs :} Remplacez mentalement « donc » par « parce que » ; « cependant » par « mais ». La relation logique tient-elle ?
    \item \textbf{Vérifiez les chaînes lexicales :} Le mot-clé du sujet (ex. \emph{électrification}) réapparaît-il assez souvent (ou via synonymes) pour ancrer le thème ?
    \item \textbf{Contrôlez la concordance des temps :} Résumé au présent de vérité (généralités) ou passé composé (faits datés) ? Pas de mélange anarchique.
\end{enumerate}
\end{methode}

\courssub{Le sigle : définition, formation, prononciation et règles d'usage}

\begin{definition}[Sigle et acronyme]
Un \textbf{sigle} est un mot formé par la concaténation des \textbf{initiales} (premières lettres) de plusieurs mots constituant une dénomination complexe (organisme, institution, concept technique).
\begin{itemize}
    \item \textbf{Acronyme (sigle prononcé comme un mot) :} Les lettres s'enchaînent phonétiquement. Ex. \textbf{SNI} (Stratégie Nationale de l'Information), \textbf{CNU} (Centre National Universitaire), \textbf{UNESCO}, \textbf{NASA}, \textbf{laser} (light amplification by stimulated emission of radiation), \textbf{ovni} (objet volant non identifié).
    \item \textbf{Sigle épelé (ou initialisme) :} Chaque lettre se prononce séparément. Ex. \textbf{CRTV} (C-a-R-T-V), \textbf{MINESEC} (M-I-N-E-S-E-C), \textbf{FECAFOOT}, \textbf{GICAM}, \textbf{ONU}, \textbf{USA}.
\end{itemize}
\end{definition}

\begin{propriete}[Règles orthographiques et grammaticales des sigles en français]
\begin{enumerate}
    \item \textbf{Majuscules intégrales :} \textbf{MINESEC}, \textbf{SONEL}, \textbf{CAMTEL}. (Pas de minuscules, sauf lexicalisation ancienne : \emph{laser, ovni, sida, radar}).
    \item \textbf{Pas de points abréviatifs :} On écrit \textbf{USA}, \textbf{ONU}, \textbf{CRTV} (et non U.S.A., O.N.U., C.R.T.V.).
    \item \textbf{Accord du pluriel :} \textbf{Facultatif} mais \textbf{recommandé} à l'écrit pour marquer le nombre quand le sigle n'est pas lexicalisé.
        \begin{itemize}
            \item \emph{Les \textbf{MINESEC} de la sous-région se réunissent.} (Pluriel marqué)
            \item \emph{Deux \textbf{CRTV} diffusent le match.} (Pluriel marqué)
            \item \emph{Des \textbf{laser} puissants.} (Lexicalisé : pluriel obligatoire en \emph{-s}).
        \end{itemize}
    \item \textbf{Genre :} Le sigle prend le genre du noyau de la dénomination complète.
        \begin{itemize}
            \item \textbf{La} SNI (la \emph{Stratégie}) — \textbf{Le} CNU (le \emph{Centre}) — \textbf{La} CRTV (la \emph{Radio} / la \emph{Télévision}).
        \end{itemize}
    \item \textbf{Article :} Souvent zéro à l'oral technique (\emph{« MINESEC a décidé »}), mais \emph{le/la} MINESEC à l'écrit formel.
\end{enumerate}
\end{propriete}

\begin{popfigure}
\begin{tikzpicture}[
    node distance=0.8cm and 1cm,
    acro/.style={draw=popblue, thick, fill=popblue!10, rounded corners, minimum width=3.5cm, minimum height=1cm, align=center, font=\small\bfseries},
    epele/.style={draw=poporange, thick, fill=poporange!10, rounded corners, minimum width=3.5cm, minimum height=1cm, align=center, font=\small\bfseries},
    def/.style={font=\footnotesize, align=left, text width=5cm},
    arrow/.style={-Stealth, popmuted}
]
% Titre colonnes
\node[acro] (AcroTitle) at (0,3) {ACRONYMES (prononcés en mot)};
\node[epele] (EpeleTitle) at (6,3) {SIGLES ÉPELÉS (lettre par lettre)};
% Acronymes
\node[acro, below=of AcroTitle] (SNI) {SNI};
\node[def, right=of SNI] (SNIdef) {Stratégie Nationale de l'Information};
\node[acro, below=of SNI] (CNU) {CNU};
\node[def, right=of CNU] (CNUdef) {Centre National Universitaire};
\node[acro, below=of CNU] (UNESCO) {UNESCO};
\node[def, right=of UNESCO] (UNESCOdef) {Organisation Nations Unies Éducation Science Culture};
\node[acro, below=of UNESCO] (LASER) {LASER $\to$ laser};
\node[def, right=of LASER] (LASERdef) {Lexicalisé : nom commun};
% Sigles épelés
\node[epele, below=of EpeleTitle] (MINESEC) {MINESEC};
\node[def, right=of MINESEC] (MINESECdef) {Ministère des Enseignements Secondaires};
\node[epele, below=of MINESEC] (MINEDUB) {MINEDUB};
\node[def, right=of MINEDUB] (MINEDUBdef) {Ministère de l'Éducation de Base};
\node[epele, below=of MINEDUB] (SONEL) {SONEL (hist.) / ENEO};
\node[def, right=of SONEL] (SONELdef) {Société Nationale d'Électricité / Energy of Cameroon};
\node[epele, below=of SONEL] (CAMTEL) {CAMTEL};
\node[def, right=of CAMTEL] (CAMTELdef) {Cameroon Telecommunications};
\node[epele, below=of CAMTEL] (FECAFOOT) {FECAFOOT};
\node[def, right=of FECAFOOT] (FECAFOOTdef) {Fédération Camerounaise de Football};
\node[epele, below=of FECAFOOT] (GICAM) {GICAM};
\node[def, right=of GICAM] (GICAMdef) {Groupement Inter-patronal du Cameroun};
% Flèches de classification
\draw[arrow] (AcroTitle.south) -- (SNI.north);
\draw[arrow] (EpeleTitle.south) -- (MINESEC.north);
\end{tikzpicture}
\legende{Nuage de sigles camerounais classés : acronymes (prononciation syllabique) vs sigles épelés (lettre par lettre). Maîtrisez leur graphie, genre et pluriel.}
\end{popfigure}

\begin{savaistu}[« Laser » et « Ovni » : de sigles à noms communs]
\textbf{Laser} (1960) = \textbf{L}ight \textbf{A}mplification by \textbf{S}timulated \textbf{E}mission of \textbf{R}adiation. \textbf{Ovni} (années 50) = \textbf{O}bjet \textbf{V}olant \textbf{N}on \textbf{I}dentifié. Leur usage massif a gommé la majuscule et imposé le pluriel en \emph{-s} (\emph{des lasers, des ovnis}). C'est le destin de tout sigle très fréquent : il \textbf{se lexicalise}. Au Cameroun, \textbf{sida} (Syndrome d'Immuno-Déficit Acquis) a suivi ce chemin : s'écrit en minuscule, invariable au pluriel (\emph{la lutte contre le sida}). En rédaction technique, conservez les majuscules pour les sigles institutionnels (MINESEC, FECAFOOT) qui ne sont pas lexicalisés.
\end{savaistu}

\courssub{Sigles courants de la vie camerounaise : usage correct dans une copie}

Le Probatoire / Baccalauréat technique attend une orthographe irréprochable des sigles institutionnels. Voici le référentiel minimal à connaître par cœur.

\begin{center}
\begin{tabularx}{\linewidth}{|l|X|l|X|}\hline
\rowcolor{popblueL}
\textbf{Sigle} & \textbf{Dénomination complète} & \textbf{Genre} & \textbf{Contexte d'usage typique} \\ \hline
\textbf{MINESEC} & Ministère des Enseignements Secondaires & Masc. & Réformes curriculaires, examens, notes de service. \\ \hline
\textbf{MINEDUB} & Ministère de l'Éducation de Base & Masc. & École primaire, alphabétisation, cantines scolaires. \\ \hline
\textbf{MINSANTE} & Ministère de la Santé Publique & Masc. & Campagnes vaccination, hôpitaux, paludisme, COVID. \\ \hline
\textbf{MINEPAT} & Ministère de l'Économie, de la Planification et de l'Aménagement du Territoire & Masc. & Budget, PND, indicateurs macroéconomiques. \\ \hline
\textbf{ENEO} & Energy of Cameroon (ex-SONEL) & Masc. & Électricité, factures, coupures, barrages. \\ \hline
\textbf{CAMTEL} & Cameroon Telecommunications & Masc. & Téléphone fixe, internet, fibre optique. \\ \hline
\textbf{FECAFOOT} & Fédération Camerounaise de Football & Fém. & Lions Indomptables, CAN, championnats locaux. \\ \hline
\textbf{GICAM} & Groupement Inter-patronal du Cameroun & Masc. & Patronat, dialogue social, climat des affaires. \\ \hline
\textbf{CNPS} & Caisse Nationale de Prévoyance Sociale & Fém. & Retraites, accidents travail, cotisations. \\ \hline
\textbf{ARMP} & Agence de Régulation des Marchés Publics & Fém. & Appels d'offres, transparence, contrats publics. \\ \hline
\end{tabularx}
\end{center}

\begin{exoresolu}[Correction d'un paragraphe truffé d'erreurs sur les sigles]
\textbf{Énoncé :} Recopiez et corrigez le paragraphe suivant : \emph{« Le mineSEC a signé une convention avec la CNPS pour la couverture sociale des enseignants. La FECAFOOT a aussi demandé l'appui du GICAM pour organiser la CAN. Le CAMTEL va fournir la connexion internet. »}

\textbf{Solution détaillée :}
\begin{enumerate}
    \item \textbf{mineSEC $\to$ MINESEC} (Majuscules obligatoires).
    \item \textbf{la CNPS} : Correct (genre féminin : \emph{la Caisse}).
    \item \textbf{La FECAFOOT} : Correct (genre féminin : \emph{la Fédération}).
    \item \textbf{le GICAM} : Préférer \textbf{le GICAM} (masculin : \emph{le Groupement}) ou \textbf{GICAM} sans article (style télégraphique admis en résumé).
    \item \textbf{Le CAMTEL} : Correct (masculin : \emph{le Cameroon Telecommunications} / \emph{l'opérateur}).
    \item \textbf{Ponctuation :} Pas de points entre les lettres.
\end{enumerate}
\textbf{Version corrigée :} \emph{« Le MINESEC a signé une convention avec la CNPS pour la couverture sociale des enseignants. La FECAFOOT a aussi demandé l'appui du GICAM pour organiser la CAN. CAMTEL (ou Le CAMTEL) va fournir la connexion internet. »}
\end{exoresolu}

\courssub{Pourquoi cohésion et cohérence sont les critères majeurs de réussite d'un résumé}

\begin{aretenir}[Critères d'évaluation du résumé (Grille type Bac/Probatoire)]
\begin{itemize}
    \item \textbf{Fidélité au sens (Cohérence) :} 4 à 6 pts. Le résumé rend-il la thèse, les arguments principaux et leur articulation logique sans contresens ni omission d'étape cruciale ?
    \item \textbf{Qualité de l'expression / Cohésion :} 3 à 5 pts. Le texte produit est-il fluide ? Les connecteurs sont-ils pertinents ? Les reprises (pronominales/nominales) assurent-elles la continuité référentielle sans ambiguïté ? La concordance des temps est-elle respectée ?
    \item \textbf{Concision et reformulation :} 3 à 4 pts. Le taux de compression est-il respecté (1/3 à 1/4) ? Le vocabulaire est-il propre (pas de copier-coller) ?
    \item \textbf{Orthographe / Sigles :} 1 à 2 pts. Zéro faute sur les sigles institutionnels (MINESEC, FECAFOOT, etc.), accords, majuscules.
\end{itemize}
\end{aretenir}

\begin{exemplebox}[Mini-résumé incohérent $\to$ Corrigé (CAN 2021)]
\textbf{Texte source (résumé) :} \emph{« Le Cameroun a organisé la CAN 2021. Des stades ont été construits. L'économie a été boostée. Il y a eu des critiques sur les coûts. »}

\textbf{Version élève (incohérente) :} \emph{« La CAN 2021 a eu lieu au Cameroun. \textbf{Par conséquent}, des stades ont été construits. \textbf{Mais}, l'économie a été boostée. \textbf{Or}, il y a eu des critiques. »}
\begin{itemize}
    \item \textbf{Erreurs cohérence :} « Par conséquent » (conséquence) entre organisation et construction ? Non, la construction est \emph{antérieure} (condition). « Mais » (opposition) entre construction et boost économique ? Non, relation de cause à effet. « Or » (renforcement/opposition) mal placé.
\end{itemize}

\textbf{Version corrigée (cohésive \& cohérente) :} \emph{« Pour organiser la CAN 2021, le Cameroun a construit des stades. \textbf{Cet investissement} a dynamisé l'économie. \textbf{Toutefois}, les coûts ont suscité des critiques. »}
\begin{itemize}
    \item \textbf{Outils cohésion :} \emph{Pour} (but), \emph{Cet investissement} (reprise nominale), \emph{Toutefois} (concession/adversatif).
    \item \textbf{Logique cohérence :} But $\to$ Moyens $\to$ Effet positif $\to$ Point de vue critique (équilibre).
\end{itemize}
\end{exemplebox}

\begin{exercice}[Entraînement : Cohésion, Cohérence, Sigles]
\begin{enumerate}
    \item \textbf{Analyse cohésive :} Dans l'extrait ci-dessous, soulignez les reprises (pronominales/nominales), encadrez les connecteurs, identifiez la chaîne lexicale principale.
    \emph{« Le MINEDUB a lancé le programme « École propre » dans tous les établissements publics du pays. Ce projet ambitieux vise l'assainissement des cours et des salles de classe par une mobilisation hebdomadaire. Il mobilise activement les élèves, les enseignants et les parents d'élèves autour d'un objectif commun. En effet, l'hygiène rigoureuse prévient efficacement les maladies hydriques et parasitaires qui freinent la scolarité. Cependant, le manque d'eau potable dans de nombreuses localités rurales complique considérablement la mise en œuvre quotidienne de ce programme salutaire. »}
    \item \textbf{Correction de cohérence :} Le résumé suivant brise la logique causale. Réécrivez-le en restaurant la cohérence (choisissez les bons connecteurs : \emph{car / donc / mais / pour / bien que}).
    \emph{« La CNPS augmente les pensions. \_\_\_\_\_\_ l'inflation érode le pouvoir d'achat. \_\_\_\_\_\_ les retraités restent pauvres. \_\_\_\_\_\_ cette mesure, le gouvernement veut les soulager. »}
    \item \textbf{Orthographe sigles :} Écrivez correctement : \emph{le minesec ; la fecafoot ; des camtel ; le gicam ; une unesco ; des sida.}
    \item \textbf{Production :} Résumez en 40 mots ($\pm$10\%) le paragraphe de l'exercice 1 en respectant cohésion, cohérence et sigles.
\end{enumerate}
\end{exercice}

\begin{corrige}[Exercice n°1 à 4]
\begin{enumerate}
    \item \textbf{Reprises :} \emph{Ce projet ambitieux} (nominal, = programme) ; \emph{Il} (pronominal, = Ce projet). \textbf{Connecteurs :} \emph{En effet} (explication/justification) ; \emph{Cependant} (concession/opposition). \textbf{Chaîne lexicale :} \emph{projet / assainissement / hygiène / maladies / eau / programme} (champ santé/environnement/éducation).
    \item \textbf{Version cohérente :} \emph{« La CNPS augmente les pensions \textbf{car} l'inflation érode le pouvoir d'achat. \textbf{Bien que} cette mesure vise à les soulager, les retraités restent pauvres. »} (Logique : Cause $\to$ Mesure $\to$ Concession sur l'insuffisance).
    \item \textbf{MINESEC ; la FECAFOOT ; des CAMTEL ; le GICAM ; l'UNESCO (acronyme) ; le sida (invariable, minuscule).}
    \item \textbf{Proposition de résumé (env. 39 mots) :} \emph{« Le MINEDUB lance « École propre » pour assainir les établissements. Ce projet mobilise la communauté éducative car l'hygiène prévient les maladies. Cependant, le manque d'eau en zone rurale en entrave l'application effective. »}
\end{enumerate}
\end{corrige}

\begin{methode}[Synthèse : Check-list finale avant de rendre votre résumé]
\begin{enumerate}
    \item \textbf{Sens (Cohérence) :} Ai-je gardé la thèse + les 2-3 arguments clés + l'exemple crucial ? L'ordre est-il logique ? Pas de contresens ?
    \item \textbf{Liens (Cohésion) :} Chaque pronom a-t-il son antécédent \emph{dans mon texte} ? Les connecteurs reflètent-ils la \emph{vraie} relation logique (cause, conséq., oppos., concess.) ? Les temps sont-ils stables ?
    \item \textbf{Sigles :} Tous en MAJUSCULES, sans points, genre correct (le/la), pluriel -s si nécessaire (Les MINESEC).
    \item \textbf{Longueur :} Compte des mots $\approx$ 25-33\% du texte source.
    \item \textbf{Reformulation :} Pas de phrases copiées-collées. Vocabulaire personnel.
\end{enumerate}
\end{methode}

\coursec{Le Vieux nègre et la médaille : contexte, tonalités et bilan de l'œuvre}

Dans la section précédente, nous avons vu comment la \cle{cohésion} (les liens grammaticaux et lexicaux entre les phrases) et la \cle{cohérence} (l'unité de sens du texte) permettent à un lecteur de suivre un propos sans se perdre. Ces notions vont maintenant nous servir d'outils de lecture : pour comprendre un roman, il ne suffit pas de suivre l'histoire phrase après phrase, il faut aussi saisir ce que l'œuvre \emph{veut dire} dans son époque, et \emph{comment} elle le dit. C'est précisément l'objet de cette section : situer \emph{Le Vieux nègre et la médaille} de Ferdinand Oyono dans son contexte, analyser ses tonalités satirique et polémique, puis dresser le bilan de notre étude de l'œuvre.

\courssub{Inscription de l'œuvre dans son contexte}

\textbf{Le Cameroun sous administration française.} Après la défaite de l'Allemagne en 1916, le Cameroun, ancienne colonie allemande, est partagé entre la France et la Grande-Bretagne. La partie française est administrée d'abord sous mandat de la Société des Nations (1922), puis sous tutelle de l'ONU (1946). Concrètement, cela signifie que le pays est dirigé par des administrateurs coloniaux français : commandants de cercle, chefs de subdivision, gendarmes, missionnaires. Les populations africaines sont soumises à des obligations lourdes : impôts, travaux forcés, corvées, portage. L'administration prétend apporter la « civilisation », mais elle repose sur l'inégalité et l'humiliation quotidienne.

\textbf{La littérature anticoloniale des années 1950.} Dans les années qui précèdent les indépendances africaines, une génération d'écrivains africains de langue française prend la plume pour dénoncer la colonisation : Mongo Beti (\emph{Le Pauvre Christ de Bomba}, 1956), Camara Laye, Sembène Ousmane, et bien sûr Ferdinand Oyono avec \emph{Une vie de boy} (1956) et \emph{Le Vieux nègre et la médaille} (1956). Ces romans démontent le mythe de la « mission civilisatrice » et montrent la réalité vécue par les colonisés. Ils préparent les esprits à l'indépendance, obtenue par le Cameroun français le 1\textsuperscript{er} janvier 1960.

\begin{savaistu}[Un écrivain devenu homme d'État]
Ferdinand Oyono n'a que \textbf{27 ans} lorsqu'il publie \emph{Le Vieux nègre et la médaille} en 1956. Né en 1929 à Ngoulemakong, dans la région du Sud du Cameroun, il écrit ce roman alors qu'il est encore étudiant en France. Après l'indépendance, il mène une brillante carrière diplomatique et politique : ambassadeur du Cameroun dans plusieurs pays, représentant permanent à l'ONU, puis ministre. Son œuvre romanesque, brève (trois romans), reste pourtant l'un des sommets de la littérature anticoloniale africaine.
\end{savaistu}

\begin{popfigure}
\begin{center}
\begin{tikzpicture}[x=1.15cm]
% Axe
\draw[line width=1.2pt,popdark,->] (0,0) -- (10.6,0);
% Graduations
\foreach \x/\y in {0/1916, 2.2/1922, 4.4/1946, 7.7/1956, 9.9/1960}{
  \draw[popdark] (\x,0.08) -- (\x,-0.08);
  \node[below,font=\small\bfseries,popdark] at (\x,-0.12) {\y};
}
% Bandeau administration française
\fill[popblue!25] (0,0.25) rectangle (9.9,0.75);
\node[font=\small\itshape,popblue] at (4.95,0.5) {Administration française du Cameroun};
% Événements
\fill[poporange] (2.2,0) circle (2.5pt);
\node[above,font=\scriptsize,poporange,align=center] at (2.2,0.85) {Mandat de la\\Société des Nations};
\fill[popgreen] (4.4,0) circle (2.5pt);
\node[above,font=\scriptsize,popgreen!60!black,align=center] at (4.4,0.85) {Tutelle de l'ONU\\(après la guerre)};
\fill[popink] (7.7,0) circle (3pt);
\node[above,font=\scriptsize\bfseries,popink,align=center] at (7.7,0.85) {Publication du roman\\de F. Oyono};
\fill[popgold] (9.9,0) circle (3pt);
\node[above,font=\scriptsize\bfseries,popgold!70!black,align=center] at (9.9,0.85) {Indépendance\\du Cameroun};
\end{tikzpicture}
\end{center}
\legende{Frise chronologique : le Cameroun sous administration française (1916--1960). Le roman paraît en 1956, quatre ans avant l'indépendance.}
\end{popfigure}

\begin{aretenir}[Contexte : ce qu'il faut retenir]
\emph{Le Vieux nègre et la médaille} (1956) s'inscrit dans la \cle{littérature anticoloniale} des années 1950 : il dénonce, par la fiction, les mensonges et les humiliations de la colonisation française au Cameroun, à la veille des indépendances.
\end{aretenir}

\courssub{Résumé de l'intrigue : Meka, la médaille, l'humiliation, la prise de conscience}

Le héros, \textbf{Meka}, est un vieillard du village de Doum. Il a beaucoup donné à la France : ses terres ont été cédées à la mission catholique, ses deux fils sont morts pour la France pendant la Seconde Guerre mondiale, et il s'est lui-même laissé convaincre par le catéchisme. Pour le « récompenser », l'administration coloniale décide de lui remettre une \cle{médaille} de l'amitié franco-africaine lors de la fête du 14 juillet.

La cérémonie est un triomphe apparent : Meka, en boubou neuf, reçoit sa médaille devant tout le village, serre la main du commandant et du père Vandermayer. Mais très vite, la réalité reprend ses droits. Invité chez le commandant, il est traité avec condescendance. La nuit venue, perdu dans le quartier européen sous une pluie torrentielle, il est arrêté par des gendarmes qui ne reconnaissent pas le « héros » du matin : ils le bousculent, l'humilient et le jettent en prison comme un vulgaire délinquant. La médaille sur sa poitrine ne le protège de rien.

Cette nuit d'\cle{humiliation} provoque chez Meka une \cle{prise de conscience} : il comprend que la médaille n'est qu'un leurre, que l'« amitié » des Blancs est un masque, et que sa docilité ne lui a rien rapporté. Revenu au village, il refuse désormais de se soumettre et retrouve une forme de dignité dans la révolte intérieure.

\begin{attention}[Piège de résumé]
Ne racontez pas l'intrigue comme une simple suite d'aventures. L'essentiel est l'\textbf{arc du personnage} : illusion (la médaille) $\rightarrow$ épreuve (l'humiliation) $\rightarrow$ lucidité (la prise de conscience). C'est cette progression qui donne son sens au roman.
\end{attention}

\courssub{La tonalité satirique : railler les travers et les institutions}

\begin{definition}[Tonalité satirique]
La \cle{tonalité satirique} est le registre d'un texte qui \textbf{dénonce les travers} d'une personne, d'une institution ou d'une société \textbf{en les tournant en ridicule}. La satire fait rire, mais c'est un rire qui corrige : elle vise à réveiller la lucidité du lecteur. Dans le roman, Oyono raille l'administration coloniale (le commandant, les gendarmes), la mission religieuse (le père Vandermayer) et même la naïveté de certains villageois.
\end{definition}

La cérémonie du 14 juillet est un morceau de satire achevé. Tout y est grotesque : les discours pompeux sur l'« amitié franco-africaine », les notables africains qui se pavanent, le commandant qui décore Meka comme on récompense un animal fidèle. Le narrateur souligne l'écart entre la solennité affichée et la réalité mesquine : une médaille de pacotille pour une vie de sacrifices. Le rire du lecteur naît de cet \textbf{écart}, et ce rire juge.

\courssub{La tonalité polémique : contester violemment une idée ou un pouvoir}

\begin{definition}[Tonalité polémique]
La \cle{tonalité polémique} est le registre d'un texte qui \textbf{attaque frontalement} une idée, une doctrine ou un pouvoir pour les réfuter. Contrairement à la satire, elle ne cherche pas d'abord à faire rire : elle conteste, accuse, révolte. Dans le roman, la dimension polémique apparaît dans la dénonciation directe de l'injustice coloniale : la médaille est démasquée comme un instrument de domination, et la « mission civilisatrice » comme un mensonge.
\end{definition}

\begin{attention}[Erreur fréquente : satire ou polémique ?]
Ne confondez pas les deux tonalités. La \textbf{satire} passe par le \emph{rire} : elle ridiculise pour corriger (registre comique, ironique). La \textbf{polémique} est une \emph{attaque frontale} : elle réfute et indigne (registre oratoire, accusateur). Dans le roman d'Oyono, les deux se mêlent : la cérémonie du 14 juillet est plutôt \emph{satirique} (on rit du grotesque), tandis que la nuit de la pluie et la réflexion finale de Meka sont plutôt \emph{polémiques} (on s'indigne de l'injustice). À l'examen, citez toujours la tonalité \emph{et} le passage qui la justifie.
\end{attention}

\textbf{Les procédés communs.} Satire et polémique utilisent des procédés d'amplification et de déformation :

\begin{center}
\begin{tabularx}{\linewidth}{|l|X|X|}
\hline
\rowcolor{popblueL} \textbf{Procédé} & \textbf{Exemple dans le roman} & \textbf{Effet produit} \\
\hline
\cle{Ironie} & Les éloges de l'« amitié franco-africaine » pendant la cérémonie, alors que Meka sera jeté en prison le soir même. & Dire le contraire de ce qu'on pense : le lecteur mesure l'hypocrisie du discours colonial. \\
\hline
\cle{Caricature} & Le commandant, vaniteux et suffisant ; le père Vandermayer, obsédé par les statistiques de conversions. & Exagération des défauts : les représentants de l'ordre colonial deviennent ridicules. \\
\hline
\cle{Antithèse} & Le triomphe du matin (médaille, discours, fête) opposé à la honte de la nuit (pluie, gendarmes, prison). & Opposition brutale : elle révèle le mensonge de la « récompense ». \\
\hline
\cle{Hyperbole} & La pluie diluvienne, l'errance interminable de Meka dans le quartier européen. & Exagération : elle dramatise l'épreuve et souligne l'absurdité de la situation. \\
\hline
\cle{Grotesque} & Meka décoré comme un « bon nègre », exhibé puis abandonné. & Mélange de rire et de malaise : le lecteur rit, puis s'indigne. \\
\hline
\end{tabularx}
\end{center}

\begin{methode}[Identifier une tonalité en trois questions]
Pour déterminer la tonalité d'un passage, procédez toujours ainsi :
\begin{enumerate}
\item \textbf{Repérer la cible} : qui ou quoi le texte attaque-t-il ? (une personne, une institution, une idée)
\item \textbf{Repérer le procédé} : comment l'attaque-t-il ? (ironie, caricature, antithèse, hyperbole, invective...)
\item \textbf{Repérer l'effet recherché} : que veut provoquer l'auteur ? (le rire correcteur $\rightarrow$ satire ; l'indignation, la révolte $\rightarrow$ polémique)
\end{enumerate}
Rédigez ensuite votre réponse en une phrase complète : « Ce passage est de tonalité \ldots car l'auteur s'en prend à \ldots par le procédé de \ldots, afin de \ldots ».
\end{methode}

\courssub{Analyse guidée de deux passages clés}

\textbf{La cérémonie du 14 juillet.} Ce passage est le sommet satirique du roman. La cible est double : l'administration (qui instrumentalise Meka pour sa propagande) et les notables africains (qui se laissent duper). Le procédé dominant est l'\cle{ironie} : plus le discours officiel est grandiloquent, plus la récompense paraît dérisoire. La \cle{caricature} des personnages (le commandant gonflé d'importance, le père Vandermayer comptant ses « brebis ») achève de rendre l'ordre colonial ridicule. L'effet recherché : un rire lucide, qui démasque la supercherie.

\begin{exemplebox}[Commentaire guidé : Meka sous la pluie, médaille sur la poitrine]
\textbf{Situation.} La nuit du 14 juillet, Meka, ivre et égaré dans le quartier réservé aux Européens, est surpris par une pluie violente. Il porte encore sa médaille. Des gendarmes l'arrêtent : la médaille ne lui sert de rien ; il est frappé, insulté, emprisonné.

\textbf{Analyse.}
\begin{itemize}
\item \textbf{Cible} : le système colonial tout entier, qui honore l'Africain en public et le brutalise dès que la fête est finie.
\item \textbf{Procédés} : l'\emph{antithèse} (gloire du jour / honte de la nuit), l'\emph{hyperbole} (la pluie diluvienne, l'errance sans fin), le \emph{symbole} (la médaille, décor inutile, devient le signe de la trahison).
\item \textbf{Tonalité} : dominante \emph{polémique} — le rire s'efface devant l'indignation. La scène accuse : voilà ce que vaut l'« amitié » coloniale.
\item \textbf{Effet} : le lecteur éprouve pitié et colère ; il partage la prise de conscience de Meka.
\end{itemize}
\textbf{Phrase de synthèse.} « Ce passage, de tonalité polémique, dénonce l'hypocrisie coloniale en opposant, par antithèse, la cérémonie triomphale à l'arrestation brutale de Meka : la médaille, symbole vidé de sens, révèle que la reconnaissance des Blancs n'était qu'un leurre. »
\end{exemplebox}

\courssub{Évaluation et bilan de l'étude de l'œuvre}

\begin{exoresolu}[Questionnaire de synthèse (extrait corrigé)]
\textbf{Question 1.} Quel est le thème central du roman ?

\textbf{Question 2.} Quelle évolution le personnage de Meka connaît-il ?

\textbf{Question 3.} Quelle est la portée de l'œuvre aujourd'hui ?
\tcblower
\textbf{Réponse 1.} Le thème central est la \emph{dénonciation de la colonisation} : à travers le sort de Meka, Oyono montre que la « mission civilisatrice » et l'« amitié franco-africaine » ne sont que des masques de la domination et de l'humiliation.

\textbf{Réponse 2.} Meka passe de la \emph{soumission naïve} (il croit aux promesses des Blancs et accepte la médaille) à la \emph{lucidité révoltée} : l'humiliation de la nuit du 14 juillet lui ouvre les yeux, et il retrouve sa dignité en refusant désormais de se soumettre.

\textbf{Réponse 3.} Le roman invite à rester vigilant face aux \emph{fausses récompenses} et aux discours flatteurs du pouvoir. Sa résonance avec le Cameroun d'aujourd'hui est forte : il rappelle le prix de la dignité, la nécessité de la mémoire historique et l'importance de juger les actes plutôt que les paroles.
\end{exoresolu}

\begin{exercice}[À vous de jouer]
\begin{enumerate}
\item Relevez dans la cérémonie du 14 juillet deux procédés satiriques et expliquez leur effet.
\item La nuit de la pluie est-elle un passage satirique ou polémique ? Justifiez en trois phrases avec la méthode cible--procédé--effet.
\item Rédigez un court paragraphe (8 à 10 lignes) intitulé « Ce que le roman révèle de la colonisation ».
\end{enumerate}
\end{exercice}

\begin{corrige}[Exercice : indications]
\begin{enumerate}
\item On attend par exemple l'ironie des discours officiels (louer l'« amitié » tout en méprisant les Africains) et la caricature du commandant ou du père Vandermayer ; effet : ridiculiser l'ordre colonial et éveiller la lucidité du lecteur.
\item Passage à dominante \emph{polémique} : la cible est l'injustice coloniale, les procédés sont l'antithèse et l'hyperbole, l'effet recherché est l'indignation, non le rire.
\item Le paragraphe doit articuler : l'hypocrisie du discours colonial, la réalité des humiliations (travail, mépris, violence), et la prise de conscience comme première étape de la libération.
\end{enumerate}
\end{corrige}

\begin{aretenir}[Bilan de l'étude de l'œuvre]
\emph{Le Vieux nègre et la médaille} révèle que la colonisation reposait sur un \textbf{mensonge} : celui de la fraternité et de la civilisation, démenti par l'humiliation quotidienne des colonisés. Par la satire (le rire qui corrige) et la polémique (l'attaque qui révolte), Oyono transforme un fait divers local — un vieillard décoré puis jeté en prison — en \textbf{acte d'accusation universel}. L'œuvre garde aujourd'hui toute sa force : elle nous apprend à démasquer les discours flatteurs du pouvoir et à défendre la dignité de chaque homme.
\end{aretenir}

Cette étude de l'œuvre achevée, nous disposons désormais d'un texte riche de sens et d'arguments : c'est le point de départ idéal pour aborder, dans la prochaine section, la \cle{structure argumentative du texte}, étape indispensable avant d'apprendre à contracter et à résumer un texte.

\coursec{La structure argumentative du texte}

Après avoir exploré les tonalités satirique et polémique dans \emph{Le Vieux nègre et la médaille} et dressé le bilan de l'étude de l'œuvre, nous abordons maintenant une compétence transversale indispensable : l'analyse de la structure argumentative. Que l'on doive résumer un éditorial, une dissertation ou un discours, c'est le « squelette » logique du texte qu'il faut préserver. Comprendre comment s'articulent thèse, arguments et exemples permet de produire un résumé fidèle, cohérent et bien proportionné.

\courssub{Qu'est-ce que l'argumentation ?}

\begin{definition}[Argumentation]
L'\cle{argumentation} est un discours qui vise à faire adhérer un destinataire à une \cle{thèse} (une opinion, une affirmation) en mobilisant des \cle{arguments} étayés d'\cle{exemples} ou de faits. Elle repose sur une intention persuasive ou démonstrative.
\end{definition}

Contrairement à l'explication (qui montre \emph{comment} ou \emph{pourquoi} un phénomène se produit) ou à la narration (qui raconte une suite d'événements), l'argumentation répond à la question : \textbf{« Que veut me faire croire ou faire l'auteur, et pourquoi ? »}.

\begin{exemplebox}[Distinction rapide]
\begin{itemize}
    \item \textbf{Explication :} « La route Yaoundé–Douala se dégrade à cause de l'érosion et du trafic lourd. » (Cause $\to$ Effet)
    \item \textbf{Argumentation :} « Il faut réhabiliter d'urgence la route Yaoundé–Douala \emph{car} elle est vitale pour l'économie et son état actuel cause des accidents mortels. » (Thèse + Arguments)
\end{itemize}
\end{exemplebox}

\courssub{Les composantes de la structure argumentative}

Toute argumentation construite repose sur une architecture hiérarchique à trois étages.

\begin{definition}[Thèse, Argument, Exemple]
\begin{itemize}
    \item \textbf{Thèse :} Proposition centrale que l'auteur soutient (ou combat s'il s'agit d'une antithèse). Elle est souvent explicite (dans l'introduction ou la conclusion) mais peut être implicite.
    \item \textbf{Argument :} Raison logique, fait, valeur ou autorité invoquée pour appuyer la thèse. Un texte comporte généralement 2 à 4 arguments principaux.
    \item \textbf{Exemple (ou illustration) :} Cas concret, statistique, citation, anecdote, référence historique qui \emph{matérialise} l'argument et le rend tangible.
\end{itemize}
\end{definition}

\begin{aretenir}[Hiérarchie à respecter dans le résumé]
Dans un résumé, la thèse et les arguments sont \textbf{indispensables} (ils forment le squelette). Les exemples sont \textbf{secondaires} : on ne garde que les plus frappants, et souvent on les fusionne ou on les résume en une clause. \textbf{Erreur fatale :} supprimer un argument pour garder un exemple.
\end{aretenir}

\begin{popfigure}
\begin{tikzpicture}[
    level distance=2.2cm,
    sibling distance=3.5cm,
    every node/.style={draw, rounded corners, fill=popcream, thick, inner sep=4pt, font=\small, align=center},
    edge from parent/.style={draw=popdark, thick, ->},
    level 1/.style={sibling distance=5cm},
    level 2/.style={sibling distance=2.5cm}
]
\node [fill=popblue!20, font=\bfseries, align=center]{Thèse\\ (Idée directrice)}
    child { node [fill=popgreen!20] {Argument 1}
        child { node [fill=poporange!20, font=\itshape] {Exemple 1.1} }
        child { node [fill=poporange!20, font=\itshape] {Exemple 1.2} }
    }
    child { node [fill=popgreen!20] {Argument 2}
        child { node [fill=poporange!20, font=\itshape] {Exemple 2.1} }
    }
    child { node [fill=popgreen!20] {Argument 3}
        child { node [fill=poporange!20, font=\itshape] {Exemple 3.1} }
        child { node [fill=poporange!20, font=\itshape] {Exemple 3.2} }
    };
\node[below=1.5cm, font=\small\bfseries, text=popdark, align=center] {Arbre de la structure argumentative : la thèse domine les arguments, qui portent les exemples.};
\end{tikzpicture}
\legende{Arbre TikZ de la structure argumentative : thèse au sommet, arguments au niveau intermédiaire, exemples en feuilles.}
\end{popfigure}

\courssub{Les principaux types d'arguments}

Connaître la nature d'un argument aide à évaluer sa force et à le reformuler sans le trahir.

\begin{propriete}[Typologie des arguments (recensement Bac)]
\begin{itemize}
    \item \textbf{Argument d'autorité :} S'appuie sur un expert, un texte de loi, un proverbe, une personnalité reconnue. \emph{Ex. : « Selon l'OMS, la vaccination réduit la mortalité de 70\,\%. »}
    \item \textbf{Argument d'expérience (ou de fait) :} Repose sur une observation vérifiable, une statistique, un témoignage direct. \emph{Ex. : « En 2023, 12 accidents mortels ont été recensés sur ce tronçon. »}
    \item \textbf{Argument de valeur :} Mobilise des valeurs partagées (justice, solidarité, patriotisme, éthique). \emph{Ex. : « Laisser nos routes se dégrader, c'est trahir l'idéal d'émergence. »}
    \item \textbf{Argument logique (adéquation) :} Déduit une conséquence nécessaire d'un principe admis. \emph{Ex. : « Si l'État investit dans l'entretien, les coûts de réparation future baisseront. »}
    \item \textbf{Argument pragmatique (d'efficacité) :} Montre l'utilité, le gain ou le risque évité. \emph{Ex. : « Réhabiliter cette route fera gagner 2h par trajet aux transporteurs. »}
\end{itemize}
\end{propriete}

\begin{attention}[Piège : confondre type d'argument et connecteur]
Un connecteur (\emph{car, donc, cependant}) signale une \emph{relation logique} ; le \emph{type d'argument} désigne la \emph{nature du contenu} (autorité, fait, valeur...). Au Bac, on peut vous demander : « Quel type d'argument est utilisé ligne 5 ? » $\to$ Réponse attendue : « Argument d'expérience (statistique) », et non « Argument de cause ».
\end{attention}

\courssub{Marqueurs de la structure et connecteurs logiques}

La progression argumentative est balisée par des connecteurs qui révèlent l'enchaînement des idées. Les repérer permet de segmenter le texte en unités argumentatives.

\begin{center}
\begin{tabularx}{\linewidth}{|l|X|X|}
\hline
\rowcolor{popblueL}
\textbf{Relation logique} & \textbf{Connecteurs fréquents} & \textbf{Effet dans l'argumentation} \\ \hline
Addition / Enchaînement & \emph{de plus, par ailleurs, en outre, également, d'une part... d'autre part} & Ajoute un argument nouveau, renforce la thèse \\ \hline
Opposition / Concession & \emph{mais, cependant, pourtant, certes... cependant, bien que} & Intègre un contre-argument (réfutation) ou nuance \\ \hline
Cause & \emph{car, parce que, puisque, comme, du fait que} & Introduit l'argument (la raison) \\ \hline
Conséquence & \emph{donc, par conséquent, c'est pourquoi, ainsi, de sorte que} & Introduit la thèse ou une conclusion partielle \\ \hline
Illustration & \emph{par exemple, ainsi, notamment, c'est le cas de} & Signale un exemple \\ \hline
Conclusion & \emph{en définitive, en somme, pour conclure, au total} & Récapitule la thèse confirmée \\ \hline
\end{tabularx}
\end{center}

\begin{methode}[Repérer la structure argumentative en 3 questions]
\begin{enumerate}
    \item \textbf{Que défend l'auteur ?} $\to$ Cherchez la thèse (souvent en introduction ou conclusion, parfois au milieu). Formulez-la en une phrase affirmative.
    \item \textbf{Avec quels arguments ?} $\to$ Segmentez le corps du texte : chaque paragraphe (ou groupe de paragraphes) développe généralement un argument principal. Notez le type de chaque argument.
    \item \textbf{Quels exemples ?} $\to$ Pour chaque argument, identifiez l'illustration concrète (chiffre, cas, citation). Décidez si elle est \emph{essentielle} (chiffre clé, exemple unique) ou \emph{accessoire} (anecdote de soutien).
\end{enumerate}
\end{methode}

\courssub{Les grands schémas de raisonnement}

Au-delà des connecteurs, l'argumentation suit une logique globale.

\begin{propriete}[Types de raisonnement]
\begin{itemize}
    \item \textbf{Raisonnement déductif :} Du général au particulier. \emph{Principe général $\to$ Application au cas d'espèce.} (Ex. : « Tout homme est mortel. Socrate est un homme. Donc Socrate est mortel. »)
    \item \textbf{Raisonnement inductif :} Du particulier au général. \emph{Observation de cas $\to$ Généralisation.} (Ex. : « Les routes X, Y, Z sont dégradées. Donc le réseau routier national est en mauvais état. »)
    \item \textbf{Raisonnement par analogie :} Rapprochement de deux situations semblables pour en tirer une même conclusion. (Ex. : « Comme le Cameroun a réussi la CAN 2021, il peut réussir l'organisation des Jeux de la Francophonie. »)
\end{itemize}
\end{propriete}

\begin{propriete}[Procédés de réfutation (antithèse)]
L'auteur peut anticiper les objections :
\begin{itemize}
    \item \textbf{Concession + retournement :} \emph{« Certes, le coût est élevé, mais l'inaction coûtera plus cher. »}
    \item \textbf{Réduction à l'absurde :} Montrer que l'opinion adverse mène à une conséquence inacceptable.
    \item \textbf{Contre-argument factuel :} Apporter une preuve qui contredit l'objection.
\end{itemize}
\end{propriete}

\courssub{Application : analyse d'un éditorial camerounais}

Prenons un éditorial type (texte simulé, représentatif du corpus de presse) sur le thème \emph{« L'urgence de la réhabilitation de l'axe Yaoundé–Douala »}.

\begin{exoresolu}[Analyse de structure argumentative]
\textbf{Énoncé :} Voici un extrait d'éditorial. Dégagez la thèse, les arguments (avec leur type), les exemples et les connecteurs structurants.

\textit{« \textbf{Notre poumon économique s'asphyxie.} (1) \textbf{Car} la route Yaoundé–Douala, artère vitale qui draine 60\,\% du fret national, est devenue un parcours du combattant. (2) \textbf{De plus}, l'état de la chaussée engendre une hécatombe silencieuse : selon le Ministère des Transports, 47 morts rien qu'au premier semestre 2024. (3) \textbf{Certes}, le gouvernement a lancé le projet de doublement de la voie. (4) \textbf{Mais} les délais s'éternisent et le budget alloué reste insuffisant face à l'ampleur des travaux. (5) \textbf{Par conséquent}, il est impératif de déclarer l'état d'urgence routière et de mobiliser des financements innovants (PPP, diaspora). (6) \textbf{En définitive}, la prospérité du pays passe par des infrastructures à la hauteur de nos ambitions. (7) »}

\tcblower
\textbf{Solution détaillée :}

\begin{center}
\begin{tabularx}{\linewidth}{|l|X|X|l|}
\hline
\rowcolor{popblueL}
\textbf{Segment} & \textbf{Contenu} & \textbf{Rôle structurel} & \textbf{Type d'argument} \\ \hline
(1) & « Notre poumon économique s'asphyxie. » & \textbf{Thèse} (métaphorique) & — \\ \hline
(2) & « \textbf{Car} la route... draine 60\,\% du fret... » & \textbf{Argument 1} : Importance économique & \textbf{Argument de fait / chiffre} (autorité : statistique) \\ \hline
(3) & « \textbf{De plus}, l'état... 47 morts... » & \textbf{Argument 2} : Coût humain & \textbf{Argument d'expérience / statistique officielle} (autorité : Ministère) \\ \hline
(4)-(5) & « \textbf{Certes}... \textbf{Mais}... délais... budget insuffisant. » & \textbf{Réfutation (Antithèse)} : Concession sur l'effort + Retour à la thèse (insuffisance) & \textbf{Argument pragmatique} (efficacité de l'action actuelle) \\ \hline
(6) & « \textbf{Par conséquent}, il est impératif... » & \textbf{Conclusion / Appel à l'action} (reformulation thèse + proposition) & \textbf{Argument de valeur} (prospérité, ambition) \\ \hline
(7) & « \textbf{En définitive}, la prospérité... » & \textbf{Conclusion générale} & — \\ \hline
\end{tabularx}
\end{center}

\textbf{Structure globale :} Thèse $\to$ Argument 1 (économie) + Exemple (60\,\%) $\to$ Argument 2 (sécurité) + Exemple (47 morts) $\to$ Concession/Réfutation (projet existant mais lent) $\to$ Conclusion impérative (urgence, financement) $\to$ Conclusion valeur.
\end{exoresolu}

\begin{popfigure}
\begin{tikzpicture}[
    box/.style={draw=popdark, thick, rounded corners, inner sep=6pt, font=\small, align=center, minimum height=1.8cm},
    arrow/.style={->, thick, draw=popdark, shorten >=2pt, shorten <=2pt}
]
% Zones colorées verticales
\node[box, fill=popblue!15, minimum width=4.5cm, align=center] (th) at (0,0){\textbf{ZONE THÈSE}\\(1) « Notre poumon économique s'asphyxie. »};
\node[box, fill=popgreen!15, minimum width=4.5cm, align=center] (a1) at (5.5,0){\textbf{ZONE ARGUMENT 1}\\(2) « Car... 60\,\% fret »\\Type : Fait / Chiffre};
\node[box, fill=popgreen!15, minimum width=4.5cm, align=center] (a2) at (11,0){\textbf{ZONE ARGUMENT 2}\\(3) « De plus... 47 morts »\\Type : Expérience / Stat. off.};
\node[box, fill=poporange!15, minimum width=4.5cm, align=center] (ref) at (16.5,0){\textbf{ZONE RÉFUTATION}\\(4-5) « Certes... Mais... »\\Concession + Contre-argument};
\node[box, fill=poppurple!15, minimum width=4.5cm, align=center] (conc) at (22,0){\textbf{ZONE CONCLUSION}\\(6-7) « Par conséquent... En définitive »\\Appel à l'action + Valeur};

\draw[arrow] (th.east) -- (a1.west) node[midway, above, font=\tiny, text=popdark] {Car};
\draw[arrow] (a1.east) -- (a2.west) node[midway, above, font=\tiny, text=popdark] {De plus};
\draw[arrow] (a2.east) -- (ref.west) node[midway, above, font=\tiny, text=popdark] {Certes / Mais};
\draw[arrow] (ref.east) -- (conc.west) node[midway, above, font=\tiny, text=popdark] {Par conséquent};

% Flèches vers exemples (simplifié)
\node[font=\tiny, text=popmuted] at (5.5, -1.5) {Exemple : 60\%};
\node[font=\tiny, text=popmuted] at (11, -1.5) {Exemple : 47 morts};
\node[font=\tiny, text=popmuted] at (16.5, -1.5) {Exemple : Projet existant};
\end{tikzpicture}
\legende{Schéma d'un éditorial annoté : zones colorées (Thèse / Arguments / Réfutation / Conclusion) avec connecteurs et types d'arguments.}
\end{popfigure}

\courssub{Lien direct avec la contraction de texte (résumé)}

C'est le point crucial pour l'examen : \textbf{le résumé est la projection fidèle de la structure argumentative.}

\begin{aretenir}[Règle d'or du résumé argumentatif]
\begin{enumerate}
    \item La \textbf{thèse} source devient le \textbf{sujet} (ou la phrase d'ouverture) du résumé.
    \item Les \textbf{arguments} principaux deviennent les \textbf{idées forces} du résumé (un paragraphe ou une phrase complexe par argument).
    \item Les \textbf{exemples} sont \textbf{filtrés} : on ne garde que l'indicateur chiffré clé ou l'illustration unique. On écrit « \emph{47 morts au 1er semestre 2024} » et non le récit détaillé.
    \item La \textbf{réfutation} (concession + retournement) doit apparaître pour montrer l'honnêteté intellectuelle de l'auteur.
    \item Les \textbf{connecteurs} logiques sont conservés ou remplacés par leurs équivalents neutres (\emph{car $\to$ parce que ; par conséquent $\to$ donc}).
\end{enumerate}
\end{aretenir}

\begin{exemplebox}[Résumé simulé de l'éditorial ci-dessus (≈ 60 mots)]
\textit{« L'éditorial alerte sur l'asphyxie économique due à la dégradation de la route Yaoundé–Douala, artère supportant 60\,\% du fret national. Il argue du coût humain (47 morts en six mois selon le Ministère des Transports). Si un projet de doublement existe, son exécution lente et sous-financée impose de décréter l'urgence routière et d'innover dans le financement, condition de la prospérité nationale. »}

\textbf{Analyse :} Thèse conservée (asphyxie économique). 2 arguments fusionnés (économie + sécurité) avec leurs chiffres clés. Réfutation résumée (projet existe mais insuffisant). Conclusion action/valeur conservée. Connecteurs : \emph{Il argue de... Si... impose de... condition de...}
\end{exemplebox}

\courssub{Exercice d'entraînement : le téléphone portable en classe}

\begin{exercice}[Structure argumentative : « Pour ou contre le portable en classe ? »]
Voici un court texte argumentatif. Appliquez la \textbf{Méthode en 3 questions} pour en extraire la structure.

\textit{« Interdire le téléphone portable au lycée est une erreur pédagogique. \textbf{D'une part}, cet outil est devenu un accès portable au savoir : dictionnaires, encyclopédies, applications de révision y sont gratuits. \textbf{Par exemple}, l'application « Khan Academy » permet à un élève de Terminale de réviser ses maths dans le bus. \textbf{D'autre part}, l'interdiction totale crée une fracture numérique : les élèves défavorisés n'ont souvent que leur téléphone comme ordinateur. \textbf{Certains diront} que cela distrait. \textbf{Mais} c'est à l'école d'éduquer à l'usage responsable, non de fuir l'outil. \textbf{En somme}, intégrer le numérique, c'est préparer l'avenir. »}

\begin{enumerate}
    \item Quelle est la thèse ? (Formulez en une phrase affirmative).
    \item Combien d'arguments principaux ? Identifiez-les et donnez leur type.
    \item Quels sont les exemples ? Lequel est indispensable ?
    \item Repérez la concession et le retournement.
    \item Proposez un résumé en 35 mots maximum.
\end{enumerate}
\end{exercice}

\begin{corrige}[Exercice n°1]
\begin{enumerate}
    \item \textbf{Thèse :} L'interdiction du portable au lycée est une erreur ; il faut l'intégrer pédagogiquement.
    \item \textbf{2 arguments principaux :}
        \begin{itemize}
            \item Arg 1 : Outil d'accès au savoir (Type : \textbf{Pragmatique / Utilité}).
            \item Arg 2 : Lutte contre la fracture numérique / équité (Type : \textbf{Valeur / Justice sociale}).
        \end{itemize}
    \item \textbf{Exemples :} Khan Academy (Arg 1) ; « seul ordinateur des défavorisés » (Arg 2). \textbf{Indispensable :} La référence à la fracture numérique (argument de valeur fort) ; Khan Academy est illustratif mais remplaçable par « applications éducatives gratuites ».
    \item \textbf{Concession :} « Certains diront que cela distrait. » \textbf{Retournement :} « Mais c'est à l'école d'éduquer à l'usage responsable... »
    \item \textbf{Résumé (33 mots) :} « L'auteur refuse l'interdiction du portable en classe : c'est un outil d'accès au savoir (ex. Khan Academy) et un levier d'équité pour les élèves sans ordinateur. La distraction s'éduque, elle ne s'interdit pas. Intégrer le numérique prépare l'avenir. »
\end{enumerate}
\end{corrige}

\courssub{Conseil d'expert pour le Probatoire / Baccalauréat}

\begin{savaistu}[La structure argumentative, critère noté au Bac]
Au Probatoire et au Baccalauréat technique, l'épreuve de contraction de texte (résumé) est notée sur \textbf{la fidélité au contenu ET au plan}. Le correcteur vérifie :
\begin{itemize}
    \item La thèse est-elle clairement identifiée et placée en ouverture ?
    \item Les arguments principaux sont-ils tous présents, dans l'ordre, sans confusion avec les exemples ?
    \item La concession/réfutation est-elle mentionnée (souvent oubliée par les candidats) ?
    \item Les connecteurs logiques assurent-ils la cohérence du résumé ?
\end{itemize}
\textbf{Astuce :} Avant d'écrire, tracez l'arbre (Thèse $\to$ Arg 1, Arg 2... $\to$ Exemples clés) sur votre brouillon. C'est votre « plan de vol ». Ne commencez la rédaction qu'une fois l'arbre validé.
\end{savaistu}

\begin{attention}[Erreur classique : « L'exemple tue l'argument »]
Beaucoup d'élèves écrivent : « L'auteur dit que Khan Academy aide les élèves. » $\to$ \textbf{Zéro point pour l'argument.} On a gardé l'exemple, perdu l'argument (accès au savoir / utilité).
\textbf{Correct :} « L'auteur souligne l'utilité pédagogique du portable (accès aux ressources gratuites type Khan Academy). »
\textbf{Règle :} L'argument est le \emph{message}, l'exemple est la \emph{preuve}. Le résumé transmet le message, il signale la preuve.
\end{attention}

\courssub{Synthèse de la section}

\begin{definition}[Structure argumentative (vue synthétique)]
C'est l'organisation hiérarchisée (\textbf{Thèse} $\to$ \textbf{Arguments} $\to$ \textbf{Exemples}) reliée par des \textbf{connecteurs logiques} (addition, opposition, cause, conséquence) et animée par un \textbf{raisonnement} (déductif, inductif, analogique) qui peut intégrer une \textbf{réfutation}. Maîtriser cette architecture est la condition \emph{sine qua non} d'un résumé fidèle, cohérent et bien noté.
\end{definition}

La section suivante (« Techniques de rédaction du résumé ») transformera cette analyse structurelle en gestes rédactionnels concrets : reformulation, neutralisation, gestion de la longueur, liaison des phrases.

\coursec{Techniques de rédaction du résumé}

Dans la section précédente, nous avons appris à démonter la \cle{structure argumentative} d'un texte : repérer la thèse, suivre les arguments, identifier les exemples et les articulations logiques. Ce travail d'analyse n'est pas une fin en soi : il est la condition indispensable de l'exercice que nous abordons maintenant, la \cle{contraction de texte}. En effet, on ne peut condenser fidèlement que ce que l'on a d'abord parfaitement compris. Le résumé est l'épreuve reine de la production d'écrits au Probatoire et au Baccalauréat technique : il vérifie à la fois votre compréhension de lecture et votre maîtrise de l'expression écrite. Cette section vous donne, étape par étape, toutes les techniques pour réussir.

\courssub{Qu'est-ce qu'un résumé ? Définition et règles d'or}

Commençons par une situation de la vie courante. Un ami n'a pas pu assister à une réunion d'une heure du conseil de classe. Il vous demande : « Raconte-moi l'essentiel. » Vous ne pouvez pas tout redire : vous sélectionnez les décisions importantes, vous laissez de côté les anecdotes, les répétitions, les digressions, et vous rapportez le tout en quelques phrases, sans donner votre propre avis. Vous venez, sans le savoir, de faire un résumé. L'exercice scolaire repose exactement sur le même principe, appliqué à un texte écrit.

\begin{definition}[Le résumé]
Le \cle{résumé} est la reproduction \textbf{fidèle} et \textbf{condensée} des idées essentielles d'un texte. Il en respecte le sens, l'ordre et l'équilibre, sans aucun ajout, sans aucune appréciation personnelle, et il est rédigé entièrement avec les mots du résuméur. Sa longueur est fixée par la consigne : au Baccalauréat technique, elle correspond environ au \textbf{quart} du texte original.
\end{definition}

Cette définition contient quatre règles d'or qu'il faut graver dans sa mémoire :

\begin{enumerate}
\item \textbf{La règle du volume} : le résumé fait environ le quart du texte. Pour un texte de 400 mots, on attend environ 100 mots.
\item \textbf{La règle de fidélité} : le résumé ne trahit pas la pensée de l'auteur. Il conserve sa thèse, ses arguments, et même l'ordre de présentation des idées.
\item \textbf{La règle de neutralité} : aucune appréciation personnelle (« je trouve que », « l'auteur a raison », « cet exemple est faible ») ne doit apparaître. Le résuméur s'efface totalement.
\item \textbf{La règle de la reformulation} : le résumé est rédigé avec vos propres mots. Aucune citation longue n'est admise ; le « copier-coller » de phrases entières est une faute grave.
\end{enumerate}

\begin{savaistu}[La règle du quart]
La norme du quart ($\dfrac{1}{4}$ du volume du texte) est la règle traditionnelle de l'épreuve camerounaise de contraction de texte au Baccalauréat. Le sujet précise toujours le nombre de mots attendu (par exemple : « Résumez ce texte en 100 mots »), avec une tolérance de $\pm 10\,\%$. Un résumé de 90 ou 110 mots pour 100 mots demandés reste donc accepté ; un résumé de 70 ou 150 mots est sanctionné.
\end{savaistu}

\courssub{Étape 1 — La lecture analytique : comprendre avant de condenser}

Aucun résumé ne peut être réussi sans une lecture méthodique. Cette première étape, silencieuse, occupe souvent un tiers du temps de l'épreuve : c'est un investissement rentable.

\begin{methode}[La lecture analytique en trois passages]
\begin{enumerate}
\item \textbf{Première lecture globale} (sans stylo) : lire le texte en entier pour saisir le sens général. À la fin, formuler mentalement en une phrase : « De quoi parle ce texte ? » C'est le \cle{thème}.
\item \textbf{Deuxième lecture active} (crayon à la main) : souligner ou surligner de manière raisonnée. On repère : la \cle{thèse} (l'idée principale défendue, souvent au début ou à la fin), les \textbf{arguments} (introduits par « d'abord », « en effet », « ensuite », « enfin »), les \textbf{connecteurs logiques} qui structurent le propos (« mais », « cependant », « donc », « par conséquent »).
\item \textbf{Troisième lecture de contrôle} : vérifier que chaque paragraphe peut être ramené à une idée unique. On écrit en marge, pour chaque paragraphe, une étiquette de 3 à 5 mots. Ces étiquettes formeront le plan du résumé.
\end{enumerate}
\end{methode}

\begin{attention}[Le surlignage n'est pas un coloriage]
L'erreur classique consiste à surligner des phrases entières, voire des paragraphes complets. Un surlignage raisonné ne retient que les \textbf{mots porteurs de sens} : le sujet de l'idée, le verbe d'action, le connecteur. Si plus d'un quart du texte est surligné, c'est que rien n'a été sélectionné : vous avez recopié le texte au lieu de l'analyser.
\end{attention}

\courssub{Étape 2 — Les quatre techniques de contraction}

Une fois les idées essentielles dégagées, il faut réduire le volume. Quatre techniques, à combiner, permettent cette réduction. Ce sont les outils du résuméur, comme le rabot et la scie sont les outils du menuisier.

\textbf{1. La suppression.} C'est la technique la plus rentable. On élimine purement et simplement :
\begin{itemize}
\item les \textbf{exemples} et illustrations (souvent introduits par « par exemple », « ainsi », « c'est le cas de ») ;
\item les \textbf{répétitions} et paraphrases internes (l'auteur dit deux fois la même chose sous deux formes) ;
\item les \textbf{détails} secondaires : chiffres précis, dates accessoires, anecdotes, énumérations descriptives ;
\item les \textbf{incises} et formules de politesse du discours (« il convient de rappeler que », « comme chacun sait »).
\end{itemize}

\textbf{2. La substitution.} On remplace une expression longue par un mot unique de même sens : « d'une manière rapide » devient « rapidement » ; « les personnes qui exercent le métier d'enseigner » devient « les enseignants » ; « mettre fin à » remplace « faire cesser d'une manière définitive ».

\textbf{3. La fusion.} On regroupe en une seule phrase deux ou trois phrases voisines qui expriment des idées proches, souvent à l'aide d'un connecteur (« et », « car », « donc »). Deux phrases : « Les élèves arrivent en retard. Ils manquent souvent le début des cours » fusionnent en : « Les élèves retardataires manquent le début des cours. »

\textbf{4. La généralisation.} On remplace une énumération par un terme générique : « les machettes, les houes, les pioches et les pelles » devient « les outils agricoles » ; « les pères, les mères, les enfants et les grands-parents » devient « la famille ».

\begin{center}
\begin{tabularx}{\linewidth}{|l|X|X|}\hline
\rowcolor{popblueL} \textbf{Technique} & \textbf{Phrase originale} & \textbf{Phrase contractée} \\ \hline
Suppression & L'agriculture, par exemple dans les villages de l'Ouest où l'on cultive le maïs depuis des générations, nourrit les familles. & L'agriculture nourrit les familles. \\ \hline
Substitution & Les jeunes quittent le village d'une manière massive. & Les jeunes quittent massivement le village. \\ \hline
Fusion & Les routes sont dégradées. Le transport des récoltes devient difficile. & Les routes dégradées rendent le transport des récoltes difficile. \\ \hline
Généralisation & Le maïs, le manioc, les arachides et les haricots se vendent mal. & Les produits agricoles se vendent mal. \\ \hline
\end{tabularx}
\end{center}

\courssub{Étape 3 — La reformulation : dire avec ses propres mots}

La reformulation est le passage délicat de l'exercice. Il s'agit de rendre la pensée de l'auteur avec votre propre vocabulaire et vos propres constructions de phrases, sans jamais en changer le sens. Deux garde-fous s'imposent :

\begin{itemize}
\item \textbf{Interdiction du copier-coller} : reprendre des phrases entières du texte, même bien choisies, est sanctionné. Seuls les termes techniques irremplaçables (« exode rural », « démocratie ») peuvent être conservés tels quels.
\item \textbf{Interdiction de la trahison} : reformuler ne signifie pas déformer. Si l'auteur écrit « l'école reste un facteur d'intégration », on ne peut pas résumer « l'école intègre parfaitement » : la nuance de prudence (« reste un facteur ») fait partie du sens.
\end{itemize}

\begin{attention}[La paraphrase n'est pas un résumé]
L'erreur la plus fréquente et la plus coûteuse est la \cle{paraphrase} : réécrire le texte phrase par phrase, avec d'autres mots, mais en conservant à peu près le même nombre de mots. Le correcteur le repère immédiatement : le volume n'a pas été réduit. Un résumé \textbf{supprime} des idées secondaires ; une paraphrase les conserve toutes en les déguisant. Contrôlez-vous : si votre production dépasse nettement le quart du texte, vous avez paraphrasé, pas résumé.
\end{attention}

\courssub{Étape 4 — Le comptage des mots et l'ajustement}

Le résumé doit respecter le nombre de mots imposé, avec une tolérance de $\pm 10\,\%$. Il faut donc compter, puis annoncer le total à la fin de la copie (par exemple : « 98 mots »).

\begin{methode}[Compter vite et juste]
\begin{enumerate}
\item Compter les mots de la première ligne complète (par exemple 10 mots).
\item Multiplier par le nombre de lignes complètes : 9 lignes $\times$ 10 mots $= 90$ mots.
\item Ajouter les mots de la dernière ligne incomplète : $90 + 8 = 98$ mots.
\item Si le total dépasse la fourchette, supprimer un détail ou fusionner deux phrases ; s'il est trop court, réintégrer une idée essentielle sacrifiée.
\end{enumerate}
\end{methode}

\courssub{Exercice guidé : de 120 mots à 30 mots}

Appliquons maintenant l'ensemble de la méthode à un paragraphe, comme on le ferait au tableau avec la classe.

\begin{exoresolu}[Contraction d'un paragraphe, étape par étape]
\textbf{Énoncé.} Résumez en environ 30 mots le paragraphe suivant (120 mots) :

« Dans de nombreux villages du Cameroun, par exemple dans les régions de l'Ouest et du Nord-Ouest, les jeunes gens âgés de vingt à trente ans abandonnent progressivement les champs de leurs parents pour se rendre dans les grandes villes comme Douala ou Yaoundé, attirés par l'espoir, souvent déçu, d'y trouver un emploi salarié et une vie plus moderne. Ce départ massif, qui s'accentue d'année en année, prive les campagnes de bras valides au moment même des semailles et des récoltes, si bien que les terres, autrefois cultivées avec soin par des générations de paysans expérimentés, restent désormais en friche et que la production agricole diminue dangereusement. »

\tcblower
\textbf{Solution détaillée.}

\textbf{Étape 1 — Lecture analytique.} Thème : l'exode rural au Cameroun. Idées essentielles : (a) les jeunes quittent les villages pour les villes ; (b) cause : l'espoir d'un emploi ; (c) conséquence : les terres sont abandonnées et la production baisse. À supprimer : les exemples (« régions de l'Ouest », « Douala ou Yaoundé »), les détails (« vingt à trente ans », « semailles et récoltes », « générations de paysans expérimentés »), les répétitions (« d'année en année », « autrefois... désormais »).

\textbf{Étape 2 — Contraction.}
\begin{itemize}
\item Suppression des exemples et détails : « par exemple dans les régions de l'Ouest et du Nord-Ouest », « âgés de vingt à trente ans », « comme Douala ou Yaoundé », « souvent déçu » disparaissent.
\item Substitution : « abandonnent progressivement les champs de leurs parents » devient « quittent les campagnes ».
\item Fusion : les deux dernières propositions (« les terres restent en friche » + « la production diminue ») fusionnent en une seule conséquence.
\item Généralisation : « un emploi salarié et une vie plus moderne » devient « un emploi ».
\end{itemize}

\textbf{Étape 3 — Rédaction.} « Au Cameroun, les jeunes quittent les campagnes pour les villes, espérant y trouver un emploi. Cet exode rural prive les villages de main-d'œuvre : les terres restent en friche et la production agricole diminue. »

\textbf{Étape 4 — Comptage.} 31 mots. La consigne (30 mots $\pm 10\,\%$, soit 27 à 33 mots) est respectée. Chaque coupure est justifiée : aucune idée essentielle n'a été perdue, aucune idée n'a été ajoutée.
\end{exoresolu}

\begin{popfigure}
\begin{center}
\begin{tikzpicture}[font=\small, scale=0.9]
\fill[popblue!35] (0,0) -- (8,0) -- (6.6,-1.8) -- (1.4,-1.8) -- cycle;
\fill[poporange!40] (1.9,-2.1) -- (6.1,-2.1) -- (5.2,-3.6) -- (2.8,-3.6) -- cycle;
\fill[popgreen!45] (3.2,-3.9) -- (4.8,-3.9) -- (4.4,-5.4) -- (3.6,-5.4) -- cycle;
\node at (4,-0.9) {\textbf{Texte original : 400 mots}};
\node at (4,-2.85) {\textbf{Idées essentielles}};
\node at (4,-4.65) {\textbf{Résumé : 100 mots}};
\draw[->,very thick,popdark] (8.6,-0.5) -- (8.6,-5) node[midway,right,align=left] {suppression\\substitution\\fusion\\généralisation};
\end{tikzpicture}
\end{center}
\legende{Schéma en entonnoir de la contraction : le texte de 400 mots est filtré par la lecture analytique (qui dégage les idées essentielles), puis condensé par les quatre techniques jusqu'au résumé d'environ 100 mots, soit le quart du volume initial.}
\end{popfigure}

\courssub{Les erreurs à éviter : bilan de la méthode}

Pour terminer, dressons la liste des fautes qui coûtent le plus de points à l'examen, afin de les éliminer définitivement de vos copies :

\begin{itemize}
\item \textbf{La paraphrase} : réécrire tout le texte sans réduire le volume (voir l'encadré d'attention ci-dessus).
\item \textbf{L'ajout d'idées} : introduire des exemples ou des arguments qui ne figurent pas dans le texte, même s'ils sont vrais et pertinents.
\item \textbf{Le jugement personnel} : commenter, approuver ou critiquer la thèse de l'auteur. Le résumé est neutre ; le commentaire, lui, sera un autre exercice.
\item \textbf{Le déséquilibre entre les parties} : résumer longuement le premier paragraphe et expédier la fin. Le résumé doit respecter les proportions du texte : si un tiers du texte traite des conséquences, un tiers du résumé doit leur être consacré.
\item \textbf{Le non-respect du volume} : oublier de compter et d'indiquer le nombre de mots.
\end{itemize}

\begin{methode}[Récapitulatif : la méthode du résumé en 5 étapes]
\begin{enumerate}
\item \textbf{Lire} le texte deux ou trois fois, en surlignant de manière raisonnée thème, thèse et connecteurs.
\item \textbf{Dégager le plan} : une idée par paragraphe, notée en marge.
\item \textbf{Contracter} au brouillon par suppression, substitution, fusion et généralisation.
\item \textbf{Rédiger} avec ses propres mots, dans l'ordre du texte, sans avis personnel.
\item \textbf{Compter} les mots, ajuster à $\pm 10\,\%$ et indiquer le total.
\end{enumerate}
\end{methode}

\begin{exercice}[À vous de jouer]
Résumez en environ 25 mots le passage suivant (100 mots) : « Le téléphone portable, que l'on trouve désormais partout, dans les grandes villes comme dans les villages les plus reculés, a profondément transformé la vie des Camerounais. Grâce à lui, par exemple, un commerçant de Bafoussam peut passer commande à son fournisseur de Douala sans se déplacer, une mère de famille peut recevoir de l'argent envoyé par son fils resté en ville, et les élèves échangent des informations sur leurs cours. Cet outil, devenu indispensable, a ainsi réduit les distances et facilité les échanges quotidiens. »
\end{exercice}

\begin{corrige}[Exercice]
Pistes de correction : idées essentielles — (a) le portable s'est généralisé au Cameroun ; (b) il facilite le commerce, les transferts d'argent et la communication ; (c) il a réduit les distances. Suppression des exemples localisés (Bafoussam, Douala), généralisation des trois usages. Proposition de résumé (26 mots) : « Généralisé au Cameroun, le téléphone portable facilite le commerce, les transferts d'argent et la communication : devenu indispensable, il a réduit les distances et simplifié la vie quotidienne. »
\end{corrige}

La maîtrise de ces techniques de contraction acquise, il reste à les mettre en œuvre sur un texte complet : c'est l'objet de la prochaine section, où nous établirons le plan d'un texte entier, rédigerons le résumé en temps limité, puis confronterons et améliorerons nos productions.

\coursec{Plan du texte, rédaction, confrontation et amélioration}

Dans la section précédente, nous avons appris les \cle{techniques de rédaction du résumé} : repérer la thèse, sélectionner les arguments essentiels, supprimer les exemples et les répétitions, reformuler avec ses propres mots et compter les mots. Mais une technique, aussi bien connue soit-elle, ne suffit pas : beaucoup de candidats au Probatoire et au Baccalauréat technique échouent non pas parce qu'ils ignorent les règles, mais parce qu'ils se lancent dans la rédaction sans organisation préalable. Cette dernière section vous donne la \cle{démarche complète de production} : construire le plan du texte, rédiger le résumé à partir de ce plan, confronter sa production à celle des camarades, puis l'améliorer. C'est cette boucle « planifier -- produire -- confronter -- améliorer » qui transforme un élève moyen en bon résumeur.

\courssub{Élaborer le plan du texte}

\begin{definition}[Plan du texte]
Le \cle{plan du texte} est la représentation schématique de la structure argumentative d'un texte : il dégage les \textbf{grandes parties} (les étapes du raisonnement de l'auteur) et, si nécessaire, les \textbf{sous-parties} (les arguments qui soutiennent chaque étape). Il ne contient que des idées, jamais d'exemples ni de citations.
\end{definition}

L'intuition est simple : un texte argumentatif est comme une maison. Avant de décrire la maison en quelques phrases, il faut d'abord voir combien elle a de pièces et comment elles s'enchaînent. Le plan, ce sont les pièces ; le résumé, c'est la visite guidée rapide de ces pièces.

Pour construire le plan, on s'appuie sur la \cle{structure argumentative} étudiée dans la section 4 : la thèse (ce que l'auteur veut démontrer), les arguments (ce qui prouve la thèse) et la position finale (ce que l'auteur conclut). Les \textbf{connecteurs logiques} sont de précieux indicateurs : « d'abord », « ensuite », « en effet », « cependant », « enfin », « c'est pourquoi » signalent souvent le passage d'une partie à une autre.

\begin{methode}[Construire le plan du texte]
\begin{enumerate}
\item Relire le texte en entier et formuler en une phrase la \textbf{thèse} de l'auteur.
\item Découper le texte en \textbf{moments du raisonnement} : à chaque fois que l'auteur change d'idée force, on ouvre une nouvelle partie. Principe fondamental : \emph{une idée = une partie}.
\item Donner à chaque partie un \textbf{titre bref} (quelques mots) qui résume l'idée, sans recopier une phrase du texte.
\item Ranger sous chaque partie les \textbf{arguments} qui la soutiennent (sous-parties), en éliminant exemples, chiffres et anecdotes.
\item \textbf{Vérifier que le plan « tient debout » sans le texte} : lire uniquement les titres des parties. Si l'enchaînement est logique et si l'on comprend le cheminement de l'auteur, le plan est bon. Sinon, il faut fusionner ou redécouper.
\end{enumerate}
\end{methode}

\begin{attention}[Erreur fréquente : rédiger sans plan]
Beaucoup d'élèves rédigent le résumé directement, au fil de la lecture. Résultat : les premières lignes du texte occupent la moitié du résumé, la fin est expédiée en une phrase, et la thèse est déformée. Un résumé \textbf{déséquilibré} est presque toujours le signe d'un plan absent. Prenez le temps de construire le plan : cinq minutes de planification économisent quinze minutes de ratures.
\end{attention}

\courssub{Du plan au résumé : la rédaction}

Une fois le plan établi, la rédaction devient presque mécanique : \textbf{chaque partie du plan devient une ou deux phrases du résumé}. Le résumé complet comporte trois moments.

\begin{propriete}[Structure du résumé rédigé]
\begin{itemize}
\item \textbf{Introduction brève} (1 à 2 phrases) : présenter le \textbf{thème} du texte et la \textbf{thèse} de l'auteur. Formule utile : « Dans ce texte, l'auteur s'interroge sur... et soutient que... ».
\item \textbf{Développement} : rendre compte des \textbf{arguments}, dans l'ordre du texte, en une ou deux phrases par partie du plan, reliées par des connecteurs logiques (d'abord, ensuite, en effet, cependant, enfin).
\item \textbf{Conclusion} (1 phrase) : formuler la \textbf{position finale} de l'auteur, sa conclusion ou sa proposition.
\end{itemize}
Le résumé est rédigé à la troisième personne, au présent, sans jugement personnel, et respecte le nombre de mots imposé (en général le quart du texte, avec une marge de 10~\%).
\end{propriete}

\begin{exemplebox}[D'un plan à un résumé : l'entrepreneuriat des jeunes au Cameroun]
Texte argumentatif de 480 mots sur l'entrepreneuriat des jeunes. \textbf{Plan dégagé :}
\begin{enumerate}
\item Thèse : l'entrepreneuriat est une voie sérieuse contre le chômage des jeunes.
\item Argument 1 : l'emploi salarié public et privé ne suffit plus à absorber les diplômés.
\item Argument 2 : les jeunes disposent d'atouts (créativité, maîtrise du numérique, réseaux).
\item Argument 3 : des obstacles existent (financement, formation), mais des solutions émergent (incubateurs, microfinance).
\item Conclusion de l'auteur : l'État et les jeunes doivent miser sur l'entrepreneuriat.
\end{enumerate}
\textbf{Résumé obtenu (environ 120 mots) :} « Ce texte traite du chômage des jeunes au Cameroun et soutient que l'entrepreneuriat constitue une voie efficace pour le combattre. L'auteur montre d'abord que les emplois salariés, publics comme privés, ne peuvent plus absorber le nombre croissant de diplômés. Il fait ensuite valoir que les jeunes possèdent des atouts précieux : la créativité, la maîtrise des outils numériques et de solides réseaux relationnels. Certes, reconnaît-il, le manque de financement et de formation freine les initiatives ; mais des solutions apparaissent, telles que les incubateurs et la microfinance. En définitive, l'auteur invite l'État et les jeunes eux-mêmes à miser résolument sur l'entrepreneuriat. »
\end{exemplebox}

Observez la correspondance : les cinq parties du plan sont devenues cinq blocs de phrases, reliées par des connecteurs (« d'abord », « ensuite », « certes... mais », « en définitive »). Aucun exemple, aucun chiffre du texte d'origine n'a été conservé.

\courssub{Production individuelle en classe}

\begin{experience}[Séance de production : résumé d'un texte de 400 à 500 mots]
\textbf{Matériel :} un texte argumentatif inédit de 400 à 500 mots (éditorial, article de presse, extrait d'essai), une feuille de brouillon, une copie, un dictionnaire.

\textbf{Protocole (durée effective ~1 h 15, prévoir 1 h 30 avec consignes et installation) :}
\begin{enumerate}
\item Lecture silencieuse du texte, deux fois (5 min).
\item Soulignage de la thèse et des arguments ; annotation des connecteurs (10 min).
\item Construction du plan sur le brouillon (10 min).
\item Rédaction du résumé au propre, en visant le quart du texte, soit 100 à 125 mots (30 min).
\item Comptage précis des mots et ajustement (suppression ou ajout) (10 min).
\item Relecture : grammaire, orthographe, cohésion (10 min).
\end{enumerate}

\textbf{Observations attendues :} les élèves qui suivent le protocole produisent des résumés équilibrés et dans la fourchette de mots ; ceux qui sautent l'étape du plan dépassent souvent la limite ou déséquilibrent leur production.

\textbf{Interprétation :} le plan est l'outil de régulation de la production : il garantit la fidélité à la structure du texte et la répartition équilibrée des mots.
\end{experience}

\begin{exoresolu}[Comptage et ajustement d'un résumé]
Énoncé. Un texte comporte 460 mots ; la consigne impose un résumé d'environ 115 mots (marge de 10~\% acceptée). Un élève produit un résumé de 148 mots. Que doit-il faire ?

\tcblower
Solution détaillée. La fourchette admise va de 104 à 126 mots environ (115 $\pm$ 10~\%). Le résumé de 148 mots est donc trop long d'au moins 22 mots. L'élève applique les techniques de contraction : (1) il supprime un exemple qu'il avait gardé (« comme le montre le cas d'une jeune couturière de Bafoussam » : $-$10 mots) ; (2) il remplace un groupe de mots par un mot unique (« les personnes qui n'ont pas de travail » devient « les chômeurs » : $-$5 mots) ; (3) il fusionne deux phrases par un connecteur (« L'auteur évoque le financement. Il parle aussi de la formation. » devient « L'auteur évoque le financement et la formation. » : $-$6 mots). Nouveau total : 127 mots, proche de la limite ; une dernière compression (suppression d'un adjectif ou d'un adverbe) ramène le résumé à 125 mots. Le résumé est conforme.
\end{exoresolu}

\courssub{Confrontation des productions : la correction croisée}

Le résumé rédigé n'est pas terminé : il doit être \cle{confronté} à d'autres lectures du même texte. En classe, les copies sont échangées et chaque élève corrige la production d'un camarade à l'aide d'une \textbf{grille critériée}.

\begin{methode}[Correction croisée en deux lectures]
\begin{enumerate}
\item \textbf{Première lecture : la fidélité.} Le texte original sous les yeux, vérifier que la thèse est bien celle de l'auteur, qu'aucun argument essentiel ne manque, qu'aucune idée étrangère au texte n'a été ajoutée, qu'aucun jugement personnel n'apparaît. Annoter au crayon dans la marge : « thèse déformée », « argument 3 absent », « avis personnel ».
\item \textbf{Deuxième lecture : la langue et la forme.} Vérifier l'orthographe, la grammaire, la cohésion (connecteurs, reprises nominales), puis compter les mots. Annoter au crayon les fautes sans les corriger à la place du camarade.
\item \textbf{Attribution de la note} selon la grille, puis \textbf{retour oral} : expliquer à l'auteur au moins un point fort et un point à améliorer.
\end{enumerate}
\end{methode}

\begin{center}
\begin{tabularx}{\linewidth}{|l|X|c|}
\hline
\rowcolor{popblueL} \textbf{Critère} & \textbf{Indicateur} & \textbf{Barème} \\
\hline
Fidélité au texte & Thèse exacte, arguments essentiels présents, aucune idée ajoutée, aucun avis personnel & /6 \\
\hline
Condensation & Respect du nombre de mots imposé, suppression des exemples et répétitions, reformulation personnelle & /4 \\
\hline
Cohésion et cohérence & Enchaînement logique des idées, connecteurs appropriés, reprises nominales et pronominales correctes & /4 \\
\hline
Langue & Orthographe, grammaire, ponctuation, vocabulaire précis & /4 \\
\hline
Présentation et comptage & Texte lisible, nombre de mots indiqué et exact & /2 \\
\hline
\rowcolor{popgoldL} \textbf{Total} & & \textbf{/20} \\
\hline
\end{tabularx}
\end{center}

\begin{savaistu}[Confrontation et esprit critique]
La confrontation des productions n'est pas un simple échange de copies : elle développe l'\textbf{esprit critique} que l'approche par les compétences (APC), au cœur des programmes camerounais, place au centre des apprentissages. En jugeant le travail d'un autre selon des critères explicites, l'élève apprend à objectiver ses propres erreurs. Celui qui sait corriger le résumé d'un camarade sait bientôt corriger le sien : c'est le principe de l'\emph{évaluation formative}, qui fait de l'erreur un outil de progrès et non une sanction.
\end{savaistu}

\courssub{Amélioration, compte rendu et remédiation}

\begin{definition}[Amélioration de la production]
L'\cle{amélioration} est la réécriture de sa propre production à partir des remarques reçues pendant la confrontation. Elle s'accompagne d'une \textbf{justification orale} : l'élève explique pourquoi il a supprimé tel passage, reformulé telle phrase, choisi tel connecteur. Justifier ses choix de contraction, c'est prouver qu'ils sont délibérés et non le fruit du hasard.
\end{definition}

La démarche d'amélioration suit trois étapes : reprendre sa copie et les annotations du correcteur ; décider, remarque par remarque, si on l'accepte (et on corrige) ou si on la discute (et on s'explique) ; réécrire la version définitive en veillant à ce que le comptage reste dans la fourchette. La version améliorée est ensuite comparée à la première : l'écart entre les deux mesure le progrès réel.

\begin{popfigure}
\begin{center}
\begin{tikzpicture}[node distance=0.55cm and 0.9cm,
etape/.style={rectangle, rounded corners=3pt, draw=popblue, thick, fill=popblueL, text width=2.6cm, minimum height=1.1cm, align=center, font=\small},
fleche/.style={-{Stealth[length=3mm]}, thick, popdark}]
\node[etape] (t) {\textbf{Texte} argumentatif (400--500 mots)};
\node[etape, right=of t] (p) {\textbf{Plan} : thèse, parties, arguments};
\node[etape, right=of p] (b) {\textbf{Brouillon} puis résumé rédigé};
\node[etape, below=of b] (c) {\textbf{Confrontation} : correction croisée critériée};
\node[etape, left=of c] (a) {\textbf{Version améliorée} justifiée oralement};
\draw[fleche] (t) -- (p);
\draw[fleche] (p) -- (b);
\draw[fleche] (b) -- (c);
\draw[fleche] (c) -- (a);
\draw[fleche, dashed, poporange] (a.north) -- ++(0,0.55) -| node[pos=0.25, above, font=\scriptsize, text=poporange]{nouvelle production} (p.north);
\end{tikzpicture}
\end{center}
\legende{Le processus complet de production du résumé : du texte au plan, du plan au brouillon rédigé, puis à la confrontation et à la version améliorée. La flèche en pointillés rappelle que les acquis de l'amélioration servent à la production suivante.}
\end{popfigure}

Le \textbf{compte rendu de la séquence} clôt le travail : collectivement, la classe dresse au tableau la synthèse des acquis (savoir construire un plan, rédiger en trois moments, compter les mots, corriger selon une grille) et des difficultés persistantes. Ce bilan oral, guidé par le professeur, fixe ce qui est définitivement acquis et ce qui doit être retravaillé.

La \textbf{remédiation} propose alors des exercices ciblés selon les difficultés repérées.

\begin{exercice}[Remédiation : trois ateliers ciblés]
\begin{enumerate}
\item \textbf{Atelier comptage} (pour les résumés hors fourchette) : réduire de 150 à 120 mots un résumé fourni, en indiquant à chaque suppression la technique utilisée.
\item \textbf{Atelier reformulation} (pour les résumés qui recopient le texte) : reformuler cinq phrases du texte source sans reprendre plus de trois mots consécutifs de l'original.
\item \textbf{Atelier thèse} (pour les résumés infidèles) : à partir de trois courts éditoriaux, formuler en une phrase la thèse de chaque auteur, puis comparer avec celle d'un camarade.
\end{enumerate}
\end{exercice}

\begin{corrige}[Exercice : ateliers de remédiation]
\textbf{Atelier 1.} Les suppressions doivent porter d'abord sur les exemples et les répétitions, jamais sur la thèse ni sur un argument essentiel. Chaque mot retiré doit être justifié : « exemple supprimé », « synonyme unique », « fusion de phrases ».

\textbf{Atelier 2.} Exemple : « Les jeunes diplômés se heurtent à l'absence d'emplois dans le secteur public » peut devenir « Le secteur public n'offre plus de débouchés aux diplômés » : l'idée est conservée, la formulation est personnelle.

\textbf{Atelier 3.} La thèse se formule avec un verbe d'opinion ou d'affirmation : « l'auteur soutient que... », « l'auteur dénonce... », « l'auteur invite à... ». Si deux camarades formulent des thèses différentes, on retourne au texte pour trancher : c'est le texte qui a le dernier mot.
\end{corrige}

\begin{aretenir}[L'essentiel de la section]
\begin{itemize}
\item Le \textbf{plan du texte} dégage les grandes parties de la structure argumentative : une idée = une partie ; il doit « tenir debout » sans le texte.
\item Chaque partie du plan devient \textbf{une ou deux phrases} du résumé : introduction (thème + thèse), développement (arguments), conclusion (position finale).
\item La \textbf{confrontation} se fait par correction croisée en deux lectures (fidélité, puis langue), annotée au crayon, avec une grille critériée sur 20.
\item L'\textbf{amélioration} est une réécriture justifiée oralement ; le compte rendu collectif fixe les acquis ; la remédiation cible les difficultés : comptage, reformulation, repérage de la thèse.
\end{itemize}
\end{aretenir}

\newpage

\coursec{Activité d'intégration}

\coursec{Contraction de texte : rédiger un résumé de texte}

\courssub{Objectifs de l'activité}

À l'issue de cette activité d'intégration, chaque élève sera capable de :

\begin{itemize}
\item repérer la \cle{structure argumentative} d'un texte de presse (thèse, arguments, exemples, conclusion) ;
\item construire le \cle{plan du texte} à résumer avant toute rédaction ;
\item appliquer les quatre techniques de \cle{contraction} : suppression, substitution, fusion, généralisation ;
\item rédiger un résumé \cle{cohésif} et \cle{cohérent}, rédigé dans sa propre langue, sans copier-coller de phrases du texte ;
\item respecter une \cle{contrainte de longueur} (100 mots $\pm$ 10) et justifier ses choix de contraction ;
\item évaluer la production d'un pair à l'aide d'une grille critériée et améliorer sa propre copie.
\end{itemize}

\begin{attention}[Compétence visée au Baccalauréat technique]
La contraction de texte est une épreuve classique : on demande de résumer un texte argumentatif en un quart de sa longueur, sans commentaire personnel, sans citation, avec un comptage exact. Cette activité reproduit ces conditions.
\end{attention}

\courssub{Situation}

La rédaction du journal lycéen \emph{Le Clairon de Douala} prépare son numéro spécial « Débats de la nation ». Le chef de rubrique charge chaque élève de la classe de Terminale technique de résumer en \textbf{100 mots ($\pm$ 10)} l'éditorial suivant, paru dans un quotidien camerounais. Le résumé retenu sera publié dans la rubrique « L'essentiel en cent mots ».

\textbf{Texte à résumer (400 mots) :}

\begin{quote}
\emph{Faut-il rendre obligatoire l'enseignement des langues nationales au secondaire ?}

Depuis plusieurs années, la question revient avec insistance dans le débat public camerounais : nos lycées et collèges doivent-ils inscrire à leur emploi du temps, de manière obligatoire, l'enseignement des langues nationales ? À notre sens, la réponse ne souffre aucune hésitation : oui, et nous allons le démontrer.

D'abord, l'enseignement des langues nationales est un puissant facteur d'enracinement culturel. Un jeune qui maîtrise le ewondo, le duala, le fulfuldé ou le bassa comprend mieux les valeurs, les proverbes et l'histoire de sa communauté. À Douala comme à Maroua, combien de collégiens s'expriment couramment en français ou en anglais mais sont incapables de saluer leurs grands-parents dans la langue de leur village ? Cette rupture linguistique est une rupture identitaire. Or, une nation ne se construit pas sur des citoyens déracinés.

Ensuite, la recherche en éducation l'a montré : l'enfant apprend mieux lorsqu'il est d'abord instruit dans une langue qu'il maîtrise. Les notions abstraites, en mathématiques ou en sciences, sont plus facilement assimilées lorsque l'élève peut les reformuler dans sa langue première. Les expériences menées dans certaines écoles pilotes du Nord et de l'Est du pays l'ont confirmé : les élèves initiés à l'écrit dans une langue nationale progressent ensuite plus vite en français, langue officielle d'enseignement.

Enfin, promouvoir nos langues, c'est aussi préserver un patrimoine menacé. Sur plus de deux cent cinquante langues parlées au Cameroun, beaucoup risquent de disparaître avec la génération de nos aînés. Chaque langue qui meurt emporte avec elle des savoirs médicinaux, des récits, des techniques agricoles uniques. L'école, qui a largement contribué à marginaliser ces langues, a le devoir de participer à leur sauvegarde.

Certains objectent que l'enseignement des langues nationales coûterait cher : formation des maîtres, manuels scolaires, choix délicat des langues à enseigner dans des villes cosmopolites. D'autres craignent qu'il ne fragilise l'unité nationale ou ne retarde la maîtrise du français et de l'anglais, langues de la mobilité sociale. Ces objections sont réelles mais ne résistent pas à l'examen : le coût de la formation est un investissement, non une dépense perdue ; quant à l'unité nationale, elle se nourrit du respect des identités et non de leur négation.

En définitive, rendre obligatoire l'enseignement des langues nationales au secondaire, c'est offrir à la jeunesse camerounaise des racines solides, des apprentissages facilités et un patrimoine sauvegardé. L'État, les enseignants et les familles doivent s'y engager sans tarder.
\end{quote}

\courssub{Consignes}

\begin{enumerate}
\item Lisez deux fois l'éditorial. Soulignez la \cle{thèse} défendue, entourez les connecteurs logiques qui marquent la progression (d'abord, ensuite, enfin, certains, en définitive).
\item Complétez le tableau de la structure argumentative : thèse, argument 1 + exemple, argument 2 + exemple, argument 3 + exemple, objections réfutées, conclusion.
\item Établissez le \cle{plan du texte} en quatre ou cinq rubriques courtes, rédigées dans votre propre langue.
\item Rédigez une première version du résumé en appliquant les quatre techniques : \cle{suppression} (exemples, répétitions, énumérations), \cle{substitution} (un mot pour un groupe de mots), \cle{fusion} (deux phrases en une), \cle{généralisation} (le terme générique pour la liste des termes spécifiques).
\item Comptez les mots de votre résumé. Ajustez jusqu'à obtenir entre 90 et 110 mots. Indiquez le nombre exact entre parenthèses.
\item Échangez votre copie avec un voisin. Corrigez sa production à l'aide de la grille critériée ci-dessous, puis justifiez oralement vos annotations.
\item Rédigez la version définitive améliorée et présentez à la classe deux choix de contraction que vous avez opérés.
\end{enumerate}

\begin{center}
\begin{tabularx}{\linewidth}{|l|X|c|}
\hline
\rowcolor{popblueL} \textbf{Critère} & \textbf{Attendu} & \textbf{Points} \\
\hline
Fidélité & La thèse et les trois arguments essentiels sont restitués sans déformation ni opinion personnelle. & /5 \\
\hline
Sélection & Les exemples, énumérations et répétitions sont supprimés ; les idées essentielles conservées. & /4 \\
\hline
Reformulation & Aucune phrase copiée du texte ; substitutions et fusions visibles. & /4 \\
\hline
Cohésion & Connecteurs logiques, reprises pronominales, progressions thématiques assurant un texte fluide. & /4 \\
\hline
Contrainte & 90 à 110 mots, comptage indiqué, orthographe et syntaxe correctes. & /3 \\
\hline
\multicolumn{2}{|r|}{\textbf{Total}} & \textbf{/20} \\
\hline
\end{tabularx}
\end{center}

\courssub{Ressources à mobiliser}

\begin{aretenir}[Boîte à outils du résumé]
\begin{itemize}
\item \cle{Structure argumentative} : thèse $\rightarrow$ arguments $\rightarrow$ exemples $\rightarrow$ réfutation des objections $\rightarrow$ conclusion. Dans un résumé, on garde la thèse et les arguments ; on sacrifie les exemples.
\item \cle{Progression du texte} : progression thématique (un thème annoncé puis repris), progression par connecteurs (d'abord, ensuite, enfin), progression argumentative (de l'argument le plus simple au plus fort).
\item \cle{Cohésion et cohérence} : reprises pronominales, substituts lexicaux, champs lexicaux, connecteurs logiques garantissent un résumé lisible et logique.
\item \cle{Techniques de contraction} : suppression, substitution, fusion, généralisation. La règle d'or : un quart du texte, dans sa propre langue, sans commentaire.
\item \cle{Sigles} : s'ils apparaissent dans le texte, on les conserve tels quels (ex. MINESEC) ou on les explicite une fois.
\end{itemize}
\end{aretenir}

\courssub{Démarche et corrigé commenté}

\begin{exoresolu}[Résolution guidée de la tâche d'intégration]
Résumer l'éditorial « Faut-il rendre obligatoire l'enseignement des langues nationales au secondaire ? » en 100 mots ($\pm$ 10), selon les consignes 1 à 7.
\tcblower
\textbf{Étape 1 — Repérage de la structure argumentative.} La thèse est explicite dès l'introduction : « oui », l'enseignement des langues nationales doit être obligatoire au secondaire. Les connecteurs balisent trois arguments : \emph{d'abord} l'enracinement culturel (exemple : les collégiens incapables de saluer leurs grands-parents), \emph{ensuite} la facilitation des apprentissages (exemple : les écoles pilotes du Nord et de l'Est), \emph{enfin} la sauvegarde du patrimoine (exemple : les savoirs médicinaux). Le paragraphe « Certains objectent\ldots » réfute deux objections (coût, unité nationale). La conclusion appelle à l'engagement de l'État, des enseignants et des familles.

\textbf{Étape 2 — Plan du texte.} (1) Thèse : rendre obligatoire l'enseignement des langues nationales. (2) Argument culturel. (3) Argument pédagogique. (4) Argument patrimonial. (5) Réfutation des objections et conclusion.

\textbf{Étape 3 — Application des techniques.} \emph{Suppression} : on élimine les exemples (villages, écoles pilotes, savoirs médicinaux) et les énumérations de langues (ewondo, duala, fulfuldé, bassa). \emph{Généralisation} : « ewondo, duala, fulfuldé, bassa » devient « les langues nationales ». \emph{Substitution} : « une nation ne se construit pas sur des citoyens déracinés » devient « l'identité nationale en dépend ». \emph{Fusion} : les deux objections et leur réfutation se condensent en une phrase.

\textbf{Étape 4 — Résumé rédigé (modèle conforme).}
\begin{quote}
L'éditorialiste soutient que l'enseignement des langues nationales doit devenir obligatoire au secondaire. Il développe trois arguments : d'abord, cet enseignement assure l'enracinement culturel de la jeunesse, empêchant la rupture identitaire ; ensuite, il facilite les apprentissages scolaires, car l'enfant assimile mieux les concepts abstraits dans sa langue première, ce qui améliore par la suite sa maîtrise du français ; enfin, il préserve un patrimoine linguistique menacé, chaque langue disparaissant emportant des savoirs uniques. L'auteur réfute les objections financières et unitaires : le coût est un investissement et l'unité nationale se construit par le respect des identités. Il appelle l'État, les enseignants et les familles à s'engager résolument dans cette réforme. (110 mots)
\end{quote}

\textbf{Étape 5 — Comptage.} On compte chaque mot séparé par une espace ; les mots composés à trait d'union comptent pour un. On inscrit « (110 mots) » à la fin.

\textbf{Étape 6 — Correction croisée.} La grille permet de vérifier : fidélité (aucune opinion ajoutée), sélection (exemples absents), reformulation (aucune phrase copiée), cohésion (connecteurs : « il développe », « d'abord », « ensuite », « enfin », « L'auteur réfute », « Il appelle »), contrainte de longueur.

\textbf{Étape 7 — Justification orale.} Exemple : « J'ai généralisé la liste des quatre langues en "langues nationales" car la liste est un exemple ; j'ai fusionné les deux objections car elles relèvent du même mouvement de réfutation ; j'ai substitué "rupture linguistique est une rupture identitaire" par "empêchant la rupture identitaire". »
\end{exoresolu}

\begin{attention}[Pièges fréquents]
Ne donnez jamais votre avis dans un résumé ; ne copiez aucune phrase entière du texte ; ne conservez pas les exemples au détriment des arguments ; vérifiez toujours le comptage final, une pénalité est prévue hors de la fourchette.
\end{attention}

\courssub{Bilan de l'activité}

Cette tâche authentique a permis de mobiliser en une seule séquence l'ensemble des savoir-faire de la leçon. Le repérage des connecteurs a confirmé que la \cle{progression} d'un texte explicatif et argumentatif est balisée et prévisible, ce qui facilite la construction du plan. Le tableau de la structure argumentative a montré que thèse et arguments forment l'ossature du texte, tandis que exemples et énumérations en sont la chair, sacrifiable dans une contraction. Les quatre techniques (suppression, substitution, fusion, généralisation) se sont révélées complémentaires : aucune ne suffit seule à atteindre le quart du texte. Enfin, la correction croisée et la réécriture ont illustré qu'un bon résumé est toujours le fruit d'une \cle{confrontation} et d'une amélioration, et non d'un premier jet. Pour le Baccalauréat technique, les élèves s'entraîneront en conditions réelles : même tâche, en 45 minutes chronométrées, sans dictionnaire.

\newpage

\coursec{Méthodes et exercices résolus}

\courssub{Méthodes et exercices résolus}

\begin{methode}[Analyser un texte explicatif : progression, cohésion et cohérence]
\textbf{Objectif :} Identifier le type de progression, les marqueurs de cohésion et vérifier la cohérence globale pour comprendre la logique du texte avant de le résumer.
\begin{enumerate}
    \item \textbf{Lire globalement} pour saisir le thème, la thèse (si argumentatif) ou l'objet expliqué.
    \item \textbf{Repérer la progression} : 
    \begin{itemize}
        \item \emph{Chronologique} : indicateurs temporels (d'abord, ensuite, enfin, dates).
        \item \emph{Logique (déductive/inductive)} : connecteurs logiques (donc, car, par conséquent, or).
        \item \emph{Thématique / Énumérative} : mots de liaison d'ajout (de plus, ensuite, enfin, premièrement).
        \item \emph{Analogique / Comparative} : (de même, de façon similaire, au contraire).
    \end{itemize}
    \item \textbf{Analyser la cohésion} (liens linguistiques) : 
    \begin{itemize}
        \item \emph{Référence} : pronoms, déterminants, synonymes (reprise).
        \item \emph{Conjonction / Adverbes de liaison} : mais, cependant, ainsi que.
        \item \emph{Isotopies} : champs lexicaux récurrents.
    \end{itemize}
    \item \textbf{Vérifier la cohérence} (logique du sens) : le texte respecte-t-il les connaissances du monde ? Les énoncés sont-ils compatibles entre eux ?
    \item \textbf{Noter les sigles} : repérer les sigles (ex~: \cle{ONU}, \cle{CEMAC}, \cle{IA}) et leur développement à la première occurrence pour les conserver ou les expliquer dans le résumé.
\end{enumerate}
\cle{Réflexe Bac} : Un résumé fidèle respecte la progression de l'original. Ne jamais réorganiser les idées.
\end{methode}

\begin{exoresolu}[Analyse de la structure d'un texte explicatif]
\textbf{Énoncé :} Lisez l'extrait ci-dessous et répondez aux questions.
\begin{quote}
\textbf{Le développement durable au Cameroun}
\emph{Premièrement}, la gestion des forêts tropicales constitue un enjeu majeur. \emph{En effet}, le bassin du Congo, dont le Cameroun occupe une part significative, est le deuxième poumon vert de la planète. \emph{Par conséquent}, la déforestation incontrôlée menace la biodiversité. \emph{De plus}, l'exploitation minière, \emph{notamment} dans l'Est, dégrade les sols et pollue les cours d'eau. \emph{Cependant}, des initiatives comme le Plan National de Développement Durable (PNDD) visent à inverser la tendance. \emph{Ainsi}, la promotion de l'agroforesterie permet de concilier agriculture et préservation forestière. \emph{Enfin}, l'implication des populations locales reste la condition \emph{sine qua non} de la réussite.
\end{quote}
\begin{enumerate}
    \item Identifiez le type de progression de ce texte. Justifiez par deux indices linguistiques.
    \item Relevez trois procédés de cohésion (un par catégorie : référence, connecteur, isotopie).
    \item Quel est le sigle présent ? Donnez son développement probable.
\end{enumerate}
\tcblower
\textbf{Solution détaillée :}
\begin{enumerate}
    \item \textbf{Type de progression :} \textbf{Thématique / Énumérative (argumentative)}. 
    \textbf{Justifications :}
    \begin{itemize}
        \item L'usage d'organisateurs textuels énumératifs : \og \emph{Premièrement} \fg{}, \og \emph{De plus} \fg{}, \og \emph{Enfin} \fg{} qui structurent le discours en points distincts s'ajoutant les uns aux autres.
        \item L'absence de chronologie stricte (pas de dates, pas de succession d'actions dans le temps) mais une addition d'arguments (forêts, mines, solutions, condition de réussite).
    \end{itemize}
    \item \textbf{Procédés de cohésion :}
    \begin{itemize}
        \item \emph{Référence (anaphore) :} \og \emph{dont} \fg{} (relatif, antécédent : \og le bassin du Congo \fg{}) ; \og \emph{la tendance} \fg{} (reprise nominale de \og inverser la tendance \fg{} / déforestation).
        \item \emph{Connecteur logique :} \og \emph{Par conséquent} \fg{} (conséquence), \og \emph{Cependant} \fg{} (opposition/concession), \og \emph{Ainsi} \fg{} (conséquence/résultat).
        \item \emph{Isotopie (champ lexical) :} Champ lexical de l'environnement / écologie : \og forêts tropicales, bassin du Congo, poumon vert, biodiversité, déforestation, sols, cours d'eau, agroforesterie, préservation forestière \fg{}.
    \end{itemize}
    \item \textbf{Sigle :} \cle{PNDD}. \textbf{Développement probable :} \textbf{Plan National de Développement Durable}. (Au Cameroun, c'est le document de référence). Dans un résumé, on l'écrirait \og Plan National de Développement Durable (PNDD) \fg{} à la première occurrence.
\end{enumerate}
\textbf{Conclusion :} Ce texte suit une progression thématique énumérative, cohérente par son champ lexical environnemental et cohésive par ses connecteurs logiques et ses reprises référentielles.
\end{exoresolu}

\begin{methode}[Identifier la structure argumentative d'un texte]
\textbf{Objectif :} Découper le texte en mouvements logiques (thèse, arguments, exemples, conclusion) pour construire le plan du résumé.
\begin{enumerate}
    \item \textbf{Délimiter les paragraphes} (ou séquences si texte continu) : un paragraphe = une idée force.
    \item \textbf{Trouver la thèse} (souvent en début ou fin de texte, parfois implicite). Question : \og Que veut prouver l'auteur ? \fg{}
    \item \textbf{Isoler les arguments} : idées principales qui soutiennent la thèse. Cherchez les connecteurs : \og car \fg{}, \og parce que \fg{}, \og puisque \fg{}, \og en effet \fg{}.
    \item \textbf{Repérer les illustrations/exemples} : \og par exemple \fg{}, \og c'est le cas de \fg{}, \og notamment \fg{}. \textbf{Attention :} Dans une contraction (résumé), les exemples sont \textbf{supprimés} ou \textbf{très condensés} (seule la valeur argumentative compte).
    \item \textbf{Identifier les concessions / contre-arguments} : \og certes \fg{}, \og il est vrai que \fg{}, \og cependant \fg{}. Ils nuancent la thèse.
    \item \textbf{Formuler la conclusion} : reprise de la thèse, ouverture, appel à l'action.
    \item \textbf{Construire le plan détaillé} : Thèse $\rightarrow$ Argument 1 (+ illustration condensée) $\rightarrow$ Argument 2... $\rightarrow$ Concession $\rightarrow$ Conclusion.
\end{enumerate}
\cle{Règle d'or} : Le plan du résumé \textbf{calque} le plan du texte original. On ne réorganise pas.
\end{methode}

\begin{exoresolu}[Structuration argumentative et plan de résumé]
\textbf{Énoncé :} Étudiez la structure argumentative du texte suivant pour en proposer un plan de résumé.
\begin{quote}
\textbf{L'école face au numérique}
\og Certes, l'introduction des tablettes en classe modernise les pratiques pédagogiques. \emph{Cependant}, elle ne saurait remplacer l'enseignant. \emph{D'abord}, la relation humaine reste le cœur de l'apprentissage : l'empathie, l'autorité bienveillante, l'adaptation immédiate à l'élève en difficulté sont irremplaçables. \emph{Ensuite}, l'outil numérique, sans guidage, devient source de distraction ; \emph{par exemple}, les réseaux sociaux détournent l'attention des apprenants. \emph{Par ailleurs}, la fracture numérique exclut les élèves défavorisés ne possédant pas de connexion à domicile. \emph{En définitive}, le numérique doit rester un \emph{outil au service de la pédagogie}, et non une fin en soi. \fg
\end{quote}
\begin{enumerate}
    \item Quelle est la thèse défendue ? Est-elle explicite ou implicite ?
    \item Identifiez la concession d'entrée. Quel en est le rôle ?
    \item Listez les arguments principaux (formulés en nominale).
    \item Relevez l'exemple. Doit-il figurer intégralement dans le résumé ?
    \item Proposez un plan détaillé pour un résumé de ce texte (4 à 5 mouvements).
\end{enumerate}
\tcblower
\textbf{Solution détaillée :}
\begin{enumerate}
    \item \textbf{Thèse :} \og Le numérique doit rester un outil au service de la pédagogie et non remplacer l'enseignant. \fg{} Elle est \textbf{explicite} (dernière phrase : \og En définitive... \fg{}).
    \item \textbf{Concession :} \og Certes, l'introduction des tablettes en classe modernise les pratiques pédagogiques. \fg{} \textbf{Rôle :} Admettre un point positif pour mieux faire accepter la thèse principale (stratégie de l'aveu partiel / \emph{concessio}).
    \item \textbf{Arguments (formulation nominale pour le plan) :}
    \begin{enumerate}
        \item La primauté de la relation humaine (empathie, autorité, adaptation).
        \item Le risque de distraction sans encadrement (réseaux sociaux).
        \item La fracture numérique (exclusion des élèves défavorisés).
    \end{enumerate}
    \item \textbf{Exemple :} \og par exemple, les réseaux sociaux détournent l'attention \fg{}. \textbf{Traitement au résumé :} \textbf{Non}, il ne figure pas intégralement. On le condense : \og risque de distraction (réseaux sociaux) \fg{} ou on le supprime si la contrainte de longueur est forte, en gardant l'argument \og distraction \fg{}.
    \item \textbf{Plan détaillé du résumé :}
    \begin{enumerate}
        \item \textbf{Concession / Thèse introductive :} L'apport des tablettes (modernisation) ne doit pas masquer l'irremplaçabilité de l'enseignant.
        \item \textbf{Argument 1 :} La relation humaine (empathie, adaptation) est le cœur de l'apprentissage.
        \item \textbf{Argument 2 :} L'outil numérique non guidé génère de la distraction.
        \item \textbf{Argument 3 :} La fracture numérique aggrave les inégalités sociales.
        \item \textbf{Conclusion / Thèse finale :} Le numérique : outil au service de la pédagogie, non fin en soi.
    \end{enumerate}
\end{enumerate}
\textbf{Conclusion :} La structure est classique : Concession $\rightarrow$ Thèse (différée) $\rightarrow$ Arguments $\rightarrow$ Conclusion réaffirmant la thèse. Le plan du résumé suit cet ordre exact.
\end{exoresolu}

\begin{methode}[Techniques de rédaction du résumé (Contraction) : Sélection, Reformulation, Condensation]
\textbf{Objectif :} Produire un texte court (généralement 1/4 à 1/3 de l'original), fidèle, autonome et bien écrit.
\begin{enumerate}
    \item \textbf{Sélection (Le \og quoi \fg{}) :} 
    \begin{itemize}
        \item \textbf{Garder :} Thèse, arguments principaux, articulations logiques (connecteurs), données chiffrées clés, noms propres essentiels, sigles développés.
        \item \textbf{Supprimer :} Exemples détaillés, anecdotes, répétitions, métaphores, adjectifs qualificatifs subjectifs, commentaires de l'auteur (\og je pense que \fg{}), incises.
    \end{itemize}
    \item \textbf{Reformulation / Transposition (Le \og comment \fg{}) :} 
    \begin{itemize}
        \item \emph{Discours direct $\rightarrow$ Indirect} : \og Il dit : "Je viens" \fg{} $\rightarrow$ \og Il dit qu'il vient \fg{}.
        \item \emph{Verbe $\rightarrow$ Nom (nominalisation) :} \og Il a décidé de partir \fg{} $\rightarrow$ \og Sa décision de partir \fg{} / \og Son départ \fg{}. \textbf{Gain de place majeur.}
        \item \emph{Subordination $\rightarrow$ Préposition / Participe :} \og Parce qu'il pleut, il rentre \fg{} $\rightarrow$ \og Du fait de la pluie, il rentre \fg{} / \og Il pleuvant, il rentre \fg{}.
        \item \emph{Remplacement lexical :} Synonymes plus précis/génériques, hyperonymes (\og chien, chat \fg{} $\rightarrow$ \og animaux \fg{}).
        \item \emph{Suppression des mots-outils :} \og Le fait que \fg{} $\rightarrow$ \og Que \fg{} ; \og Au niveau de \fg{} $\rightarrow$ \og À \fg{} / \og Dans \fg{}.
    \end{itemize}
    \item \textbf{Condensation (Fusion) :} Regrouper plusieurs phrases en une seule phrase complexe.
    \item \textbf{Vérification finale :} 
    \begin{itemize}
        \item Respect de la progression originale ?
        \item Neutralité (pas d'opinion personnelle) ?
        \item Exactitude de l'information ?
        \item Respect de la longueur imposée ($\pm 10\%$) ?
        \item Correction linguistique (orthographe, syntaxe, ponctuation) ?
    \end{itemize}
\end{enumerate}
\cle{Formule magique} : \textbf{Nominaliser + Subordonner + Généraliser = Contracter.}
\end{methode}

\begin{exoresolu}[Application des techniques de contraction]
\textbf{Énoncé :} Contractez le paragraphe suivant en \textbf{une seule phrase} de 30 mots maximum, en conservant l'information essentielle.
\begin{quote}
\og Le directeur de l'école a convoqué les parents d'élèves. Il leur a expliqué que les résultats du premier trimestre étaient catastrophiques. Il a ajouté que l'absentéisme était la cause principale de cet échec. Il a proposé des cours de rattrapage pendant les vacances. \fg
\end{quote}
\tcblower
\textbf{Solution détaillée :}
\textbf{Analyse de l'original (56 mots) :}
\begin{itemize}
    \item Sujets : Directeur / Parents / Résultats / Absentéisme / Cours.
    \item Actions : Convoquer, Expliquer (résultats catastrophiques), Ajouter (cause = absentéisme), Proposer (rattrapage).
\end{itemize}
\textbf{Étapes de contraction :}
\begin{enumerate}
    \item \textbf{Sélection :} Garder : Directeur, Parents, Résultats catastrophiques, Cause (absentéisme), Solution (rattrapage vacances). Supprimer : \og premier trimestre \fg{} (implicite), \og cause principale \fg{} $\rightarrow$ \og cause \fg{}.
    \item \textbf{Nominalisation :} 
    \begin{itemize}
        \item \og Le directeur a convoqué \fg{} $\rightarrow$ \og La convocation des parents par le directeur \fg{} (ou \og Le directeur convoque \fg{}).
        \item \og Il a expliqué que les résultats étaient catastrophiques \fg{} $\rightarrow$ \og L'explication des résultats catastrophiques \fg{}.
        \item \og Il a ajouté que l'absentéisme était la cause \fg{} $\rightarrow$ \og L'imputation de l'échec à l'absentéisme \fg{}.
        \item \og Il a proposé des cours \fg{} $\rightarrow$ \og La proposition de cours de rattrapage \fg{}.
    \end{itemize}
    \item \textbf{Fusion / Subordination :} Relier par des prépositions (\og pour \fg{}, \og afin de \fg{}, \og en imputant \fg{}).
\end{enumerate}
\textbf{Propositions de phrase unique (environ 25 mots) :}
\begin{itemize}
    \item \og Le directeur convoque les parents pour expliquer l'échec catastrophique, imputable à l'absentéisme, et propose des cours de rattrapage pendant les vacances. \fg{} (26 mots) \checkmark
    \item \og Convoquant les parents, le directeur impute les résultats catastrophiques à l'absentéisme et propose un rattrapage pendant les vacances. \fg{} (21 mots) \checkmark
\end{itemize}
\textbf{Vérification :} Fidèle ? Oui. Progression respectée ? Oui. Longueur ? Oui. Neutre ? Oui.
\end{exoresolu}

\begin{methode}[Rédiger le résumé : Du plan à la rédaction finale (épreuve type Bac)]
\textbf{Objectif :} Passer du plan détaillé au texte continu, fluide et cohérent, dans la limite de mots imposée.
\begin{enumerate}
    \item \textbf{Rédiger le brouillon (étape obligatoire) :} Écrire \og au fil de l'eau \fg{} en suivant le plan, sans compter les mots précisément, en appliquant les techniques (nominalisation, discours indirect).
    \item \textbf{Compter les mots :} Compter \textbf{tous les mots} (mots grammaticaux inclus : \og le, de, à \fg{} comptent pour 1). Noter le total en marge.
    \item \textbf{Calibrer (Ajuster la longueur) :}
    \begin{itemize}
        \item \emph{Trop long :} Couper les détails résiduels, fusionner davantage les phrases, nominaliser plus radicalement, supprimer les connecteurs lourds (\og de plus \fg{} $\rightarrow$ \og ; \fg{}).
        \item \emph{Trop court :} Développer légèrement un argument clé, réintégrer un connecteur logique important, préciser un chiffre ou un nom propre omis.
    \end{itemize}
    \item \textbf{Soigner les articulations (Cohésion) :} Vérifier les connecteurs logiques (\og Ainsi \fg{}, \og Toutefois \fg{}, \og En outre \fg{}), les reprises pronominales (\og ce dernier \fg{}, \og cette mesure \fg{}), la ponctuation.
    \item \textbf{Relire à voix basse (ou mentalement) :} Traquer les fautes d'accord, les anglicismes, les répétitions, les phrases sans verbe (sauf choix stylistique assumé).
    \item \textbf{Recopier au propre :} Écriture lisible, paragraphes respectés (si texte long) ou bloc unique (si court). \textbf{Indiquer le nombre de mots exact à la fin} (ex~: \og [123~mots] \fg{}).
\end{enumerate}
\cle{Consigne Bac} : \og Vous rédigerez un résumé de \textbf{X mots} ($\pm 10\%$). Tout excédent ou déficit significatif sera pénalisé. \fg{}
\end{methode}

\begin{exoresolu}[Rédaction complète d'un résumé calibré]
\textbf{Énoncé :} Résumez le texte suivant en \textbf{80~mots $\pm 10\%$} (soit 72 à 88 mots).
\begin{quote}
\textbf{La jeunesse camerounaise et l'entrepreneuriat}
\og Au Cameroun, le chômage des jeunes diplômés atteint des proportions alarmantes. Selon l'Institut National de la Statistique (INS), le taux de sous-emploi des 15-35 ans dépasse 70~\%. Face à cette situation, l'État encourage l'auto-emploi via le Programme d'Appui à l'Insertion Professionnelle (PAIP). Ce programme offre des formations techniques et un financement sous forme de microcrédits. Toutefois, l'accès au financement reste difficile : les banques exigent des garanties que les jeunes n'ont pas. De plus, la fiscalité pesante étouffe les jeunes pousses. Pour réussir, il faut simplifier l'accès au crédit et alléger la fiscalité des start-ups. \fg
\end{quote}
\tcblower
\textbf{Solution détaillée :}
\textbf{Étape 1 : Analyse \& Plan (Mental / Brouillon)}
\begin{itemize}
    \item \textbf{Thème :} Jeunesse / Chômage / Entrepreneuriat au Cameroun.
    \item \textbf{Progression :} Constat (Chômage/INS) $\rightarrow$ Réponse État (PAIP : formation + microcrédit) $\rightarrow$ Obstacles (Garanties bancaires, Fiscalité) $\rightarrow$ Solutions (Simplifier crédit, Alléger fiscalité).
    \item \textbf{Sigles :} INS (Institut National de la Statistique), PAIP (Programme d'Appui à l'Insertion Professionnelle).
\end{itemize}
\textbf{Étape 2 : Rédaction Brouillon (Application techniques)}
\og Au Cameroun, le chômage des jeunes diplômés est alarmant (INS : sous-emploi > 70~\% chez les 15-35 ans). L'État réagit par le PAIP (formations, microcrédits). Mais l'accès au crédit bute sur l'exigence de garanties bancaires inaccessibles, et la fiscalité étouffe les start-ups. La réussite impose de simplifier l'accès au financement et d'alléger la fiscalité. \fg
\textbf{Étape 3 : Comptage \& Calibrage}
Comptage : \og Au (1) Cameroun, (2) le (3) chômage (4) des (5) jeunes (6) diplômés (7) est (8) alarmant (9) (INS (10) : (11) sous-emploi (12) > (13) 70~\% (14) chez (15) les (16) 15-35 (17) ans). (18) L'État (19) réagit (20) par (21) le (22) PAIP (23) (formations, (24) microcrédits). (25) Mais (26) l'accès (27) au (28) crédit (29) bute (30) sur (31) l'exigence (32) de (33) garanties (34) bancaires (35) inaccessibles, (36) et (37) la (38) fiscalité (39) étouffe (40) les (41) start-ups. (42) La (43) réussite (44) impose (45) de (46) simplifier (47) l'accès (48) au (49) financement (50) et (51) d'alléger (52) la (53) fiscalité. (54) \fg{} $\rightarrow$ \textbf{54 mots}.
\textbf{Problème :} Trop court (min 72 mots). Il faut développer.
\textbf{Étape 4 : Développement contrôlé (Réintégrer précision + connecteurs)}
\og Au Cameroun, le chômage des jeunes diplômés atteint des proportions alarmantes : selon l'INS, le sous-emploi dépasse 70~\% chez les 15-35 ans. Pour y remédier, l'État a mis en place le PAIP, offrant formations techniques et microcrédits. Toutefois, l'efficacité de ce dispositif est freinée par l'exigence de garanties bancaires inaccessibles aux jeunes et par une fiscalité lourde qui étouffe les start-ups. Ainsi, la réussite de l'insertion professionnelle passe impérativement par la simplification de l'accès au crédit et l'allègement de la fiscalité. \fg
\textbf{Étape 5 : Comptage Final}
Au(1) Cameroun,(2) le(3) chômage(4) des(5) jeunes(6) diplômés(7) atteint(8) des(9) proportions(10) alarmantes(11) :(12) selon(13) l'INS,(14) le(15) sous-emploi(16) dépasse(17) 70~\%(18) chez(19) les(20) 15-35(21) ans.(22) Pour(23) y(24) remédier,(25) l'État(26) a(27) mis(28) en(29) place(30) le(31) PAIP,(32) offrant(33) formations(34) techniques(35) et(36) microcrédits.(37) Toutefois,(38) l'efficacité(39) de(40) ce(41) dispositif(42) est(43) freinée(44) par(45) l'exigence(46) de(47) garanties(48) bancaires(49) inaccessibles(50) aux(51) jeunes(52) et(53) par(54) une(55) fiscalité(56) lourde(57) qui(58) étouffe(59) les(60) start-ups.(61) Ainsi,(62) la(63) réussite(64) de(65) l'insertion(66) professionnelle(67) passe(68) impérativement(69) par(70) la(71) simplification(72) de(73) l'accès(74) au(75) crédit(76) et(77) l'allègement(78) de(79) la(80) fiscalité.(81)
\textbf{Total : 81 mots.} \checkmark (Dans l'intervalle 72-88).
\textbf{Version finale :} (Recopie au propre avec indication \og [81~mots] \fg{}).
\end{exoresolu}

\begin{methode}[Étudier "Le Vieux Nègre et la Médaille" : Contexte, Tonalités (Satirique/Polémique) et Bilan]
\textbf{Objectif :} Maîtriser les connaissances requises sur l'œuvre pour l'oral ou l'écrit (dissertation, commentaire, résumé d'extrait).
\begin{enumerate}
    \item \textbf{Contexte historique \& biographique :} 
    \begin{itemize}
        \item \textbf{Auteur :} Ferdinand \cle{Oyono} (Cameroun, 1929-2010). Diplomate, haut fonctionnaire.
        \item \textbf{Date :} 1956 (période coloniale, veille des indépendances).
        \item \textbf{Genre :} Roman court / Nouvelle / Récit initiatique à rebours.
        \item \textbf{Thème central :} La désillusion coloniale, l'aliénation, la vanité des décorations coloniales, le choc des cultures.
    \end{itemize}
    \item \textbf{Intrigue (Résumé synthétique) :} Meka, vieux chef respecté, reçoit la médaille de la \cle{Légion d'honneur}. Il prépare une grande fête. Le jour J, l'administrateur ne vient pas ; un sous-officier remet la médaille vite fait. Meka réalise l'humiliation : la médaille ne vaut rien, il a été dupé par le système colonial qu'il a servi. Il meurt peu après, lucide.
    \item \textbf{Tonalité \cle{Satirique} :} 
    \begin{itemize}
        \item \emph{Cible :} L'administration coloniale (absurde, méprisante), les notables complices (Meka d'abord), les Blancs (stéréotypes : l'administrateur paresseux, le prêtre hypocrite).
        \item \emph{Procédés :} Ironie (gap entre attente Meka et réalité), hyperboles (préparatifs grandioses pour rien), sous-entendus, portrait caricatural (le commandant \og gros \fg{}, \og rougeaud \fg{}).
    \end{itemize}
    \item \textbf{Tonalité \cle{Polémique} :} 
    \begin{itemize}
        \item \emph{Charge directe :} Dénonciation du racisme, de l'exploitation (travail forcé, impôts), de la \og mission civilisatrice \fg{} mensongère.
        \item \emph{Portée :} Politique. Le texte devient un réquisitoire contre le colonialisme. La médaille = symbole de l'aliénation.
    \end{itemize}
    \item \textbf{Bilan / Évaluation de l'œuvre :} 
    \begin{itemize}
        \item \emph{Modernité :} Langue simple, orale, humour noir (précurseur).
        \item \emph{Universalisme :} Au-delà du Cameroun, critique de toute domination.
        \item \emph{Limite :} Vision très masculine (femmes en arrière-plan), fin tragique sans espoir de révolte collective explicite (mais lucidité individuelle).
    \end{itemize}
\end{enumerate}
\cle{Point méthode} : Pour un résumé d'extrait de cette œuvre : garder la progression narrative (Attente $\rightarrow$ Cérémonie $\rightarrow$ Déception $\rightarrow$ Lucidité/Mort), neutraliser l'ironie (la restituer par le vocabulaire : \og prétendue cérémonie \fg{}, \og remise expéditive \fg{}).
\end{methode}

\begin{exoresolu}[Analyse tonalités et résumé d'extrait "Le Vieux Nègre et la Médaille"]
\textbf{Énoncé :} Lisez l'extrait (fin du roman, mort de Meka). Répondez aux questions.
\begin{quote}
\og ... Meka regarda la médaille. Elle ne brillait plus. C'était un petit morceau de métal sans valeur, comme ceux qu'on donne aux enfants pour jouer. Il pensa à ses fils, à ses petits-fils, à sa tribu. Il avait tout donné : ses terres, sa paix, son honneur. Pour ça ? Pour un bout de ferraille ? La nuit tombait sur la case. Meka ferma les yeux. Il ne rouvrit plus. \fg
\end{quote}
\begin{enumerate}
    \item Identifiez deux procédés relevant de la \textbf{tonalité satirique} et deux de la \textbf{tonalité polémique} dans cet extrait.
    \item Rédigez un résumé de cet extrait en \textbf{30~mots $\pm 10\%$}.
\end{enumerate}
\tcblower
\textbf{Solution détaillée :}
\begin{enumerate}
    \item \textbf{Tonalité Satirique :}
    \begin{itemize}
        \item \emph{Dépréciation / Ironie :} \og Elle ne brillait plus. C'était un petit morceau de métal sans valeur, comme ceux qu'on donne aux enfants pour jouer. \fg{} (Réduction du symbole glorieux à un jouet d'enfant).
        \item \emph{Contraste grotesque :} \og Il avait tout donné... Pour ça ? Pour un bout de ferraille ? \fg{} (Disproportion massive entre le sacrifice et la récompense).
    \end{itemize}
    \textbf{Tonalité Polémique :}
    \begin{itemize}
        \item \emph{Accusation directe / Réquisitoire :} \og Il avait tout donné : ses terres, sa paix, son honneur. \fg{} (Liste des spoliations subies par le colonisé au profit du colonisateur).
        \item \emph{Symbolique de l'aliénation :} La médaille (\og bout de ferraille \fg{}) incarne la duperie du système colonial qui achète la soumission pour une pacotille.
    \end{itemize}
    \item \textbf{Résumé (30~mots $\pm 10\%$ = 27 à 33 mots) :}
    \textbf{Brouillon :} Meka, lucide avant de mourir, réalise que la médaille reçue ne vaut rien face aux sacrifices consentis (terres, honneur). Il meurt désabusé, victime de la duperie coloniale.
    \textbf{Comptage :} Meka(1), lucide(2), avant(3), de(4), mourir,(5) réalise(6), que(7), la(8), médaille(9), reçue(10), ne(11), vaut(12), rien(13), face(14), aux(15), sacrifices(16), consentis(17), (terres,(18), honneur).(19) Il(20), meurt(21), désabusé,(22), victime(23), de(24), la(25), duperie(26), coloniale(27). $\rightarrow$ \textbf{27 mots}. \checkmark (Limite basse atteinte).
    \textbf{Version finale :} \og Meka, lucide avant de mourir, réalise la nullité de sa médaille face aux sacrifices consentis (terres, honneur). Il s'éteint, victime de la duperie coloniale. \fg{} [27~mots]
\end{enumerate}
\end{exoresolu}

\begin{methode}[Confronter et améliorer les productions (Auto-correction / Co-évaluation)]
\textbf{Objectif :} Développer l'esprit critique sur sa propre production ou celle d'un pair pour gagner en qualité (fidélité, concision, correction).
\begin{enumerate}
    \item \textbf{Échange / Lecture croisée :} Lire le résumé du camarade (ou relire le sien après une pause).
    \item \textbf{Grille d'évaluation critique (4 critères) :}
    \begin{enumerate}
        \item \textbf{Fidélité / Contenu :} Thèse présente ? Arguments principaux tous là ? Pas d'invention ? Pas d'omission majeure ? Progression respectée ?
        \item \textbf{Concision / Contraction :} Longueur respectée ? Exemples supprimés ? Nominalisations utilisées ? Pas de répétitions ?
        \item \textbf{Cohérence / Cohésion :} Connecteurs logiques justes ? Références claires (pronoms) ? Texte compréhensible sans l'original ?
        \item \textbf{Correction linguistique :} Orthographe (accords, lexique), Syntaxe (phrases complètes, pas d'anacoluthes), Ponctuation, Registre (soutenu/neutre).
    \end{enumerate}
    \item \textbf{Annotation :} Surligner en \textcolor{popgreen}{vert} (bon), \textcolor{poporange}{orange} (à améliorer), \textcolor{poppink}{rose} (erreur). Écrire une \og suggestion précise \fg{} en marge (ex~: \og Nominalisez ici \fg{}, \og Connecteur logique manquant \fg{}).
    \item \textbf{Réécriture ciblée :} Ne réécrire que les passages signalés. Ne pas tout recopier si le reste est bon.
    \item \textbf{Validation finale :} Recompter les mots. Relire à voix haute.
\end{enumerate}
\cle{Recommandation} : Accepter la critique comme un outil de progression, pas comme une attaque. La confrontation est l'étape où l'on gagne le plus de points.
\end{methode}

\begin{exoresolu}[Analyse et correction d'une production d'élève]
\textbf{Énoncé :} Voici une proposition de résumé (consigne : 80~mots $\pm 10\%$) pour le texte sur \og La jeunesse camerounaise et l'entrepreneuriat \fg{} (voir exercice résolu n°4). Identifiez 3 défauts majeurs et proposez une version corrigée.
\begin{quote}
\og \textbf{Proposition élève :} Dans ce texte, l'auteur nous dit que au Cameroun les jeunes diplômés sont au chômage. L'INS dit que c'est plus de 70 pour cent. L'État a fait le PAIP pour aider avec des formations et de l'argent. Mais les banques demandent des garanties que les jeunes n'ont pas. Aussi les impôts sont trop chers pour les jeunes entreprises. Donc il faut changer ça pour que ça marche. Je pense que c'est une bonne idée. \fg
\end{quote}
\tcblower
\textbf{Solution détaillée :}
\textbf{Analyse (Grille d'évaluation) :}
\begin{enumerate}
    \item \textbf{Fidélité / Neutralité (Échec) :} Présence de \og Dans ce texte, l'auteur nous dit que \fg{} (métadiscours inutile), \og Je pense que c'est une bonne idée \fg{} (\textbf{Opinion personnelle interdite}).
    \item \textbf{Concision / Style (Échec) :} Phrases courtes, juxtaposées, style oral (\og c'est plus de 70 pour cent \fg{} au lieu de \og dépasse 70~\% \fg{}, \og l'argent \fg{} au lieu de \og microcrédits \fg{}, \og les impôts \fg{} au lieu de \og la fiscalité \fg{}). Longueur : ~65 mots (trop court, manque de densité).
    \item \textbf{Cohésion / Correction (Faible) :} \og que au \fg{} (hiatus, fautif : \og qu'au \fg{}), \og Aussi \fg{} mal employé (conséquence vs addition), ponctuation faible.
    \item \textbf{Contenu :} Sigles INS/PAIP non développés à la 1ère occurrence (ou mal intégrés). Argument \og terres/paix/honneur \fg{} absent (normal, pas dans l'extrait fourni ici, mais l'argument \og fiscalité \fg{} est là).
\end{enumerate}
\textbf{3 Défauts majeurs à corriger :}
\begin{enumerate}
    \item Supprimer toute trace d'énonciation subjective (\og l'auteur dit \fg{}, \og je pense \fg{}).
    \item Nominaliser et préciser le vocabulaire (microcrédits, fiscalité, sous-emploi).
    \item Améliorer la cohésion (connecteurs : \og Toutefois \fg{}, \og Ainsi \fg{}) et corriger \og qu'au \fg{}.
\end{enumerate}
\textbf{Version corrigée (calibrée ~80 mots) :}
\og Au Cameroun, le sous-emploi des jeunes diplômés (15-35 ans) dépasse 70~\% selon l'INS. L'État y répond par le PAIP (formations, microcrédits). Toutefois, l'exigence de garanties bancaires bloque l'accès au financement, tandis qu'une fiscalité lourde étouffe les start-ups. La réussite de l'insertion professionnelle exige donc la simplification de l'accès au crédit et l'allègement de la fiscalité. \fg{} [78~mots] \checkmark
\end{exoresolu}

\begin{methode}[Gérer les sigles et abréviations dans le résumé]
\textbf{Objectif :} Traiter correctement les sigles (acronymes, initialismes) pour la clarté et le gain de place.
\begin{enumerate}
    \item \textbf{Repérage :} Identifier tous les sigles dans le texte source (ex~: \cle{ONU}, \cle{PIB}, \cle{CEMAC}, \cle{IA}, \cle{ONG}).
    \item \textbf{Règle de la première occurrence :} Au premier emploi dans \textbf{votre} résumé, écrire le \textbf{nom complet} suivi du sigle entre parenthèses. Ex~: \og Organisation des Nations Unies (ONU) \fg{}.
    \item \textbf{Emplois suivants :} Utiliser \textbf{uniquement le sigle}.
    \item \textbf{Cas des sigles "transparents" (notoriété publique) :} Si le sigle est plus connu que le nom complet (ex~: \cle{TV}, \cle{radio}, \cle{USB}, \cle{GPS}, \cle{VIH}, \cle{ONU} souvent), on \textbf{peut} l'utiliser seul dès le début, mais le développer est plus sûr au Bac.
    \item \textbf{Contraction par le sigle :} Remplacer une longue appellation répétée par son sigle \textbf{dès la 2ème occurrence} dans le résumé. C'est un outil de condensation puissant.
    \item \textbf{Attention :} Ne pas inventer de sigles. Ne pas mettre de points entre les lettres (ex~: \og O.N.U. \fg{} $\rightarrow$ \og ONU \fg{}).
\end{enumerate}
\cle{Astuce} : Si le texte source utilise un sigle sans le développer, et que vous le connaissez, développez-le dans le résumé (preuve de culture). Si vous ne le connaissez pas, gardez le sigle tel quel.
\end{methode}

\begin{exoresolu}[Traitement des sigles en situation de contraction]
\textbf{Énoncé :} Contractez le paragraphe suivant en 40~mots max. Gérez rigoureusement les sigles.
\begin{quote}
\og L'Organisation Mondiale de la Santé (OMS) a publié un rapport alarmant. L'OMS y dénonce la résistance aux antibiotiques. Cette résistance, selon l'OMS, menace les acquis de la médecine moderne. L'Organisation Mondiale de la Santé appelle à une mobilisation mondiale. Les États membres de l'OMS doivent agir vite. \fg
\end{quote}
\tcblower
\textbf{Solution détaillée :}
\textbf{Analyse :} Texte très répétitif (5 occurrences de l'OMS). Idéal pour la contraction par sigle.
\textbf{Stratégie :} 
\begin{enumerate}
    \item 1ère occurrence dans mon résumé : \og Organisation Mondiale de la Santé (OMS) \fg{}.
    \item Occurrences suivantes : \og OMS \fg{}.
    \item Fusionner les phrases : Rapport $\rightarrow$ Dénonciation $\rightarrow$ Menace $\rightarrow$ Appel $\rightarrow$ Action États.
\end{enumerate}
\textbf{Rédaction :}
\og L'Organisation Mondiale de la Santé (OMS) publie un rapport alarmant sur la résistance aux antibiotiques, qui menace la médecine moderne. L'OMS appelle à une mobilisation mondiale et exige une action rapide de ses États membres. \fg
\textbf{Comptage :} L'(1) Organisation(2) Mondiale(3) de(4) la(5) Santé(6) (OMS)(7) publie(8) un(9) rapport(10) alarmant(11) sur(12) la(13) résistance(14) aux(15) antibiotiques,(16) qui(17) menace(18) la(19) médecine(20) moderne.(21) L'OMS(22) appelle(23) à(24) une(25) mobilisation(26) mondiale(27) et(28) exige(29) une(30) action(31) rapide(32) de(33) ses(34) États(35) membres(36). $\rightarrow$ \textbf{36 mots}. \checkmark (Dans les 40~mots max).
\textbf{Gain :} 36 mots au lieu de ~65 mots original. Sigle géré normativement.
\end{exoresolu}

\begin{attention}[Les 5 erreurs qui coûtent des points au Bac (Contraction de texte)]
\begin{enumerate}
    \item \textbf{Le "Copier-Coller" déguisé (Plagiat / Absence de reformulation) :} Recopier des phrases entières du texte original en changeant juste 2-3 mots. \textbf{Sanction :} Note très faible (0 à 4/20~souvent), car l'épreuve évalue la \emph{reformulation} et la \emph{synthèse}. \textbf{Parade :} Systématiser la nominalisation et le discours indirect sur le brouillon.
    \item \textbf{L'apport d'idées personnelles / Jugement de valeur :} Écrire \og L'auteur a raison de dire que... \fg{}, \og C'est une bonne initiative... \fg{}, \og Heureusement, l'État a réagi... \fg{}. \textbf{Sanction :} Pénalité lourde sur le critère "Fidélité / Neutralité". Le résumé n'est pas un commentaire. \textbf{Parade :} Bannir \og je \fg{}, \og nous \fg{}, \og à mon avis \fg{}, adjectifs subjectifs (\og malheureusement \fg{}, \og heureusement \fg{}).
    \item \textbf{Le non-respect de la longueur imposée (Hors fourchette $\pm 10\%$) :} Rendre 50 mots pour 80 demandés, ou 120 mots. \textbf{Sanction :} Pénalité proportionnelle (souvent 1 à 2~points en moins sur la note globale, voire note plafonnée). \textbf{Parade :} Compter \textbf{deux fois} (brouillon + propre). Ajuster par nominalisation (si trop long) ou réintégration de connecteurs/précisions (si trop court).
    \item \textbf{La rupture de la progression / Désorganisation du plan :} Inverser l'ordre des arguments (mettre la conclusion au début), mélanger thèse et concession, regrouper des idées dispersées dans l'original. \textbf{Sanction :} Perte de points sur "Structure / Cohérence". Le correcteur suit le texte original en parallèle. \textbf{Parade :} Faire le plan détaillé \textbf{avant} d'écrire. Numéroter les paragraphes du texte source. Suivre l'ordre 1, 2, 3...
    \item \textbf{La négligence de la correction linguistique (Accords, Syntaxe, Sigles) :} Fautes d'accords (participes passés, sujets-verbes), phrases sans verbe (anacoluthes), sigles non développés à la 1ère occurrence, hiatus (\og le arbre \fg{}), anglicismes (\og opportunité \fg{} pour \og occasion \fg{}). \textbf{Sanction :} Pénalité sur la note de langue (souvent 2 à 4~points sur 20). \textbf{Parade :} Relire \textbf{spécifiquement pour la langue} (dernière relecture), à voix basse, en pointant chaque mot du doigt.
\end{enumerate}
\end{attention}

\newpage

\coursec{Exercices}

\coursec{Contraction de texte : rédiger un résumé de texte}

\courssub{Je vérifie mes connaissances}

\begin{aretenir}[Rappels essentiels : le résumé de texte]
Un \cle{résumé} est une reformulation brève, fidèle et personnelle des idées essentielles d'un texte source. Il respecte trois principes fondamentaux :
\begin{enumerate}
\item \textbf{Fidélité} : aucune idée ajoutée, aucune opinion personnelle, aucun exemple détaillé conservé.
\item \textbf{Concision} : réduction au quart du nombre de mots du texte original (marge tolérée $\pm 10\%$).
\item \textbf{Cohérence} : reconstruction logique du texte (thèse, arguments, conclusion) grâce aux connecteurs.
\end{enumerate}
\end{aretenir}

\begin{methode}[Les quatre techniques de contraction]
\begin{enumerate}
\item \textbf{La suppression} : éliminer les exemples, détails, répétitions, adjectifs accessoires, incises.
\item \textbf{La substitution} : remplacer un groupe de mots par un terme plus général (hyperonyme) ou un pronom.
\item \textbf{La fusion} : regrouper plusieurs phrases en une seule en utilisant la subordination ou la coordination.
\item \textbf{La généralisation} : remplacer une énumération par un terme collectif ou une catégorie.
\end{enumerate}
\end{methode}

\begin{exercice}[QCM : le résumé et ses règles]
Pour chaque question, choisis la seule bonne réponse.
\begin{enumerate}
\item Un résumé de texte est :
\begin{enumerate}
\item une paraphrase du texte de départ ;
\item une reformulation brève, fidèle et personnelle des idées essentielles ;
\item une copie des phrases les plus importantes du texte.
\end{enumerate}
\item Si le texte de départ compte 480 mots et que la consigne impose le quart, le résumé doit compter :
\begin{enumerate}
\item 240 mots ; \item 100 mots ; \item 120 mots.
\end{enumerate}
\item La technique qui consiste à remplacer un groupe de mots par un seul mot (ex. « le père, la mère et les enfants » $\rightarrow$ « la famille ») s'appelle :
\begin{enumerate}
\item la suppression ; \item la substitution ; \item la fusion.
\end{enumerate}
\item Dans un résumé, l'opinion personnelle du résumant :
\begin{enumerate}
\item doit apparaître en conclusion ; \item est totalement interdite ; \item peut remplacer la thèse de l'auteur.
\end{enumerate}
\end{enumerate}
\end{exercice}

\begin{exercice}[Vrai ou faux ? Justifie]
Réponds par vrai ou faux, puis justifie en une phrase.
\begin{enumerate}
\item La cohésion d'un texte repose sur les procédés grammaticaux et lexicaux (reprises, connecteurs, pronoms) qui relient les phrases entre elles.
\item La cohérence d'un texte concerne uniquement l'orthographe et la ponctuation.
\item Un texte explicatif a pour fonction principale de faire comprendre un phénomène en donnant des informations, des causes et des conséquences.
\item La tonalité polémique vise à amuser le lecteur sans attaquer personne.
\end{enumerate}
\end{exercice}

\begin{exercice}[Définitions]
Définis en une ou deux phrases chacun des termes suivants, puis illustre chacun par un exemple de ton choix :
\begin{enumerate}
\item la thèse d'un texte argumentatif ;
\item l'argument d'autorité ;
\item le sigle (et la différence entre sigle épelé et acronyme) ;
\item la généralisation (technique de contraction).
\end{enumerate}
\end{exercice}

\begin{exercice}[Comptage et proportion]
Un candidat au Baccalauréat technique doit résumer un éditorial de 600 mots.
\begin{enumerate}
\item La consigne impose le quart du texte : combien de mots le résumé doit-il contenir ?
\item Le règlement tolère une marge de 10 \% en plus ou en moins. Calcule le nombre minimal et le nombre maximal de mots acceptés.
\item Le candidat a rédigé 185 mots. Son résumé est-il conforme ? Justifie par un calcul.
\end{enumerate}
\end{exercice}

\courssub{Je m'entraîne}

\begin{aretenir}[Cohésion vs Cohérence]
\begin{itemize}
\item \textbf{Cohésion} : ensemble des marques linguistiques (connecteurs, reprises pronominales/lexicales, énoncés déictiques, temps verbaux) qui assurent la liaison \emph{formelle} entre les phrases.
\item \textbf{Cohérence} : organisation \emph{sémantique} et logique du contenu (progression thématique, enchaînement des idées, non-contradiction) qui donne du sens au texte.
\end{itemize}
Un texte peut être cohésif sans être cohérent (enchainement grammatical correct mais idées absurdes), et vice-versa.
\end{aretenir}

\begin{savaistu}[Les sigles camerounais : épelés ou acronymes ?]
Au Cameroun, l'usage administratif distingue :
\begin{itemize}
\item \textbf{Sigles épelés} (lettres prononcées une à une) : \cle{MINESEC} [mi.ne.sɛk], \cle{MINESUP}, \cle{SNI}, \cle{CRTV}.
\item \textbf{Acronymes} (prononcés comme un mot unique) : \cle{SONARA} [so.na.ʁa], \cle{CAMRAIL} [kam.ʁaj], \cle{FEICOM} [fei.kɔm], \cle{ONEF} [ɔ.nɛf], \cle{CAMTEL} [kam.tɛl], \cle{ENEO} [e.ne.o].
\end{itemize}
Règle d'accord : l'article s'accorde avec le noyau du sigle (ex. \emph{la} SONARA car société, \emph{le} CAMRAIL car chemin de fer).
\end{savaistu}

\begin{exercice}[Repérage de la structure argumentative]
Lis l'éditorial suivant, puis réponds aux questions.

\emph{« Le téléphone portable est devenu l'outil indispensable de la jeunesse camerounaise. Selon l'Agence de régulation des télécommunications, plus de 80 \% des jeunes de 15 à 25 ans possèdent un smartphone. Cet outil facilite l'accès à l'information : un élève de Ngaoundéré peut suivre un cours en ligne dispensé depuis Yaoundé. J'ai moi-même constaté, dans mon lycée, que les élèves qui utilisent leur téléphone pour réviser progressent plus vite. Mais attention : si le portable est mal utilisé, il devient un instrument de distraction, car celui qui passe ses nuits sur les réseaux sociaux ne peut pas être attentif en classe. Il faut donc éduquer les jeunes à un usage raisonnable de cet outil. »}

\begin{enumerate}
\item Releve la phrase qui exprime la thèse de l'auteur.
\item Identifie trois arguments et classe-les selon leur type : argument d'autorité, d'expérience vécue ou logique.
\item Releve deux exemples et explique pourquoi ils devront être supprimés dans un résumé.
\item Releve deux connecteurs logiques et précise la relation qu'ils expriment.
\end{enumerate}
\end{exercice}

\begin{exercice}[Contraction ciblée : les quatre techniques]
Réduis chacune des phrases suivantes en appliquant la technique indiquée, puis justifie ton choix en une phrase.
\begin{enumerate}
\item (Suppression) « Le jeune Meka, âgé, fatigué, courbé par les années, reçut enfin sa médaille. »
\item (Substitution) « Les cultivateurs, les commerçants, les enseignants et les fonctionnaires de la région ont protesté. »
\item (Fusion) « Les pluies ont été abondantes. Les récoltes ont été excellentes. »
\item (Généralisation) « Les manguiers, les goyaviers et les papayers bordaient la cour de la case. »
\item (Technique libre, à identifier) « Il est évident, il est clair, il est manifeste que la jeunesse doit se former. »
\end{enumerate}
\end{exercice}

\begin{exercice}[Cohésion et cohérence]
Voici quatre phrases extraites d'un texte, présentées dans le désordre :
\begin{enumerate}
\item[a.] C'est pourquoi l'État camerounais encourage la formation professionnelle des jeunes.
\item[b.] Le chômage des jeunes constitue aujourd'hui un problème majeur au Cameroun.
\item[c.] Ces formations permettent en effet d'acquérir un métier et de créer sa propre entreprise.
\item[d.] En effet, de nombreux diplômés ne trouvent pas d'emploi correspondant à leurs études.
\end{enumerate}
\begin{enumerate}
\item Remets ces phrases dans l'ordre logique et justifie ton classement.
\item Releve dans les phrases les procédés de cohésion (connecteurs, reprises pronominales, reprises lexicales).
\item Remplace dans la phrase c le mot « formations » par un pronom sans nuire à la clarté du texte.
\end{enumerate}
\end{exercice}

\begin{exercice}[Les sigles camerounais]
Voici une liste de sigles : MINESEC, ENEO, CRTV, SNI, CAMTEL, MINESUP, FEICOM, SONARA, CAMRAIL, ONEF.
\begin{enumerate}
\item Classe ces sigles en deux colonnes : sigles épelés et acronymes (sigles prononcés comme un mot). Justifie ton classement pour deux d'entre eux.
\item Donne la signification complète de trois sigles de la liste.
\item Emploie correctement quatre sigles dans des phrases en respectant les règles d'accord et d'usage (article, majuscules).
\end{enumerate}
\end{exercice}

\courssub{Je me prépare au Bac}

\begin{methode}[Étapes de la rédaction d'un résumé au Bac]
\begin{enumerate}
\item \textbf{Lecture active} : identifier le genre, la thèse, le plan, les arguments, les connecteurs, le ton.
\item \textbf{Analyse} : bâtir le plan détaillé (idées principales + arguments essentiels), éliminer exemples/détails.
\item \textbf{Contraction} : appliquer les 4 techniques (suppression, substitution, fusion, généralisation) pour réduire au quart.
\item \textbf{Rédaction} : réécrire avec ses mots, en respectant l'ordre des idées, les connecteurs logiques, la troisième personne.
\item \textbf{Contrôle} : compter les mots (marge $\pm 10\%$), vérifier fidélité, cohérence, correction linguistique.
\end{enumerate}
\end{methode}

\begin{exercice}[Sujet type Baccalauréat : résumé d'un texte argumentatif]
Lis le texte suivant (environ 480 mots), puis traite les questions.

\begin{itshape}«es réseaux sociaux occupent désormais une place centrale dans la vie de la jeunesse camerounaise. Facebook, WhatsApp, TikTok : ces plateformes rythment les journées de millions d'adolescents et de jeunes adultes, de Douala à Maroua. Faut-il s'en réjouir ou s'en inquiéter ?

D'abord, il faut reconnaître que les réseaux sociaux offrent des avantages réels. Ils permettent une communication rapide et peu coûteuse : un étudiant de Bafoussam peut échanger gratuitement avec son frère installé à l'étranger. Ils constituent aussi un espace d'information et d'apprentissage : de nombreux jeunes y suivent des tutoriels, des cours de langues ou des formations en entrepreneuriat. Enfin, ils offrent des opportunités économiques : certains jeunes commerçants vendent leurs produits grâce à ces plateformes et gagnent ainsi leur vie.

Cependant, ces mêmes réseaux présentent des dangers qu'on ne saurait ignorer. La perte de temps est le premier d'entre eux : combien d'élèves sacrifient leurs heures de révision à des vidéos sans intérêt ? Les résultats scolaires s'en ressentent, comme l'ont montré plusieurs chefs d'établissement lors de la dernière rentrée. Ensuite, la désinformation circule massivement : de fausses nouvelles se répandent plus vite que la vérité et trompent des jeunes peu méfiants. Enfin, certains contenus exposent les adolescents à la violence, à la fraude ou aux mauvaises rencontres.

Face à ce constat, que faire ? Interdire les réseaux sociaux serait à la fois impossible et contre-productif. La solution passe par l'éducation : les familles et l'école doivent apprendre aux jeunes à utiliser ces outils avec discernement, à vérifier les informations et à limiter leur temps de connexion. Car un outil n'est ni bon ni mauvais en soi : tout dépend de l'usage qu'on en fait.

En définitive, les réseaux sociaux sont une chance pour la jeunesse camerounaise, à condition qu'elle apprenne à les dominer au lieu de se laisser dominer par eux. »\end{itshape}

\begin{enumerate}
\item Dégage la thèse de l'auteur et présente le plan du texte en trois ou quatre grandes étapes.
\item Releve quatre arguments et précise pour chacun son type (autorité, expérience, logique, fait).
\item Rédige un résumé du texte en 120 mots (marge tolérée : plus ou moins 10 \%), en appliquant les techniques de contraction et sans aucun commentaire personnel. Indique le nombre de mots à la fin.
\end{enumerate}
\end{exercice}

\begin{exercice}[Correction d'un résumé fautif]
Voici le résumé rédigé par un candidat à partir du texte de l'exercice précédent (consigne : 120 mots, fidélité totale à l'auteur) :

\emph{« Les réseaux sociaux, comme Facebook que j'utilise tous les jours, occupent une place centrale dans la vie des jeunes. Ces plateformes offrent des avantages réels : communication rapide, information, apprentissage et opportunités économiques. Un étudiant de Bafoussam peut échanger gratuitement avec son frère installé à l'étranger. Cependant, ils présentent des dangers : perte de temps, désinformation et mauvaises rencontres. Personnellement, je pense que les réseaux sociaux devraient être interdits aux moins de 18 ans. La solution passe par l'éducation des jeunes à un usage raisonnable. Les réseaux sociaux sont donc une chance si les jeunes apprennent à les dominer. »} (138 mots)

\begin{enumerate}
\item Repère dans ce résumé les cinq erreurs commises par le candidat (paraphrase, exemple conservé, opinion personnelle contraire à la thèse de l'auteur, rupture de cohésion, mauvais comptage) et cite le passage fautif à chaque fois.
\item Explique pour chaque erreur la règle du résumé qui a été violée.
\item Propose une version améliorée du résumé, conforme à la consigne (120 mots, plus ou moins 10 \%).
\end{enumerate}
\end{exercice}

\begin{savaistu}[Ferdinand Oyono et \emph{Le Vieux nègre et la médaille}]
\cle{Ferdinand Oyono} (1929--2010), haut fonctionnaire et diplomate camerounais, publie ce roman en 1956. L'œuvre s'inscrit dans la littérature anticoloniale francophone (comme \emph{Le Pauvre Christ de Bomba} de Mongo Beti). Le Cameroun est alors sous tutelle française. Le roman dénonce l'assimilation coloniale et la collaboration de certaines élites traditionnelles. La médaille symbolise la duperie coloniale : on récompense l'aliénation (don de terres, enrôlement des fils) par une distinction honorifique vide.
\end{savaistu}

\begin{exercice}[Synthèse sur l'œuvre : satire et polémique dans \emph{Le Vieux nègre et la médaille}]
Dans le roman de Ferdinand Oyono, Meka, vieux paysan camerounais, reçoit une médaille de l'administration coloniale pour avoir « donné » ses terres à la mission et ses fils à la guerre. La cérémonie, l'attitude des autorités coloniales et celle du commandant de cercle sont décrites avec ironie.
\begin{enumerate}
\item Rappelle le contexte historique et littéraire dans lequel s'inscrit l'œuvre (époque, situation du Cameroun, position de l'auteur).
\item Releve dans le roman trois passages à tonalité satirique ou polémique que tu as étudiés. Pour chacun, précise : la cible moquée ou attaquée, le procédé utilisé (ironie, caricature, antiphrase, exagération) et l'effet produit sur le lecteur.
\item Rédige une réponse construite (introduction, développement en deux ou trois paragraphes, conclusion) à la question de synthèse suivante : « Montrez comment la satire et la polémique servent la dénonciation de la colonisation dans \emph{Le Vieux nègre et la médaille}. »
\end{enumerate}
\end{exercice}

\newpage

\coursec{Corrigés des exercices}

\coursec{Leçon 10 : Contraction de texte — Rédiger un résumé de texte}
\courssub{Exercices d'application et sujets d'examen corrigés}

\begin{corrige}[Exercice 1 — QCM : le résumé et ses règles]
\begin{enumerate}
\item \textbf{b.} Un résumé est une reformulation brève, fidèle et personnelle des idées essentielles. La paraphrase (a) conserve la structure et la longueur ; la copie (c) n'est pas une reformulation.
\item \textbf{c.} $480 \div 4 = 120$ mots. La consigne « le quart » impose une division par 4.
\item \textbf{b.} Remplacer « le père, la mère et les enfants » par « la famille » est une \cle{substitution} par hyperonyme (terme plus général). La suppression (a) efface ; la fusion (c) regroupe des phrases.
\item \textbf{b.} L'opinion personnelle est \emph{totalement interdite} dans un résumé : seule la pensée de l'auteur doit être restituée.
\end{enumerate}
\end{corrige}

\begin{corrige}[Exercice 2 — Vrai ou faux ? Justifie]
\begin{enumerate}
\item \textbf{Vrai.} La cohésion repose sur les marques linguistiques formelles (connecteurs, reprises pronominales/lexicales, déictiques, temps) qui chaînent les énoncés.
\item \textbf{Faux.} La cohérence concerne l'organisation \emph{sémantique et logique} (progression thématique, non-contradiction), non l'orthographe ni la ponctuation (qui relèvent de la correction formelle).
\item \textbf{Vrai.} Le texte explicatif a pour fonction dominante d'informer, de faire comprendre un phénomène en en donnant les causes, les mécanismes et les conséquences.
\item \textbf{Faux.} La tonalité \cle{polémique} vise à \emph{attaquer}, contester ou réfuter une thèse adverse, souvent avec virulence ; elle ne vise pas à amuser (c'est la tonalité \cle{comique} ou \cle{satirique}).
\end{enumerate}
\end{corrige}

\begin{corrige}[Exercice 3 — Définitions]
\begin{enumerate}
\item \textbf{Thèse :} Proposition principale que l'auteur défend et cherche à faire admettre dans un texte argumentatif. \textit{Exemple :} « Les réseaux sociaux sont une chance pour la jeunesse camerounaise. »
\item \textbf{Argument d'autorité :} Argument qui s'appuie sur le témoignage ou l'avis d'une personne/institution reconnue compétente dans le domaine. \textit{Exemple :} « Selon l'OMS, le tabac tue plus de 8 millions de personnes par an. »
\item \textbf{Sigle :} Abréviation formée des initiales d'une expression. \textbf{Sigle épelé :} lettres prononcées séparément (ex. \cle{MINESEC} [mi.ne.sɛk]). \textbf{Acronyme :} sigle prononcé comme un mot unique (ex. \cle{SONARA} [so.na.ʁa]).
\item \textbf{Généralisation (contraction) :} Remplacement d'une énumération ou d'une série de termes par un mot plus général (catégorie, collectif). \textit{Exemple :} « Les manguiers, les goyaviers et les papayers » $\rightarrow$ « Les arbres fruitiers ».
\end{enumerate}
\end{corrige}

\begin{corrige}[Exercice 4 — Comptage et proportion]
\begin{enumerate}
\item $600 \div 4 = \textbf{150 mots}$.
\item Marge $\pm 10\%$ de 150 = 15 mots. \textbf{Minimum :} $150 - 15 = 135$ mots. \textbf{Maximum :} $150 + 15 = 165$ mots.
\item \textbf{Non conforme.} Le résumé compte 185 mots, ce qui dépasse le maximum autorisé (165 mots) de 20 mots. Le candidat n'a pas respecté la contrainte de concision.
\end{enumerate}
\end{corrige}

\begin{corrige}[Exercice 5 — Repérage de la structure argumentative]
\begin{enumerate}
\item \textbf{Thèse :} « Il faut donc éduquer les jeunes à un usage raisonnable de cet outil. » (phrase finale, recommandation). La première phrase (« Le téléphone portable est devenu l'outil indispensable… ») pose le constat de départ.
\item \textbf{Trois arguments :}
\begin{itemize}
\item « Selon l'Agence de régulation…, plus de 80\%… » $\rightarrow$ \textbf{Argument d'autorité} (source officielle, statistique).
\item « J'ai moi-même constaté, dans mon lycée, que les élèves qui utilisent leur téléphone pour réviser progressent plus vite. » $\rightarrow$ \textbf{Argument d'expérience vécue} (témoignage personnel).
\item « Celui qui passe ses nuits sur les réseaux sociaux ne peut pas être attentif en classe. » $\rightarrow$ \textbf{Argument logique} (relation cause-effet, raisonnement déductif).
\end{itemize}
\item \textbf{Deux exemples à supprimer :}
\begin{itemize}
\item « Un élève de Ngaoundéré peut suivre un cours en ligne dispensé depuis Yaoundé. » (illustration concrète de l'accès à l'information).
\item « Dans mon lycée, les élèves qui utilisent leur téléphone pour réviser progressent plus vite. » (exemple personnel appuyé).
\end{itemize}
\textit{Justification :} Un résumé ne conserve que les idées essentielles (thèse, arguments) ; les exemples, illustrations et détails sont des \emph{détails} à éliminer (technique de la \cle{suppression}).
\item \textbf{Connecteurs logiques :}
\begin{itemize}
\item « \textbf{Mais} » (début du 3\textsuperscript{e} paragraphe) : exprime l'\textbf{opposition} / la concession (transition avantages $\rightarrow$ dangers).
\item « \textbf{Car} » (dans « car celui qui passe ses nuits… ») : exprime la \textbf{cause} / l'explication.
\item « \textbf{Donc} » (dans « Il faut donc éduquer… ») : exprime la \textbf{conséquence} / la conclusion.
\end{itemize}
\end{enumerate}
\end{corrige}

\begin{corrige}[Exercice 6 — Contraction ciblée : les quatre techniques]
\begin{enumerate}
\item \textbf{Suppression :} « Le vieux Meka reçut enfin sa médaille. » \textit{Justification :} Suppression des adjectifs accessoires (« âgé, fatigué, courbé par les années ») et de l'apposition « jeune » (ironique) qui sont des détails descriptifs.
\item \textbf{Substitution :} « Les actifs de la région ont protesté. » ou « La population active a protesté. » \textit{Justification :} Remplacement de l'énumération « cultivateurs, commerçants, enseignants, fonctionnaires » par l'hyperonyme « actifs » ou « population active ».
\item \textbf{Fusion :} « Les pluies abondantes ont entraîné d'excellentes récoltes. » \textit{Justification :} Regroupement des deux phrases par subordination (participial « abondantes » + verbe « ont entraîné ») marquant la relation cause-effet.
\item \textbf{Généralisation :} « Des arbres fruitiers bordaient la cour de la case. » \textit{Justification :} Remplacement de l'énumération spécifique « manguiers, goyaviers, papayers » par la catégorie collective « arbres fruitiers ».
\item \textbf{Suppression (redondance) :} « Il est évident que la jeunesse doit se former. » \textit{Justification :} Les trois propositions « Il est évident, il est clair, il est manifeste » sont synonymes (redondance) ; on n'en garde qu'une par \cle{suppression} des répétitions.
\end{enumerate}
\end{corrige}

\begin{corrige}[Exercice 7 — Cohésion et cohérence]
\begin{enumerate}
\item \textbf{Ordre logique :} \textbf{b $\rightarrow$ d $\rightarrow$ a $\rightarrow$ c.}
\textit{Justification :} 
\begin{itemize}
\item \textbf{b} pose le \textbf{thème/problème} (chômage des jeunes).
\item \textbf{d} apporte l'\textbf{explication} (« En effet ») : cause du problème (adéquation formation-emploi).
\item \textbf{a} tire la \textbf{conséquence} politique (« C'est pourquoi ») : action de l'État.
\item \textbf{c} justifie l'efficacité de cette action (« en effet », « permettent ») : finalité des formations.
\end{itemize}
\item \textbf{Procédés de cohésion repérés :}
\begin{itemize}
\item \textbf{Connecteurs :} « En effet » (d, c — explication), « C'est pourquoi » (a — conséquence).
\item \textbf{Reprise pronominale :} « Ces formations » (c) reprend « la formation professionnelle » (a).
\item \textbf{Reprise lexicale / champ sémantique :} « jeunes » (b, a), « emploi/entreprise » (d, c), « études/métier » (d, c).
\item \textbf{Chaînage thématique :} problème $\rightarrow$ cause $\rightarrow$ solution $\rightarrow$ validation.
\end{itemize}
\item \textbf{Remplacement :} « \textbf{Elles} permettent en effet d'acquérir un métier et de créer sa propre entreprise. » (Pronom « Elles » = antécédent « formations », féminin pluriel, sujet de la phrase).
\end{enumerate}
\end{corrige}

\begin{corrige}[Exercice 8 — Les sigles camerounais]
\begin{enumerate}
\item \textbf{Classement :}
\begin{center}
\begin{tabularx}{\linewidth}{|l|X|}\hline
\textbf{Sigles épelés (lettres)} & \textbf{Acronymes (prononcés comme un mot)} \\ \hline
MINESEC, CRTV, SNI, MINESUP & ENEO, FEICOM, SONARA, CAMRAIL, ONEF, CAMTEL \\ \hline
\end{tabularx}
\end{center}
\textit{Justifications :}
\begin{itemize}
\item \textbf{MINESEC} : se prononce [mi.ne.sɛk] (M-I-N-E-S-E-C), lettres distinctes $\rightarrow$ sigle épelé.
\item \textbf{SONARA} : se prononce [so.na.ʁa] comme un mot unique $\rightarrow$ acronyme.
\end{itemize}
\item \textbf{Significations (trois exemples) :}
\begin{itemize}
\item \cle{MINESEC} = Ministère des Enseignements Secondaires.
\item \cle{SONARA} = Société Nationale de Raffinage.
\item \cle{CAMRAIL} = Cameroon Railways (Chemin de fer du Cameroun).
\item \cle{ENEO} = Energy of Cameroon (électricité).
\item \cle{FEICOM} = Fonds Spécial d'Équipement et d'Intervention Communale.
\end{itemize}
\item \textbf{Phrases d'emploi correct (quatre exemples) :}
\begin{enumerate}
\item \textbf{Le} MINESEC a publié le calendrier des examens. (masculin : ministère)
\item \textbf{La} SONARA assure le raffinage du pétrole brut. (féminin : société)
\item \textbf{Le} CAMRAIL gère le réseau ferroviaire national. (masculin : chemin de fer)
\item \textbf{L'}\cle{ENEO} (ou \textbf{La} société ENEO) fournit l'énergie électrique. (masculin : energy / société)
\end{enumerate}
\end{enumerate}
\end{corrige}

\begin{corrige}[Exercice 9 — Sujet type Baccalauréat : résumé d'un texte argumentatif]
\textbf{Barème indicatif :} Thèse + plan (4 pts), Arguments typés (4 pts), Résumé : fidélité/concision/cohérence/correction (12 pts). Total 20 pts.

\begin{enumerate}
\item \textbf{Thèse :} Les réseaux sociaux constituent une chance pour la jeunesse camerounaise, \emph{à condition} qu'une éducation à l'usage raisonnable (discernement, vérification, limitation du temps) permette aux jeunes de les dominer au lieu de les subir.
\textbf{Plan en quatre étapes :}
\begin{enumerate}
\item \textbf{Avantages :} communication gratuite/rapide, espace d'information/apprentissage (tutoriels, langues, entrepreneuriat), opportunités économiques (commerce en ligne).
\item \textbf{Dangers :} perte de temps nuisant aux résultats scolaires, désinformation massive (fausses nouvelles), exposition à la violence/fraude/mauvaises rencontres.
\item \textbf{Solution :} Éducation (famille + école) à un usage discernant et limité ; interdiction impossible et contre-productive.
\item \textbf{Conclusion :} Les réseaux sont une chance \emph{si} la jeunesse apprend à les maîtriser.
\end{enumerate}
\item \textbf{Quatre arguments et types :}
\begin{enumerate}
\item « Un étudiant de Bafoussam peut échanger gratuitement avec son frère installé à l'étranger. » $\rightarrow$ \textbf{Argument de fait} (exemple concret, vérifiable).
\item « De nombreux jeunes y suivent des tutoriels, des cours de langues ou des formations en entrepreneuriat. » $\rightarrow$ \textbf{Argument de fait} (constat observable).
\item « Comme l'ont montré plusieurs chefs d'établissement lors de la dernière rentrée. » $\rightarrow$ \textbf{Argument d'autorité} (témoignage d'acteurs légitimes du système éducatif).
\item « De fausses nouvelles se répandent plus vite que la vérité et trompent des jeunes peu méfiants. » $\rightarrow$ \textbf{Argument logique} (mécanisme cause-effet : viralité $\rightarrow$ tromperie).
\end{enumerate}
\item \textbf{Résumé modèle (120 mots $\pm$ 10\% = 108 à 132 mots) :}
\begin{quote}
Les réseaux sociaux (Facebook, WhatsApp, TikTok) occupent une place centrale dans la vie des jeunes Camerounais. Ils offrent de réels avantages : communication gratuite et rapide, accès à l'information et à l'apprentissage (tutoriels, langues, entrepreneuriat), opportunités économiques pour de jeunes commerçants. Cependant, ils présentent des dangers majeurs : perte de temps au détriment des études, propagation massive de fausses nouvelles, exposition à la violence, à la fraude et aux mauvaises rencontres. Interdire ces plateformes serait impossible et contre-productif. La solution réside dans l'éducation : familles et école doivent apprendre aux jeunes à les utiliser avec discernement, à vérifier les informations et à limiter leur temps de connexion. En définitive, les réseaux sociaux sont une chance pour la jeunesse, à condition qu'elle apprenne à les dominer au lieu de s'en laisser dominer.
\end{quote}
\textbf{Nombre de mots : 122 mots} (conforme : $108 \le 122 \le 132$).
\textit{Points de vigilance :} Fidélité à la thèse (chance conditionnée), suppression des exemples détaillés (Bafoussam, chefs d'établissement), fusion des listes d'avantages/dangers, connecteurs logiques conservés (Cependant, En définitive), 3\textsuperscript{e} personne, pas d'opinion personnelle.
\end{enumerate}
\end{corrige}

\begin{corrige}[Exercice 10 — Correction d'un résumé fautif]
\begin{enumerate}
\item \textbf{Cinq erreurs repérées :}
\begin{enumerate}
\item \textbf{Paraphrase / intrusion personnelle :} « \emph{comme Facebook que j'utilise tous les jours} » (ajout d'expérience personnelle absente du texte source).
\item \textbf{Exemple conservé (détail) :} « \emph{Un étudiant de Bafoussam peut échanger gratuitement avec son frère installé à l'étranger.} » (illustration concrète à supprimer).
\item \textbf{Opinion personnelle contraire à la thèse :} « \emph{Personnellement, je pense que les réseaux sociaux devraient être interdits aux moins de 18 ans.} » (l'auteur écrit : « Interdire serait impossible et contre-productif »).
\item \textbf{Rupture de cohésion / accord :} « \emph{Ces plateformes offrent… Cependant, \textbf{ils} présentent…} » (accord incorrect : « plateformes » fém. pl. $\rightarrow$ « elles » ; ou reprise de « réseaux sociaux » masc. pl. mais changement de référent sans marqueur clair).
\item \textbf{Mauvais comptage :} 138 mots $>$ 132 mots (max autorisé pour 120 $\pm$ 10\%).
\end{enumerate}
\item \textbf{Règles violées :}
\begin{enumerate}
\item \cle{Fidélité} : aucune idée/opinion/exemple personnel ne doit être ajouté.
\item \cle{Concision} / \cle{Suppression} : les exemples illustrant un argument doivent être éliminés.
\item \cle{Fidélité} : le résumé doit restituer \emph{exactement} la thèse de l'auteur (éducation, pas interdiction).
\item \cle{Cohérence} / \cle{Correction} : respect de la cohésion grammaticale (accords, reprises) et logique.
\item \cle{Concision} : respect strict du nombre de mots imposé (marge $\pm 10\%$).
\end{enumerate}
\item \textbf{Version améliorée (120 mots $\pm$ 10\%) :}
\begin{quote}
Les réseaux sociaux (Facebook, WhatsApp, TikTok) sont centraux dans la vie des jeunes Camerounais. Ils apportent des avantages : communication gratuite, information, apprentissage, opportunités économiques. Mais ils posent des problèmes : perte de temps nuisant aux études, désinformation massive, exposition à la violence et à la fraude. L'interdiction est impossible et contre-productive. La solution passe par l'éducation des jeunes (famille, école) à un usage discernant, à la vérification des informations et à la limitation du temps de connexion. Ainsi, ces réseaux sont une chance si la jeunesse apprend à les maîtriser au lieu de les subir.
\end{quote}
\textbf{Nombre de mots : 118 mots} (conforme).
\end{enumerate}
\end{corrige}

\begin{corrige}[Exercice 11 — Synthèse sur l'œuvre : satire et polémique dans \emph{Le Vieux nègre et la médaille}]
\begin{enumerate}
\item \textbf{Contexte historique et littéraire :}
\begin{itemize}
\item \textbf{Époque :} 1956, Cameroun sous \emph{tutelle française} (ONU), période pré-indépendance (1960).
\item \textbf{Situation :} Administration coloniale, système d'indirect rule, collaboration de certaines chefferies traditionnelles, assimilation culturelle imposée.
\item \textbf{Position de l'auteur :} \cle{Ferdinand Oyono}, haut fonctionnaire et diplomate, écrit dans la veine \textbf{anticoloniale} (proche de \emph{Le Pauvre Christ de Bomba} de Mongo Beti). Il dénonce l'aliénation, la duperie de la « mission civilisatrice » et la complicité des élites traditionnelles.
\end{itemize}
\item \textbf{Trois passages à tonalité satirique/polémique (exemples types étudiés) :}
\begin{enumerate}
\item \textbf{La cérémonie de remise de la médaille (Ch. 4-5).}
\textit{Cible :} L'administration coloniale (commandant, prêtre, chef) et le rituel humiliant.
\textit{Procédés :} \textbf{Ironie} (description solennelle d'une mascarade), \textbf{Caricature} (gestes ridicules du commandant, discours creux), \textbf{Antiphrase} (« distinction honorifique » pour une duperie).
\textit{Effet :} Rire jaune, prise de conscience du ridicule et de la violence symbolique.
\item \textbf{Le discours du commandant de cercle (Ch. 5).}
\textit{Cible :} La rhétorique coloniale paternaliste (« nos frères noirs », « bienfaits de la France »).
\textit{Procédés :} \textbf{Antiphrase} systématique, \textbf{Euphémisme} (euphémisation de la spoliation : « donation » des terres), \textbf{Parodie} du langage officiel.
\textit{Effet :} Dénonciation du mensonge d'État, colère du lecteur face à l'hypocrisie.
\item \textbf{La scène du « don » des terres à la mission / départ des fils à la guerre (rappel narratif).}
\textit{Cible :} La collaboration du chef Meka (aliénation volontaire) et l'injustice coloniale.
\textit{Procédés :} \textbf{Exagération} (générosité démesurée de Meka), \textbf{Contraste} (richesse donnée / médaille de toc reçue), \textbf{Satire} du naïf qui croit à la reconnaissance.
\textit{Effet :} Pitié critique pour Meka, condamnation du système qui exploite la loyauté.
\end{enumerate}
\item \textbf{Réponse construite : « Montrez comment la satire et la polémique servent la dénonciation de la colonisation dans \emph{Le Vieux nègre et la médaille}. »}
\textbf{Introduction :} 
\emph{Le Vieux nègre et la médaille} (1956) de Ferdinand Oyono est un roman emblématique de la littérature anticoloniale camerounaise. Sous couvert d'une intrigue simple — le vieux chef Meka reçoit une médaille pour sa « loyauté » —, l'œuvre déploie une charge féroce contre le système colonial. La \cle{satire} (rire critique) et la \cle{polémique} (attaque frontale) sont les deux armes rhétoriques majeures qui permettent à Oyono de dévoiler la duperie de l'assimilation et la tragédie des collaborateurs.

\textbf{Développement 1 : La satire, dévoilement du ridicule et de la duperie coloniale.}
La satire opère par \textbf{ironie}, \textbf{caricature} et \textbf{antiphrase} pour transformer la cérémonie officielle en une farce grotesque. Le commandant de cercle, figure d'autorité, est réduit à un pantin récitant des platitudes ; le prêtre bénit une médaille qui ne vaut pas le métal dont elle est faite ; Meka, vêtu de son habit neuf, incarne le « bon nègre » naïf. La description minutieuse des gestes (remise de la médaille, embrassades, discours) crée un \emph{décalage} permanent entre la solennité affichée et la vacuité réelle. Ce rire amer — \emph{le rire de l'intelligence} — force le lecteur à voir l'absurdité du rituel colonial : on récompense l'aliénation (don des terres, sacrifice des fils) par un symbole vide. La satire atteint son paroxysme quand Meka, ivre de fierté, ne comprend pas qu'il est le dindon de la farce, alors que le lecteur, lui, perçoit la \textbf{duperie} totale.

\textbf{Développement 2 : La polémique, accusation directe de l'injustice et de la collaboration.}
Au-delà du rire, le texte bascule dans la \textbf{polémique} : attaque frontale, sans détours, contre les piliers du système. Le discours du commandant est une \emph{profession de foi coloniale} que le récit démontre par l'exemple de Meka : les « bienfaits » sont spoliation, la « fraternité » est domination. Oyono pointe aussi la responsabilité des élites traditionnelles (Meka, le chef Kelara) qui, par ignorance ou intérêt, servent l'oppresseur. La polémique s'incarne dans le destin des fils de Meka (morts à la guerre pour la France) et dans le sort final du vieux homme (dépouillé, humilié, mourant). L'indignation de l'auteur-narrateur transparaît dans le refus de toute complaisance : la colonisation n'est pas une « rencontre », c'est un système de \textbf{violence structurelle} que la médaille symbolise parfaitement.

\textbf{Conclusion :}
Satire et polémique sont indissociables chez Oyono : la première \emph{démasque} le mécanisme (le ridicule du pouvoir, la naïveté du collaborateur), la seconde \emph{condamne} le crime (l'exploitation, le mensonge, la mort). Ensemble, elles font du roman non seulement une chronique de la fin d'un monde, mais un \emph{acte d'accusation} lucide et définitif contre l'entreprise coloniale. Le lecteur referme le livre avec la certitude que la médaille, loin d'honorer Meka, marque au fer rouge la honte d'un système.
\end{enumerate}
\end{corrige}

\newpage

\coursec{Fiche bilan}

\begin{popfigure}
\begin{tikzpicture}[mindmap, grow cyclic, every node/.style={concept}, root concept/.append style={concept color=popblue, font=\bfseries\small}, level 1/.append style={level distance=3.4cm, sibling angle=51.5}, text=white]
\node[root concept]{Contraction de texte : le résumé}
  child[concept color=poporange]{ node[concept]{Texte explicatif} 
    child[concept color=poporange!70]{ node[concept]{Informer, expliquer} }
    child[concept color=poporange!70]{ node[concept]{Progression : thème, chaîne, dérivation} }
  }
  child[concept color=popgreen]{ node[concept]{Cohésion \& cohérence}
    child[concept color=popgreen!70]{ node[concept]{Connecteurs, reprises} }
    child[concept color=popgreen!70]{ node[concept]{Sigles} }
  }
  child[concept color=poppurple]{ node[concept]{Le Vieux nègre et la médaille}
    child[concept color=poppurple!70]{ node[concept]{Contexte colonial} }
    child[concept color=poppurple!70]{ node[concept]{Tonalités satirique et polémique} }
  }
  child[concept color=poppink]{ node[concept]{Structure argumentative}
    child[concept color=poppink!70]{ node[concept]{Thèse, arguments, exemples} }
  }
  child[concept color=popgold]{ node[concept]{Techniques du résumé}
    child[concept color=popgold!70]{ node[concept]{Lire, repérer, reformuler} }
    child[concept color=popgold!70]{ node[concept]{Supprimer, généraliser} }
  }
  child[concept color=popmuted]{ node[concept]{Plan et rédaction}
    child[concept color=popmuted!70]{ node[concept]{Plan du texte} }
    child[concept color=popmuted!70]{ node[concept]{Rédiger, confronter, améliorer} }
  };
\end{tikzpicture}
\legende{Carte mentale de la leçon : les six compétences clés pour réussir le résumé de texte au Baccalauréat technique.}
\end{popfigure}

\begin{bilan}
\textbf{Définitions clés}
\begin{itemize}
  \item \cle{Résumé} : texte bref qui restitue fidèlement les idées essentielles d'un texte, sans commentaire personnel ni citation longue.
  \item \cle{Texte explicatif} : texte qui informe et explique ; il est clair, objectif, ordonné, avec définitions, exemples et connecteurs logiques.
  \item \cle{Progression thématique} : reprise du même thème ; \cle{progression en chaîne} : la règle d'une phrase devient le thème de la suivante ; \cle{progression dérivée} : les thèmes dérivent d'un hyperthème.
  \item \cle{Cohésion} : liens formels du texte (connecteurs, reprises anaphoriques, substitutions) ; \cle{cohérence} : unité de sens et logique d'ensemble.
  \item \cle{Sigle} : mot formé des initiales d'un groupe de mots, lu lettre par lettre (ONU, MINESEC).
  \item \cle{Tonalité satirique} : rire critique qui dénonce ; \cle{tonalité polémique} : attaque vive et argumentée.
\end{itemize}

\textbf{Repères à connaître par cœur}
\begin{center}
\begin{tabularx}{\linewidth}{|l|X|}
\hline
\rowcolor{popblueL} \textbf{Notion} & \textbf{À retenir} \\
\hline
Structure argumentative & Thèse (idée défendue) $\rightarrow$ arguments $\rightarrow$ exemples $\rightarrow$ conclusion \\
\hline
Étapes du résumé & Lire (2 fois) $\rightarrow$ souligner les idées clés $\rightarrow$ établir le plan $\rightarrow$ reformuler $\rightarrow$ rédiger $\rightarrow$ vérifier \\
\hline
Techniques de contraction & Supprimer (exemples, répétitions), généraliser (liste $\rightarrow$ terme générique), reformuler (synonymes, nominalisation) \\
\hline
Règles du résumé & Fidélité, brièveté, clarté ; présent de vérité ; pas de « je » ; pas de jugement personnel \\
\hline
\emph{Le Vieux nègre et la médaille} & Ferdinand Oyono (1956) : dénonciation satirique de l'aliénation coloniale ; bilan : ironie, naïveté de Meka, critique de l'administration et des missions \\
\hline
\end{tabularx}
\end{center}

\textbf{Méthode express}
\begin{enumerate}
  \item Identifier le thème et la thèse de l'auteur.
  \item Découper le texte en grandes parties (plan).
  \item Dégager une idée essentielle par partie.
  \item Rédiger avec ses propres mots, en respectant le nombre de mots imposé.
  \item Relire : cohésion (connecteurs), cohérence, orthographe.
\end{enumerate}
\end{bilan}

\begin{aretenir}[Checklist avant le Bac]
\begin{itemize}
  \item $\square$ Je sais définir le résumé et citer ses trois qualités (fidélité, brièveté, clarté).
  \item $\square$ Je reconnais un texte explicatif et ses caractéristiques.
  \item $\square$ Je distingue les trois types de progression thématique.
  \item $\square$ Je repère les connecteurs qui assurent la cohésion du texte.
  \item $\square$ Je sais ce qu'est un sigle et je cite des exemples camerounais (MINESEC, GCE).
  \item $\square$ Je repère la thèse, les arguments et les exemples d'un texte argumentatif.
  \item $\square$ Je maîtrise les techniques de contraction : supprimer, généraliser, reformuler.
  \item $\square$ Je rédige au présent, sans « je », sans copier de phrases du texte.
  \item $\square$ Je respecte le nombre de mots demandé et je l'indique à la fin.
  \item $\square$ Je situe \emph{Le Vieux nègre et la médaille} dans le contexte colonial.
  \item $\square$ Je distingue les tonalités satirique et polémique dans l'œuvre d'Oyono.
  \item $\square$ Je relis ma copie pour corriger orthographe et accords avant de rendre.
\end{itemize}
\end{aretenir}

\findecours

\end{document}
